% Options for packages loaded elsewhere
\PassOptionsToPackage{unicode}{hyperref}
\PassOptionsToPackage{hyphens}{url}
\PassOptionsToPackage{dvipsnames,svgnames,x11names}{xcolor}
\documentclass[
  american,
  10pt,
]{report}
\usepackage{xcolor}
\usepackage[a4paper,margin=25mm]{geometry}
\usepackage{amsmath,amssymb}
\setcounter{secnumdepth}{5}
\usepackage{iftex}
\ifPDFTeX
  \usepackage[T1]{fontenc}
  \usepackage[utf8]{inputenc}
  \usepackage{textcomp} % provide euro and other symbols
\else % if luatex or xetex
  \usepackage{unicode-math} % this also loads fontspec
  \defaultfontfeatures{Scale=MatchLowercase}
  \defaultfontfeatures[\rmfamily]{Ligatures=TeX,Scale=1}
\fi
\usepackage{lmodern}
\ifPDFTeX\else
  % xetex/luatex font selection
  \setmainfont[]{TeX Gyre Pagella}
  \setsansfont[]{TeX Gyre Heros}
  \setmonofont[]{DejaVu Sans Mono}
\fi
% Use upquote if available, for straight quotes in verbatim environments
\IfFileExists{upquote.sty}{\usepackage{upquote}}{}
\IfFileExists{microtype.sty}{% use microtype if available
  \usepackage[]{microtype}
  \UseMicrotypeSet[protrusion]{basicmath} % disable protrusion for tt fonts
}{}
\makeatletter
\@ifundefined{KOMAClassName}{% if non-KOMA class
  \IfFileExists{parskip.sty}{%
    \usepackage{parskip}
  }{% else
    \setlength{\parindent}{0pt}
    \setlength{\parskip}{6pt plus 2pt minus 1pt}}
}{% if KOMA class
  \KOMAoptions{parskip=half}}
\makeatother
\usepackage{color}
\usepackage{fancyvrb}
\newcommand{\VerbBar}{|}
\newcommand{\VERB}{\Verb[commandchars=\\\{\}]}
\DefineVerbatimEnvironment{Highlighting}{Verbatim}{commandchars=\\\{\}}
% Add ',fontsize=\small' for more characters per line
\usepackage{framed}
\definecolor{shadecolor}{RGB}{248,248,248}
\newenvironment{Shaded}{\begin{snugshade}}{\end{snugshade}}
\newcommand{\AlertTok}[1]{\textcolor[rgb]{0.94,0.16,0.16}{#1}}
\newcommand{\AnnotationTok}[1]{\textcolor[rgb]{0.56,0.35,0.01}{\textbf{\textit{#1}}}}
\newcommand{\AttributeTok}[1]{\textcolor[rgb]{0.13,0.29,0.53}{#1}}
\newcommand{\BaseNTok}[1]{\textcolor[rgb]{0.00,0.00,0.81}{#1}}
\newcommand{\BuiltInTok}[1]{#1}
\newcommand{\CharTok}[1]{\textcolor[rgb]{0.31,0.60,0.02}{#1}}
\newcommand{\CommentTok}[1]{\textcolor[rgb]{0.56,0.35,0.01}{\textit{#1}}}
\newcommand{\CommentVarTok}[1]{\textcolor[rgb]{0.56,0.35,0.01}{\textbf{\textit{#1}}}}
\newcommand{\ConstantTok}[1]{\textcolor[rgb]{0.56,0.35,0.01}{#1}}
\newcommand{\ControlFlowTok}[1]{\textcolor[rgb]{0.13,0.29,0.53}{\textbf{#1}}}
\newcommand{\DataTypeTok}[1]{\textcolor[rgb]{0.13,0.29,0.53}{#1}}
\newcommand{\DecValTok}[1]{\textcolor[rgb]{0.00,0.00,0.81}{#1}}
\newcommand{\DocumentationTok}[1]{\textcolor[rgb]{0.56,0.35,0.01}{\textbf{\textit{#1}}}}
\newcommand{\ErrorTok}[1]{\textcolor[rgb]{0.64,0.00,0.00}{\textbf{#1}}}
\newcommand{\ExtensionTok}[1]{#1}
\newcommand{\FloatTok}[1]{\textcolor[rgb]{0.00,0.00,0.81}{#1}}
\newcommand{\FunctionTok}[1]{\textcolor[rgb]{0.13,0.29,0.53}{\textbf{#1}}}
\newcommand{\ImportTok}[1]{#1}
\newcommand{\InformationTok}[1]{\textcolor[rgb]{0.56,0.35,0.01}{\textbf{\textit{#1}}}}
\newcommand{\KeywordTok}[1]{\textcolor[rgb]{0.13,0.29,0.53}{\textbf{#1}}}
\newcommand{\NormalTok}[1]{#1}
\newcommand{\OperatorTok}[1]{\textcolor[rgb]{0.81,0.36,0.00}{\textbf{#1}}}
\newcommand{\OtherTok}[1]{\textcolor[rgb]{0.56,0.35,0.01}{#1}}
\newcommand{\PreprocessorTok}[1]{\textcolor[rgb]{0.56,0.35,0.01}{\textit{#1}}}
\newcommand{\RegionMarkerTok}[1]{#1}
\newcommand{\SpecialCharTok}[1]{\textcolor[rgb]{0.81,0.36,0.00}{\textbf{#1}}}
\newcommand{\SpecialStringTok}[1]{\textcolor[rgb]{0.31,0.60,0.02}{#1}}
\newcommand{\StringTok}[1]{\textcolor[rgb]{0.31,0.60,0.02}{#1}}
\newcommand{\VariableTok}[1]{\textcolor[rgb]{0.00,0.00,0.00}{#1}}
\newcommand{\VerbatimStringTok}[1]{\textcolor[rgb]{0.31,0.60,0.02}{#1}}
\newcommand{\WarningTok}[1]{\textcolor[rgb]{0.56,0.35,0.01}{\textbf{\textit{#1}}}}
\usepackage{longtable,booktabs,array}
\newcounter{none} % for unnumbered tables
\usepackage{calc} % for calculating minipage widths
% Correct order of tables after \paragraph or \subparagraph
\usepackage{etoolbox}
\makeatletter
\patchcmd\longtable{\par}{\if@noskipsec\mbox{}\fi\par}{}{}
\makeatother
% Allow footnotes in longtable head/foot
\IfFileExists{footnotehyper.sty}{\usepackage{footnotehyper}}{\usepackage{footnote}}
\makesavenoteenv{longtable}
\usepackage{graphicx}
\makeatletter
\newsavebox\pandoc@box
\newcommand*\pandocbounded[1]{% scales image to fit in text height/width
  \sbox\pandoc@box{#1}%
  \Gscale@div\@tempa{\textheight}{\dimexpr\ht\pandoc@box+\dp\pandoc@box\relax}%
  \Gscale@div\@tempb{\linewidth}{\wd\pandoc@box}%
  \ifdim\@tempb\p@<\@tempa\p@\let\@tempa\@tempb\fi% select the smaller of both
  \ifdim\@tempa\p@<\p@\scalebox{\@tempa}{\usebox\pandoc@box}%
  \else\usebox{\pandoc@box}%
  \fi%
}
% Set default figure placement to htbp
\def\fps@figure{htbp}
\makeatother
\ifLuaTeX
\usepackage[bidi=basic,shorthands=off]{babel}
\else
\usepackage[bidi=default,shorthands=off]{babel}
\fi
\ifPDFTeX
\else
\babelfont{rm}[]{TeX Gyre Pagella}
\fi
\ifLuaTeX
  \usepackage{selnolig} % disable illegal ligatures
\fi
\setlength{\emergencystretch}{3em} % prevent overfull lines
\providecommand{\tightlist}{%
  \setlength{\itemsep}{0pt}\setlength{\parskip}{0pt}}
% LaTeX preamble additions for the manual (pandoc --include-in-header).
\usepackage{xcolor}
\usepackage{tabularx}
\usepackage{longtable}
\usepackage{booktabs}
\usepackage{float}
\usepackage{etoolbox}
% Code blocks: a shaded box with a colored left rule.
\definecolor{shadecolor}{HTML}{F5F3EE}
\definecolor{rulecolor}{HTML}{E89B3C}
\usepackage{framed}
\renewenvironment{Shaded}{%
  \def\FrameCommand{{\color{rulecolor}\vrule width 2pt}\hspace{6pt}\colorbox{shadecolor}}%
  \MakeFramed{\advance\hsize-\width\FrameRestore}\small}%
  {\endMakeFramed}
% Tables: smaller font, a little breathing room.
\AtBeginEnvironment{longtable}{\small}
\setlength{\LTpre}{6pt}
\setlength{\LTpost}{6pt}
% Figures stay where they are written, and a tall one leaves room for its
% caption: pandoc's bounding macro is redefined to cap the height at 86 %
% of the text height instead of all of it.
\makeatletter
\def\fps@figure{H}
\AtBeginDocument{%
  \renewcommand*\pandocbounded[1]{%
    \sbox\pandoc@box{#1}%
    \Gscale@div\@tempa{0.86\textheight}{\dimexpr\ht\pandoc@box+\dp\pandoc@box\relax}%
    \Gscale@div\@tempb{\linewidth}{\wd\pandoc@box}%
    \ifdim\@tempb\p@<\@tempa\p@\let\@tempa\@tempb\fi
    \ifdim\@tempa\p@<\p@\scalebox{\@tempa}{\usebox\pandoc@box}%
    \else\usebox{\pandoc@box}%
    \fi}}
\makeatother
% Inline code wraps instead of overflowing the margin.
\usepackage{xurl}
\makeatletter
\renewcommand{\texttt}[1]{{\ttfamily\hyphenchar\font=-1 #1}}
\makeatother
\usepackage{bookmark}
\IfFileExists{xurl.sty}{\usepackage{xurl}}{} % add URL line breaks if available
\urlstyle{same}
\hypersetup{
  pdftitle={Cooperatio Manual},
  pdflang={en-US},
  colorlinks=true,
  linkcolor={NavyBlue},
  filecolor={Maroon},
  citecolor={Blue},
  urlcolor={NavyBlue},
  pdfcreator={LaTeX via pandoc}}

\title{Cooperatio Manual}
\usepackage{etoolbox}
\makeatletter
\providecommand{\subtitle}[1]{% add subtitle to \maketitle
  \apptocmd{\@title}{\par {\large #1 \par}}{}{}
}
\makeatother
\subtitle{Architecture, capabilities and purpose of the method and of
its instrument --- method version 0.6.0}
\author{}
\date{}

\begin{document}
\maketitle

{
\hypersetup{linkcolor=black}
\setcounter{tocdepth}{1}
\tableofcontents
}
\chapter*{Preface}\label{preface}
\addcontentsline{toc}{chapter}{Preface}

This manual describes \textbf{cooperatio}, a method for organizing a
team of people, language-model agents and code around a complex target,
and the material that distributes it: the principle, the reflexive
application (the method applied to itself), the validator and its
sensitivity test, the agent skill that scaffolds a production cell, and
the registry of executed instances. It is written for three readers:

\begin{itemize}
\tightlist
\item
  \textbf{the adopter}, a person or an orchestrator agent who wants to
  set up a production cell for a target of their own and run it to a
  verdict;
\item
  \textbf{the maintainer}, who edits the material or the instrument and
  has to keep its provenance, ledgers and pins coherent;
\item
  \textbf{the reviewer or auditor}, who wants to know what the material
  claims, what evidence stands behind each claim, and what a clone of
  the material can and cannot verify.
\end{itemize}

It is a manual, not a replacement for the material: every rule stated
here is stated with more precision in the file it comes from, and the
manual names that file. Where the manual and the material disagree, the
material wins; please report the disagreement as a note.

\textbf{Version.} The manual describes method version 0.6.0, as recorded
in \texttt{VERSION}. The counts it quotes (gates, perturbations, tests,
files) were measured on that version; later versions change them, and
the ``Migration note'' of each entry in
\texttt{CHANGELOG.\allowbreak{}md} says what a cell has to change.

\textbf{Conventions.} File names and identifiers appear in
\texttt{monospace}; paths are relative to the root of the material, or
to the folder of the reflexive application
(\texttt{02-\allowbreak{}reflexive-\allowbreak{}application/\allowbreak{}})
where the context says so. A term in \textbf{bold} at its first use is a
term of the canon
(\texttt{02-\allowbreak{}reflexive-\allowbreak{}application/\allowbreak{}01-\allowbreak{}canon.\allowbreak{}md})
and is defined in the glossary (appendix A). Claims the material itself
labels as hypotheses are quoted with their label, \texttt{[hypothesis]},
because the material's own rule applies to the manual: a claim that the
method works carries its evidence or its label.

\textbf{Diagrams.} Each diagram is written twice: as a Mermaid block in
the Markdown source, which Markdown viewers that support Mermaid render
in place, and as a Graphviz file under \texttt{diagrams/\allowbreak{}},
rendered to SVG and PDF. The PDF edition uses the Graphviz renders; the
two are drawn from the same description and differ only in layout.

\textbf{Status of the method, in the material's own words.} A working
hypothesis, with a documentary validator; the gates of a target's domain
enter as results the cell's own code produces and the validator checks;
two executed instances, one of them auditable from the material. The
manual repeats this on purpose wherever a reader could take a described
capability for a measured one.

\part{Purpose and scope}

\chapter{What cooperatio is}\label{what-cooperatio-is}

\section{The problem it addresses}\label{the-problem-it-addresses}

A complex target --- a vehicle, a book, a software library, a long
document --- is produced today by teams in which people, language models
and code all take part. The failure modes of such a team are familiar:
one very capable agent does everything and checks its own work; a
request is mistaken for a result (``I asked for medium quality, so it is
medium quality''); the same number is stated three ways in three places;
a reviewer invents a new grid every round, so nobody can say whether
version 4 is better than version 3; an error that was found once comes
back because nothing was built to detect it; and the work never ends,
because every review finds something.

Cooperatio is a method that answers each of these with a named component
and a rule, and a material that applies the method to itself so that the
rules can be measured rather than trusted.

\section{The idea in one sentence}\label{the-idea-in-one-sentence}

The principle (\texttt{01-\allowbreak{}principle.\allowbreak{}md},
section 1) states it thus:

\begin{quote}
Producing a complex artifact is compiling a graph of typed artifacts,
each with an owner, a contract and a gate, in which separate roles
communicate only through artifacts and notes, quality is measured by
what the system delivered and not by what was asked for, each review is
made by someone who did not produce, with a rubric frozen before it, so
that rounds are comparable, and every error that escapes becomes a new
gate, until a declared stopping criterion.
\end{quote}

Put another way: it replaces ``one very capable agent does everything''
with ``a small organization, with explicit contracts, in which
intelligence is raw material and code is the glue that composes, checks
and remembers''.

The organization assembled for one target is a \textbf{production cell}.
Its parts are typed \textbf{artifacts} (files with one \textbf{owner}, a
\textbf{contract}, dependencies and \textbf{gates}), narrow
\textbf{roles} described by \textbf{role cards}, a \textbf{canon} that
is the single source of numbers and terms, a \textbf{validator} that
reads every artifact against the canon, an \textbf{independent review}
with a \textbf{rubric} frozen by hash, \textbf{notes} as the only change
request between roles, and a \textbf{stopping criterion} declared by the
owner before the last round.

\section{What the material
distributes}\label{what-the-material-distributes}

The material is a folder under version control. At version 0.6.0 it
holds:

{\def\LTcaptype{none} % do not increment counter
\begin{longtable}[]{@{}
  >{\raggedright\arraybackslash}p{(\linewidth - 4\tabcolsep) * \real{0.3333}}
  >{\raggedright\arraybackslash}p{(\linewidth - 4\tabcolsep) * \real{0.3333}}
  >{\raggedright\arraybackslash}p{(\linewidth - 4\tabcolsep) * \real{0.3333}}@{}}
\toprule\noalign{}
\begin{minipage}[b]{\linewidth}\raggedright
Part
\end{minipage} & \begin{minipage}[b]{\linewidth}\raggedright
Path
\end{minipage} & \begin{minipage}[b]{\linewidth}\raggedright
What it is for
\end{minipage} \\
\midrule\noalign{}
\endhead
\bottomrule\noalign{}
\endlastfoot
The principle & \texttt{01-\allowbreak{}principle.\allowbreak{}md} & the
method: idea, scheme, eleven components, ten golden rules, what changes
by domain, the recipe in twelve steps, warning signs, limits, what would
promote or withdraw a claim, ontological assumptions \\
The reflexive application &
\texttt{02-\allowbreak{}reflexive-\allowbreak{}application/\allowbreak{}}
& the method applied to itself: a brief, a canon, an artifact graph,
four role cards, a gate table with a frozen rubric, a generic recipe in
fifteen steps, three designed instances, the experiments about the
method, the review of its first round, the lessons of every round, and
the instrument \\
The instrument &
\texttt{02-\allowbreak{}reflexive-\allowbreak{}application/\allowbreak{}80-\allowbreak{}validator/\allowbreak{}}
& \texttt{validator.\allowbreak{}py} (the gates and the self-test),
\texttt{sensitivity.\allowbreak{}py} (the gate on the gates),
\texttt{perturbation\_\allowbreak{}set.\allowbreak{}py} (the known
defects), \texttt{test\_\allowbreak{}validator.\allowbreak{}py} (the
pytest suite), and the two published reports \\
The contracts &
\texttt{02-\allowbreak{}reflexive-\allowbreak{}application/\allowbreak{}05-\allowbreak{}contracts/\allowbreak{}}
& six JSON Schemas: note, state of a note, instance registry, round
record, domain result, owner decision \\
The skill &
\texttt{skills/\allowbreak{}cooperatio-\allowbreak{}teamwork/\allowbreak{}}
& an agent skill that creates the teamwork files for a target, with a
scaffold script, templates and a review checklist \\
The instances & \texttt{instances/\allowbreak{}} &
\texttt{registry.\allowbreak{}json}, the record of the known instances
pinned by sha256, and
\texttt{004-\allowbreak{}adoption-\allowbreak{}guide/\allowbreak{}}, the
cell of the second executed instance, distributed as a record \\
The owner records & \texttt{VERSION},
\texttt{CHANGELOG.\allowbreak{}md}, \texttt{LICENSE},
\texttt{.\allowbreak{}gitattributes} & the version of the method, its
history with a migration note per version, the Apache License 2.0, and
the rule that every text file checks out with LF line breaks \\
\end{longtable}
}

\begin{figure}
\centering
\pandocbounded{\includegraphics[keepaspectratio,alt={The folders of the material and what each one holds.}]{diagrams/material-map.pdf}}
\caption{The folders of the material and what each one
holds.}\label{fig-material-map}
\end{figure}

\section{Who it is for}\label{who-it-is-for}

The brief of the reflexive application
(\texttt{02-\allowbreak{}reflexive-\allowbreak{}application/\allowbreak{}00-\allowbreak{}brief.\allowbreak{}md})
names two readers: whoever is going to instantiate the method on a new
target without having followed the first instance, and an orchestrator
agent that receives the recipe as an instruction. Both have to be able
to execute without asking what a term means --- which is why the canon
exists and why every document links to it.

The method is for a team, not for a single executor. With a single agent
playing every role the separation is weak: it writes the rubric that
will judge it. The material says so in its limits
(\texttt{01-\allowbreak{}principle.\allowbreak{}md}, section 9) and in
the skill: where there is only one agent, at least the rubric and the
review have to come from outside the producing run.

\section{The status of the evidence}\label{the-status-of-the-evidence}

The material calls itself a working hypothesis and says exactly what
that covers. The reader of this manual should carry three facts
throughout:

\begin{enumerate}
\def\labelenumi{\arabic{enumi}.}
\tightlist
\item
  \textbf{Two executed instances, both documentary, both by the same
  owner.} Instance \texttt{000-\allowbreak{}av-\allowbreak{}production},
  an audiovisual production method held by the owner and not
  distributable, from which the method was generalized; and instance
  \texttt{004-\allowbreak{}adoption-\allowbreak{}guide}, the adoption
  guide of the method itself, whose cell is distributed under
  \texttt{instances/\allowbreak{}} and whose method document is pinned
  by sha256. The three other instances in the registry (a car, a book
  about 2027, an HPC library) are designed sketches, not executions.
\item
  \textbf{One part of the scheme has never been exercised: the
  recompilation by hash.} In neither execution did an engine rebuild a
  stale artifact from the hash of its inputs; the validator resolves the
  command an artifact declares against the paths it names and does not
  run it.
\item
  \textbf{Efficacy is one-armed.} Everything the method found, it found
  under the method; no production of the same target ran without the
  gates, so there is no escape rate to compare. The counterfactual is
  experiment E2, designed and not run; the cheaper seeded-escape design
  is E4, designed and awaiting the owner's budget (chapter 10).
\end{enumerate}

The principle keeps a table, written before any data, of what would
promote each kind of claim to a result and what would withdraw it
(section 9, ``What would promote or withdraw a claim''); the manual
reproduces it in chapter 10.

\chapter{When to use it, and when
not}\label{when-to-use-it-and-when-not}

\section{The four conditions}\label{the-four-conditions}

The method holds where the target has
(\texttt{01-\allowbreak{}principle.\allowbreak{}md}, opening):

\begin{itemize}
\tightlist
\item
  \textbf{(a)} a brief that an owner can declare, with a stopping
  criterion;
\item
  \textbf{(b)} a discrete artifact, assemblable from intermediate parts;
\item
  \textbf{(c)} sources of truth separable per content type;
\item
  \textbf{(d)} dependencies between those parts that can be kept
  acyclic.
\end{itemize}

The ontological assumptions behind the four (section 10 of the
principle) tell which kinds of target fall outside:

{\def\LTcaptype{none} % do not increment counter
\begin{longtable}[]{@{}
  >{\raggedright\arraybackslash}p{(\linewidth - 4\tabcolsep) * \real{0.3333}}
  >{\raggedright\arraybackslash}p{(\linewidth - 4\tabcolsep) * \real{0.3333}}
  >{\raggedright\arraybackslash}p{(\linewidth - 4\tabcolsep) * \real{0.3333}}@{}}
\toprule\noalign{}
\begin{minipage}[b]{\linewidth}\raggedright
Assumption
\end{minipage} & \begin{minipage}[b]{\linewidth}\raggedright
A target that violates it
\end{minipage} & \begin{minipage}[b]{\linewidth}\raggedright
What happens
\end{minipage} \\
\midrule\noalign{}
\endhead
\bottomrule\noalign{}
\endlastfoot
the target is a discrete, assemblable artifact with a done moment & a
continuous process or an ongoing service & the stopping criterion would
be a fiction; a service is run, not compiled \\
the target decomposes into parts with acyclic dependencies & an
exploratory target, whose next question is decided by the last result &
it is a loop, not a graph; the method degrades to peer review with a
rubric \\
quality is measurable by instruments that do not include the producer &
a target where only taste decides & every gate becomes a contract and
the independent review is the only instrument \\
the knowledge the target needs exists, or is produced by declared roles
& an emergent target, whose content exists only after the audience
reacts & its canon would be written after the fact, which is not a
canon \\
\end{longtable}
}

\section{Four situations that cost
something}\label{four-situations-that-cost-something}

Inside the four conditions, the principle declares four recurring
situations where the method still runs but loses part of what it
promises (section 9, ``When the method fits, and what each gap costs'').
The table is a declaration of scope written before any of the cases was
measured \texttt{[hypothesis]}:

{\def\LTcaptype{none} % do not increment counter
\begin{longtable}[]{@{}
  >{\raggedright\arraybackslash}p{(\linewidth - 4\tabcolsep) * \real{0.3333}}
  >{\raggedright\arraybackslash}p{(\linewidth - 4\tabcolsep) * \real{0.3333}}
  >{\raggedright\arraybackslash}p{(\linewidth - 4\tabcolsep) * \real{0.3333}}@{}}
\toprule\noalign{}
\begin{minipage}[b]{\linewidth}\raggedright
Situation
\end{minipage} & \begin{minipage}[b]{\linewidth}\raggedright
What is lost
\end{minipage} & \begin{minipage}[b]{\linewidth}\raggedright
What it costs to run anyway
\end{minipage} \\
\midrule\noalign{}
\endhead
\bottomrule\noalign{}
\endlastfoot
a single provider of models (every role from one family) & independence
of family between producer and reviewer; what stays is independence of
session and of reasoning, and the frozen rubric & record the relation in
the round record (\texttt{reviewer\_\allowbreak{}relation}), and get the
rubric, or a reading review, from outside the run \\
a target without a registered prior & the facts about a prior that
\texttt{p-\allowbreak{}prior} checks, and a recipe already executed once
& the canon starts empty of recorded facts, and the first execution
becomes a prior; budget an extra round for the gates that do not fit \\
a short target (a page, a note) & proportion: the cost per round is
fixed and the target is small & below a few pages the method is dead
weight; a checklist and one review cost less \\
a target without an instrument of its domain & the domain gates: only
the documentary gates run, so the method degrades to a review with a
rubric & the owner writes the domain instrument first, or declares in
the brief which non-negotiable no gate measures, and the verdict says
so \\
\end{longtable}
}

\section{What it costs}\label{what-it-costs}

Per round: one production, one independent review and the maintenance of
the instruments. The only costs measured so far are those of review
runs, in experiments E1 and E3 of
\texttt{02-\allowbreak{}reflexive-\allowbreak{}application/\allowbreak{}60-\allowbreak{}experiments.\allowbreak{}md}:

{\def\LTcaptype{none} % do not increment counter
\begin{longtable}[]{@{}
  >{\raggedright\arraybackslash}p{(\linewidth - 8\tabcolsep) * \real{0.2000}}
  >{\raggedright\arraybackslash}p{(\linewidth - 8\tabcolsep) * \real{0.2000}}
  >{\raggedright\arraybackslash}p{(\linewidth - 8\tabcolsep) * \real{0.2000}}
  >{\raggedright\arraybackslash}p{(\linewidth - 8\tabcolsep) * \real{0.2000}}
  >{\raggedright\arraybackslash}p{(\linewidth - 8\tabcolsep) * \real{0.2000}}@{}}
\toprule\noalign{}
\begin{minipage}[b]{\linewidth}\raggedright
Experiment
\end{minipage} & \begin{minipage}[b]{\linewidth}\raggedright
Executor
\end{minipage} & \begin{minipage}[b]{\linewidth}\raggedright
Input tokens
\end{minipage} & \begin{minipage}[b]{\linewidth}\raggedright
Output tokens
\end{minipage} & \begin{minipage}[b]{\linewidth}\raggedright
Wall time
\end{minipage} \\
\midrule\noalign{}
\endhead
\bottomrule\noalign{}
\endlastfoot
E1 (with the rubric) & X, the producer's family & about 58,000 (most of
them the prior's method document) & 55,000 to 66,000, 90 to 98 \% of
them reasoning & 8 to 10 minutes \\
E1 & Y, another family & not separable per call & not separable & about
4 minutes, indicative \\
E3 (without the rubric and the card) & X & about 56,000 & 41,000 to
46,000 & 10 to 13 minutes \\
E3 & Y & not separable & not separable & about 4 minutes, indicative \\
\end{longtable}
}

The material is careful about these numbers: minutes taken in two
sessions of different throughput are not comparable, so the measure is
the tokens of each call, and a wall time is reported only beside the
throughput of the same run. The cost of a production round and of the
instruments has not been measured.

The method pays off when being wrong is expensive \texttt{[hypothesis]};
stated as a condition still unmeasured, the probability of an error
times its cost has to exceed the cost per round times the rounds until
the stop. For a throwaway draft the recipe prescribes only steps 1, 3
and 11 (brief, canon, one version and a read from outside).

\section{What already exists, and what this
adds}\label{what-already-exists-and-what-this-adds}

The method borrows from established practice and says so, so that a
reader can tell the vocabulary from the claim
(\texttt{01-\allowbreak{}principle.\allowbreak{}md}, section 9):

{\def\LTcaptype{none} % do not increment counter
\begin{longtable}[]{@{}
  >{\raggedright\arraybackslash}p{(\linewidth - 4\tabcolsep) * \real{0.3333}}
  >{\raggedright\arraybackslash}p{(\linewidth - 4\tabcolsep) * \real{0.3333}}
  >{\raggedright\arraybackslash}p{(\linewidth - 4\tabcolsep) * \real{0.3333}}@{}}
\toprule\noalign{}
\begin{minipage}[b]{\linewidth}\raggedright
Here
\end{minipage} & \begin{minipage}[b]{\linewidth}\raggedright
Established practice
\end{minipage} & \begin{minipage}[b]{\linewidth}\raggedright
What this adds, if anything
\end{minipage} \\
\midrule\noalign{}
\endhead
\bottomrule\noalign{}
\endlastfoot
artifact graph, rebuild by hash & build systems & the graph is of
documents and decisions too, and an owner decision is a node \\
gate with magnitude and limit & quality gates of continuous integration;
acceptance criteria & the limit carries its justification and origin
label, and a gate counts only after its perturbation \\
gate sensitivity test & mutation testing; fault injection & a condition
of every gate, recorded with its provenance; a noisy perturbation
fails \\
frozen rubric, mechanical verdict & review checklists; scoring rubrics &
frozen by hash before the review, so two rounds are comparable \\
notes routed to their origin & issue trackers; defect triage & a note
aims at the artifact where the defect was born; an escaped regression
becomes a gate or a reason \\
role cards & RACI matrices; job descriptions & the ``does not decide''
field, and the code-conversion degree per role \\
failure analysis per gate & FMEA & done for the validator's gates, for
one failure mode, the false green (chapter 13) \\
\end{longtable}
}

\chapter{The material at a glance}\label{the-material-at-a-glance}

\section{Where to start}\label{where-to-start}

The root \texttt{README.\allowbreak{}md} gives three entry points:

\begin{itemize}
\tightlist
\item
  to understand the method, read
  \texttt{01-\allowbreak{}principle.\allowbreak{}md};
\item
  to apply it to a new target, read
  \texttt{02-\allowbreak{}reflexive-\allowbreak{}application/\allowbreak{}10-\allowbreak{}generic-\allowbreak{}recipe.\allowbreak{}md};
  the three instances there show what an artifact, a gate and a rubric
  are in each domain;
\item
  to have an agent create the files, install the skill in
  \texttt{skills/\allowbreak{}cooperatio-\allowbreak{}teamwork/\allowbreak{}}
  and ask it to set up teamwork for the target.
\end{itemize}

Its ``ten-minute path'' is reproduced, with the commands verified, in
chapter 17.

\section{What an adopter carries, and what is
research}\label{what-an-adopter-carries-and-what-is-research}

The material separates the two. The fixed cost of a cell is the set of
files the scaffold writes (24 with two roles; chapter 17 lists them),
among them the instrument it copies: three Python scripts the adopter
runs and does not have to read, about 5,700 lines at version 0.6.0
(\texttt{validator.\allowbreak{}py} 4,859,
\texttt{sensitivity.\allowbreak{}py} 550,
\texttt{perturbation\_\allowbreak{}set.\allowbreak{}py} 333, measured by
line count), plus one review per round. The history that matters to an
adopter is the ``Migration note'' of each version in
\texttt{CHANGELOG.\allowbreak{}md}; the rest of the changelog --- the
panels, the probes, the reworks --- is the research.

\section{Requirements}\label{requirements}

\begin{itemize}
\tightlist
\item
  Python 3.9 or later, standard library only, for the three scripts
  (they exit 2 below it; an interpreter older than 3.7 refuses the file
  before any line runs).
\item
  \texttt{pytest} for the test suites, which are run with
  \texttt{python3 -\allowbreak{}B -\allowbreak{}m pytest -\allowbreak{}p no:\allowbreak{}cacheprovider}
  so that no bytecode cache lands in the material.
\item
  A checkout with LF line breaks: \texttt{.\allowbreak{}gitattributes}
  enforces \texttt{* text eol=lf}, and every sha256 the instrument
  computes is taken over the text read as UTF-8, byte-order mark
  dropped, line breaks as LF, so a CRLF checkout keeps the pins.
\end{itemize}

\section{Versioning and the freeze
policy}\label{versioning-and-the-freeze-policy}

\texttt{VERSION} holds the semantic version of the method;
\texttt{CHANGELOG.\allowbreak{}md} has one entry per version, with a
migration note. The round record
(\texttt{02-\allowbreak{}reflexive-\allowbreak{}application/\allowbreak{}state/\allowbreak{}round-\allowbreak{}record.\allowbreak{}jsonl})
has one line per round, naming the executor of each role by
\textbf{family} letter and effort, never by model identifier.

After 0.6.0 the instrument is frozen: the validator, the sensitivity
test and the scaffold change only for a failure that changes a verdict,
or for a new execution. The next budget goes to measuring value
(experiment E4) and portability (an execution outside the documentary
domain by someone who is not the owner, experiment E5), not to more
instrument.

\section{License}\label{license}

Apache License 2.0; the full text is in \texttt{LICENSE}.

\part{The method}

\chapter{The scheme and the loop}\label{the-scheme-and-the-loop}

\section{The scheme}\label{the-scheme}

The principle writes the method as a scheme --- a naming of the parts
and of the loop, not a formula with a testable property
(\texttt{01-\allowbreak{}principle.\allowbreak{}md}, section 3). Let
\texttt{A} be the target, \texttt{B} the brief and \texttt{F} the
corpus:

\begin{Shaded}
\begin{Highlighting}[]
\NormalTok{A = Compile( Graph( Roles, Artifacts, Gates ), B, F )  until  Stop}

\NormalTok{where:}
\NormalTok{  Artifact   = (name, owner, schema, dependencies, paths, gates, hash)}
\NormalTok{  Role       = (responsibility, inputs, outputs, decides, does not decide,}
\NormalTok{                exit gate, notes it emits, code{-}conversion degree)}
\NormalTok{  Gate       = measurable function over an artifact that returns pass/reject}
\NormalTok{               with the magnitude measured, never only a signal}
\NormalTok{  Note       = (target\_artifact, severity, text, origin, evidence, round)}
\NormalTok{  Canon      = single source of numbers, names, terms and decisions, with a}
\NormalTok{                version per section and an owner per section}
\NormalTok{  Review     = independent reading of the product with a fixed rubric,}
\NormalTok{               comparable between rounds, that returns notes and a verdict}
\NormalTok{  Stop       = criterion declared by the owner before the last round}
\end{Highlighting}
\end{Shaded}

\section{The loop}\label{the-loop}

\begin{Shaded}
\begin{Highlighting}[]
\NormalTok{compile(graph, brief, corpus):}
\NormalTok{    canon ← extract(brief, corpus)                  \# numbers, terms, names, rules}
\NormalTok{    repeat:}
\NormalTok{        for each artifact in topological order:}
\NormalTok{            if inputs changed (hash) or there is a high/blocking note against it:}
\NormalTok{                output ← owner\_role.produce(inputs, canon, open\_notes)}
\NormalTok{                for each gate of the artifact:}
\NormalTok{                    measure ← gate(output)          \# magnitude, not only a signal}
\NormalTok{                    if rejected: note(artifact, severity, measure) ; return to owner}
\NormalTok{                validator.check(output, canon)  \# terms, numbers, names, assets}
\NormalTok{        product ← assemble(final artifacts)}
\NormalTok{        review ← independent\_reviewer(product, fixed\_rubric)}
\NormalTok{        for each finding of the review:}
\NormalTok{            if it is an instrument defect (false positive): fix the instrument at the root}
\NormalTok{            if it is a regression that no gate caught: create the gate, then fix}
\NormalTok{            otherwise: note to the originating artifact}
\NormalTok{        if stopping\_criterion: record open notes and exit}
\NormalTok{    until stop}
\end{Highlighting}
\end{Shaded}

\begin{figure}
\centering
\pandocbounded{\includegraphics[keepaspectratio,alt={The compile loop of the method (principle, section 3).}]{diagrams/method-scheme.pdf}}
\caption{The compile loop of the method (principle, section
3).}\label{fig-method-scheme}
\end{figure}

Three features of the loop decide most of what follows in this manual:

\begin{itemize}
\tightlist
\item
  \textbf{Only what is stale is rebuilt.} An artifact is \textbf{stale}
  when the hash of an input changed since it passed its gates, or when a
  high or blocking note targets it. This is the design; as said in
  chapter 1, no execution has yet exercised an engine that rebuilds by
  hash.
\item
  \textbf{A gate rejects with a magnitude.} The note it emits carries
  the measure, not a sign, so the owner of the artifact knows how far
  from the limit the output landed.
\item
  \textbf{Each finding of the review has exactly three destinations.}
  Instrument defect, escaped regression, or product defect. The first
  fixes the instrument at the root; the second creates a gate
  \emph{before} the fix; the third is a note to the artifact where the
  defect was born.
\end{itemize}

\section{The orchestrator}\label{the-orchestrator}

The orchestrator is whoever executes the recipe: it passes paths (never
summaries) to the roles, runs the engine and imports the notes. It is a
\textbf{step, not a role}: it decides nothing that a card assigns to a
role or to the owner
(\texttt{10-\allowbreak{}generic-\allowbreak{}recipe.\allowbreak{}md},
``The orchestrator in pseudocode'').

\chapter{The eleven components}\label{the-eleven-components}

The principle names eleven components that do not change from one domain
to another (section 4; section 6 says what does change). The last column
is the error each one prevents, with the example the material gives.

\begin{figure}
\centering
\pandocbounded{\includegraphics[keepaspectratio,alt={The eleven components of the method around the production cell.}]{diagrams/components.pdf}}
\caption{The eleven components of the method around the production
cell.}\label{fig-components}
\end{figure}

{\def\LTcaptype{none} % do not increment counter
\begin{longtable}[]{@{}
  >{\raggedright\arraybackslash}p{(\linewidth - 4\tabcolsep) * \real{0.3333}}
  >{\raggedright\arraybackslash}p{(\linewidth - 4\tabcolsep) * \real{0.3333}}
  >{\raggedright\arraybackslash}p{(\linewidth - 4\tabcolsep) * \real{0.3333}}@{}}
\toprule\noalign{}
\begin{minipage}[b]{\linewidth}\raggedright
Component
\end{minipage} & \begin{minipage}[b]{\linewidth}\raggedright
What it is
\end{minipage} & \begin{minipage}[b]{\linewidth}\raggedright
The error it prevents (example)
\end{minipage} \\
\midrule\noalign{}
\endhead
\bottomrule\noalign{}
\endlastfoot
\textbf{Artifact graph} & each intermediate output is a typed file with
an owner, schema, dependencies and hash; the core recompiles only what
changed & work redone for nothing and, worse, work \emph{not} redone
when the input changed \\
\textbf{Role cards} & each role has a card: what it decides, what it
does not decide, inputs, outputs, exit gate, notes it emits & one role
compensating for the defect of another; the ``does not decide'' field
returns the defect to its owner as a note \\
\textbf{Contracts (schemas)} & every artifact validates against a schema
before being accepted; the schema refuses what is ambiguous & an
ambiguous position key in a hand-off sheet; an unknown field accepted at
assembly instead of at validation \\
\textbf{Gates with magnitude} & a gate returns the measure, and the
limit is declared with a justification; ``greater than zero'' lets a
broken instrument pass & an offset measured against a reference that was
itself wrong \\
\textbf{Routed notes} & a change request only by note, with target,
severity, evidence and round; only high and blocking return the
artifact; deduplication by hash & loose conversation between roles, a
lost request, the same note three times \\
\textbf{Canon} & single source of numbers, names, terms to avoid,
accepted variants, rules; each section with an owner and a version & a
figure stated one way in one section and differently in another \\
\textbf{Validator} & a verifier that reads all artifacts against the
canon and emits notes by itself & an old version of an asset cited by
the canon while the sections used the new one \\
\textbf{Independent review with a fixed rubric} & another agent, with no
access to the producers' reasoning, measures the product with the same
grid each round; three verdicts & the team approving its own work;
rounds that cannot be compared \\
\textbf{Instrument fixed at the root} & a false positive becomes a fix
to the rule, with a test that reproduces the case, and the previous
round is redone & an automatic reviewer of which 10 of 12 notes were
false, in the second pilot of the first instance
\texttt{[prior:\allowbreak{}000-\allowbreak{}av-\allowbreak{}production]} \\
\textbf{Stopping criterion} & the owner declares when to stop: the
review of the last version becomes a list of open notes, not a new
version & endless rounds \\
\textbf{Gate sensitivity test} & before trusting a gate, inject a known
defect and check that it rejects & a gate that checked numbers while the
non-negotiable was about statements; the sensitivity test found a
validator that read line by line and missed numbers in broken
sentences \\
\end{longtable}
}

What changes by domain is: which roles there are, which artifacts, what
each gate measures, what the review rubric is, which sources of truth
exist for each type of content, and what code-conversion degree the
domain already allows. The principle's table of three instantiations
(car, text, HPC) calibrates what an ``artifact'', a ``gate'' and a
``rubric'' are in each; the manual summarizes the three designed
instances in chapter 16.

\chapter{Roles, artifacts and notes}\label{roles-artifacts-and-notes}

\section{Roles and role cards}\label{roles-and-role-cards}

A \textbf{role} is a narrow responsibility exercised by an agent, a
person or code, which communicates with other roles only through
artifacts and notes. A role that neither emits notes nor has a gate is a
step, not a role.

Every role has a \textbf{role card} with seven fixed fields, which the
gate \texttt{p-\allowbreak{}card} reads as bold labels:

{\def\LTcaptype{none} % do not increment counter
\begin{longtable}[]{@{}
  >{\raggedright\arraybackslash}p{(\linewidth - 2\tabcolsep) * \real{0.5000}}
  >{\raggedright\arraybackslash}p{(\linewidth - 2\tabcolsep) * \real{0.5000}}@{}}
\toprule\noalign{}
\begin{minipage}[b]{\linewidth}\raggedright
Field
\end{minipage} & \begin{minipage}[b]{\linewidth}\raggedright
What goes there
\end{minipage} \\
\midrule\noalign{}
\endhead
\bottomrule\noalign{}
\endlastfoot
\textbf{Decides} & what the role settles alone \\
\textbf{Does not decide} & the neighboring decisions the role would be
tempted to make, and whose they are \\
\textbf{Inputs} & the artifacts it reads \\
\textbf{Outputs} & the artifacts it writes \\
\textbf{Exit gate} & the gates its output has to pass \\
\textbf{Notes it emits} & the typical target and severity \\
\textbf{Code-conversion degree} & how much of the work the code does,
and the executor of the role \\
\end{longtable}
}

The \textbf{code-conversion degree} has three values: (a) converted, the
code does it and a person or a model decides exceptions; (b) partial
under gate, the code does or measures a part; (c) contract and gate
only, the role is exercised by an agent or a person and the code only
validates the output. The degree is explicit per role and is expected to
rise over time.

The reflexive application has four roles
(\texttt{02-\allowbreak{}reflexive-\allowbreak{}application/\allowbreak{}03-\allowbreak{}roles/\allowbreak{}}):

{\def\LTcaptype{none} % do not increment counter
\begin{longtable}[]{@{}
  >{\raggedright\arraybackslash}p{(\linewidth - 6\tabcolsep) * \real{0.2500}}
  >{\raggedright\arraybackslash}p{(\linewidth - 6\tabcolsep) * \real{0.2500}}
  >{\raggedright\arraybackslash}p{(\linewidth - 6\tabcolsep) * \real{0.2500}}
  >{\raggedright\arraybackslash}p{(\linewidth - 6\tabcolsep) * \real{0.2500}}@{}}
\toprule\noalign{}
\begin{minipage}[b]{\linewidth}\raggedright
Role
\end{minipage} & \begin{minipage}[b]{\linewidth}\raggedright
Mission
\end{minipage} & \begin{minipage}[b]{\linewidth}\raggedright
Degree
\end{minipage} & \begin{minipage}[b]{\linewidth}\raggedright
Executor
\end{minipage} \\
\midrule\noalign{}
\endhead
\bottomrule\noalign{}
\endlastfoot
writer & turn the principle into documents an agent can execute; fill
the instances & (c) & a language model \\
verifier & check facts about the priors; maintain the instrument; decide
false positives; own the note ledgers & (b) & code for the number checks
and the contracts, judgment for the false positives \\
runner & run the validator and route its notes; decide nothing beyond
the rules it executes (named ``validator'' until v0.6.0) & (a) & the
script \texttt{validator.\allowbreak{}py} \\
reviewer & read the whole product with the frozen rubric; return scores,
findings and a verdict by the mechanical rule & (c) & an executor that
did not take part, in a separate session, ideally of another family \\
\end{longtable}
}

\begin{figure}
\centering
\pandocbounded{\includegraphics[keepaspectratio,alt={Roles talk only through artifacts and notes; the four roles of the reflexive application.}]{diagrams/roles-flow.pdf}}
\caption{Roles talk only through artifacts and notes; the four roles of
the reflexive application.}\label{fig-roles-flow}
\end{figure}

\section{Executors and families}\label{executors-and-families}

Separating roles makes it possible to give each role a different
executor: a person, code, or a language model of a chosen kind. The
principle (section 2, ``Using several models across the roles'') gives
the rules:

\begin{itemize}
\tightlist
\item
  \textbf{Capability follows the role.} A mechanical, fully specified
  role takes a fast, inexpensive model; a role that designs or weighs
  hard trade-offs takes the strongest available. The gate, not the
  model's confidence, decides whether the output is accepted.
\item
  \textbf{Independence across families is recommended, not required}
  \texttt{[hypothesis]}. The review is best run on a model of another
  family than the producer's; the cost is a second provider. It is a
  recommendation because the one datum is against it as a rule: in
  experiment E1 the producer's own model, in another session, caught 2
  of 2 injected defects, as the model of another family did
  (\texttt{[against:\allowbreak{}E1]}). What a pilot did show is the
  part a team can adopt without a second provider: a frozen rubric with
  a mechanical verdict rule turned a finding into a verdict
  (\texttt{[experiment:\allowbreak{}E3]}).
\item
  \textbf{Code is not a model.} Whatever a gate can measure is measured
  by code, identically whichever model produced the artifact.
\item
  \textbf{Contracts make models swappable.} Replacing the model behind a
  role must not change any other file; if it does, the contract is too
  loose.
\item
  \textbf{The assignment is auditable, and anonymous.} The round record
  names the executor of each role by \textbf{family} --- a letter kept
  stable across the material --- and by effort, never by model
  identifier. The map from letter to family stays with the owner,
  outside the material, and since the rework of v0.6.0 each chained line
  of the round record carries a salted commitment to that map (chapter
  15). A route or an alias is not a family: whether two executors differ
  is checked on what each name maps to, never on the name.
\end{itemize}

The round record also states, line by line, the reviewer's relation to
the producer (\texttt{reviewer\_\allowbreak{}relation}:
\texttt{another family}, \texttt{same family, separate session},
\texttt{same model, separate route}, or \texttt{none}), so the
recommendation can be checked against what was done.

\section{Artifacts}\label{artifacts}

An \textbf{artifact} is an intermediate or final output, typed, with a
name, a single owner, paths, a schema, dependencies, gates and content
identified by hash. An owner decision is also an artifact: it has a
name, it is recorded as an open item for the owner, and the graph knows
what depends on it (golden rule 7).

The record of an artifact, in the JSON form a cell uses
(\texttt{03-\allowbreak{}artifacts.\allowbreak{}json}):

\begin{Shaded}
\begin{Highlighting}[]
\NormalTok{\{ "name": "...", "owner": "ROLE", "paths": ["..."], "depends\_on": ["..."],}
\NormalTok{  "schema": "05{-}contracts/....schema.json" or null, "command": "..." or null,}
\NormalTok{  "gates": ["..."], "resource": "code" | "agent" | "human" | "machine",}
\NormalTok{  "proposed": false \}}
\end{Highlighting}
\end{Shaded}

Three rules about the record matter to the validator:

\begin{itemize}
\tightlist
\item
  the \texttt{gates} list is the \textbf{one source of the map from gate
  to artifact}: the validator runs each gate on the files of the
  artifacts that list it, and the gate table has to mirror that
  assignment (chapter 13);
\item
  \texttt{paths} may carry globs, and since v0.6.0 each literal path has
  to exist and each glob has to match at least once
  (\texttt{p-\allowbreak{}declared});
\item
  \texttt{command}, when present, is resolved against the folder: every
  path it names has to exist, and the command is never run. Paths and
  commands may carry \texttt{\{version\}}, which takes the layout's
  \texttt{version-\allowbreak{}label}.
\end{itemize}

The \textbf{artifact graph} is the set of artifacts with the
dependencies between them; it has to be acyclic, and the validator
checks that on every pass (\texttt{p-\allowbreak{}references}).
``Depends on'' means: the artifact reads that input and becomes stale
when it changes. A note is not a dependency.

\section{Notes}\label{notes}

A \textbf{note} is the only change request between roles. It has six
fields, all required
(\texttt{05-\allowbreak{}contracts/\allowbreak{}notes.\allowbreak{}schema.\allowbreak{}json}):

\begin{Shaded}
\begin{Highlighting}[]
\NormalTok{\{ "target\_artifact": "name in the graph", "severity": "low|medium|high|blocking",}
\NormalTok{  "text": "what is wrong and what is expected", "origin": "role or instrument",}
\NormalTok{  "evidence": "file, excerpt, measurement", "round": N \}}
\end{Highlighting}
\end{Shaded}

\textbf{Severity} has four values. Only \texttt{high} and
\texttt{blocking} send the artifact back to its owner; \texttt{low} and
\texttt{medium} accumulate as debt that the stop records. A note aims at
the artifact where the defect was born, not where it was seen; the note
router refuses a target the graph does not know and deduplicates by
identity (target, severity, text, origin, evidence).

An \textbf{open note} is a note recorded and not resolved when the
production stops: written debt, not forgotten debt (golden rule 10).
Since v0.6.0 the state of a note lives in a ledger of its own,
append-only and chained (chapter 15).

\section{Sources of truth and the
canon}\label{sources-of-truth-and-the-canon}

Every type of content has one declared \textbf{source of truth}, and it
is never the writer (golden rule 2): a fact comes from the corpus or its
primary source; a number comes from the canon or from the instrument
that measured it; exact text comes from the original. The inventory is a
table with content type, the source that errs least, who consults it and
what is forbidden as a source
(\texttt{01-\allowbreak{}sources-\allowbreak{}of-\allowbreak{}truth.\allowbreak{}md}
in a cell).

The \textbf{canon} is the single source of numbers, names, terms to
prefer and to avoid, accepted variants and rules, with an owner and a
version per section. In the reflexive application it is a Markdown page
whose tables the validator reads directly
(\texttt{01-\allowbreak{}canon.\allowbreak{}md}); in a cell it is a JSON
file (\texttt{02-\allowbreak{}canon.\allowbreak{}json}) with the lists
the validator consumes. Everything that two roles could write
differently goes into it. The canon is itself an artifact: it has an
owner, a version and dependents, and the validator reads it on every
pass.

A \textbf{contract} is the schema an artifact has to validate before it
is accepted; it rejects the ambiguous before the expensive work (golden
rule 5). The material's own contracts are listed in appendix B.

\chapter{Gates, limits and the sensitivity
test}\label{gates-limits-and-the-sensitivity-test}

\section{What a gate is}\label{what-a-gate-is}

A \textbf{gate} is a measurable function over an artifact that returns
pass or fail together with the measured \textbf{magnitude}. It has a
\textbf{limit} and a written justification, and the justification starts
with an \textbf{origin label}:

{\def\LTcaptype{none} % do not increment counter
\begin{longtable}[]{@{}
  >{\raggedright\arraybackslash}p{(\linewidth - 2\tabcolsep) * \real{0.5000}}
  >{\raggedright\arraybackslash}p{(\linewidth - 2\tabcolsep) * \real{0.5000}}@{}}
\toprule\noalign{}
\begin{minipage}[b]{\linewidth}\raggedright
Label
\end{minipage} & \begin{minipage}[b]{\linewidth}\raggedright
Meaning
\end{minipage} \\
\midrule\noalign{}
\endhead
\bottomrule\noalign{}
\endlastfoot
\texttt{[method]} & the limit follows from the method itself (for
example, zero terms outside the canon) \\
\texttt{[brief]} & the limit was set by the owner in the brief \\
\texttt{[hypothesis]} & the limit is an assumption about the domain,
plausible and justified, to be confirmed in a real instance \\
\end{longtable}
}

A row of a gate table has five columns: Gate, Artifact, Magnitude,
Limit, Justification. The validator reads the justification's label and
counts its words (at least six, read from the
\texttt{p-\allowbreak{}limit} row itself), and compares the numbers of
the Limit cell with the limit it applies.

Three distinctions the material insists on:

\begin{itemize}
\tightlist
\item
  \textbf{A tool that does not run is a tool failure, never a pass or a
  fail.} The validator exits 2 for it; the exit codes are in chapter 19.
\item
  \textbf{A gate that was not measured is pending, neither a pass nor a
  fail.} A gate whose object is a list of the canon is pending when the
  list is empty, never ``0 of 0 passed''. A run with a pending the
  caller did not accept by name exits 4.
\item
  \textbf{A signal gate} --- one that only tests existence or ``greater
  than zero'' --- lets a broken instrument through and has to gain a
  magnitude limit on both sides: an absurd value also rejects, because
  it signals a broken instrument.
\end{itemize}

A number taken from an earlier execution is a fact about it, never a
limit a new one inherits: a \texttt{[hypothesis]} limit may cite a prior
as support, never as the limit itself.

\section{Documentary gates and domain
gates}\label{documentary-gates-and-domain-gates}

The validator that ships with the material measures documents: terms,
sections, links, form, numbers stated about registered priors, the
contracts and the chains of its ledgers (the \texttt{p-\allowbreak{}*}
gates, chapter 13). A gate of a target's domain --- the identity of a
library's results, a speedup, a mass --- is a \textbf{domain gate},
named \texttt{d-\allowbreak{}<name>} in the cell's gate table. The
cell's own code measures it and writes a result file in the contract
\texttt{domain-\allowbreak{}result.\allowbreak{}schema.\allowbreak{}json};
the validator reads the result, recomputes the sha256 of its inputs,
refuses it when an input changed since it was written, turns a
\texttt{fail} into a note and a missing result into a pending, never a
pass. The validator does not run the cell's code (chapter 21).

\section{The sensitivity test}\label{the-sensitivity-test}

Before trusting a gate, inject a known defect and check that the gate
rejects it; a gate that never rejects proves nothing. The material makes
this a condition of every gate: the \textbf{gate sensitivity test}
(\texttt{sensitivity.\allowbreak{}py}) copies the whole material,
injects one \textbf{perturbation} at a time, runs the validator on the
copy and checks that the expected (document, gate) row went from pass to
fail --- and that no other row changed. A perturbation that changes
another row is \emph{noisy} and fails the report like a missed one,
unless it declares the other gate it reaches by design. The report
carries the provenance of what it was made against, and the validator's
\texttt{p-\allowbreak{}sensitivity} gate refuses a report that the
instrument cannot have produced. Chapter 14 describes the test in full.

\section{The lifecycle of a gate}\label{the-lifecycle-of-a-gate}

A gate is born from a regression that escaped every gate (golden rule
4), together with its perturbation. It lives while its object exists. It
is retired deliberately, never by neglect: when its object is gone, or
when its rule moved into a contract that refuses the defect earlier. The
row of the gate table, its assignment in the artifact list and its
perturbation leave together, and the changelog records ``retired gate:
\texttt{<id>}'', the reason and the rounds the gate stayed silent. No
gate was retired between v0.4.0 and v0.6.0.

\begin{figure}
\centering
\pandocbounded{\includegraphics[keepaspectratio,alt={The lifecycle of a gate (gate semantics, ``Gate lifecycle'').}]{diagrams/gate-lifecycle.pdf}}
\caption{The lifecycle of a gate (gate semantics, ``Gate
lifecycle'').}\label{fig-gate-lifecycle}
\end{figure}

The instrument detects half a retirement: removing a gate and keeping
its perturbation makes the perturbation missed; removing the
perturbation of a gate that still runs fails the test that every gate
has one. Removing both and regenerating the sensitivity report passes
everything; only the declaration in the changelog records that
retirement.

\chapter{The independent review}\label{the-independent-review}

\section{Who reviews, with what}\label{who-reviews-with-what}

An \textbf{independent review} is a reading of the product by someone
who did not produce it and has no access to the producer's reasoning,
with the fixed rubric, returning notes and a verdict. The reviewer
receives the files and the rubric --- never drafts, conversation or
summaries --- and the dispatch records the effective output limit of the
executor and the reason its generation ended. An ending for length, or a
report that lacks a section of the format, is a tool failure and
produces no verdict; the review runs again. This rule came from
experiment E1, where one report stopped at a limit the dispatch had not
declared and carried a verdict no reader could verify.

\section{The rubric}\label{the-rubric}

The \textbf{rubric} is the fixed grid of dimensions, anchors and verdict
rule that every review uses in every round; it is written before the
first review and frozen by hash. The generic recipe (step 8) asks for
six to ten dimensions with \textbf{rubric anchors} for scores 1, 3 and
5; a severity scale of the findings with examples from the domain; a
mechanical verdict rule; and a report format.

The reflexive application's rubric
(\texttt{04-\allowbreak{}gates-\allowbreak{}and-\allowbreak{}rubric.\allowbreak{}md},
section ``Rubric'') has eight dimensions:

{\def\LTcaptype{none} % do not increment counter
\begin{longtable}[]{@{}
  >{\raggedright\arraybackslash}p{(\linewidth - 2\tabcolsep) * \real{0.5000}}
  >{\raggedright\arraybackslash}p{(\linewidth - 2\tabcolsep) * \real{0.5000}}@{}}
\toprule\noalign{}
\begin{minipage}[b]{\linewidth}\raggedright
Dimension
\end{minipage} & \begin{minipage}[b]{\linewidth}\raggedright
What is measured
\end{minipage} \\
\midrule\noalign{}
\endhead
\bottomrule\noalign{}
\endlastfoot
D1 Fidelity to the principle & the eleven components and the golden
rules applied without contradiction; nothing about the reference case
invented \\
D2 Recipe executability & each step says what it reads, creates, does
and when it is done \\
D3 Concreteness of instances & each element filled with domain content,
not text that would serve any target \\
D4 Gates and numbers & gates with magnitude; plausible limits;
justifications that say why that number; hypotheses marked \\
D5 Internal coherence & canon terms in their sense; the same number
equal everywhere; references close \\
D6 Reflexivity & the production followed its own method and is honest
where it did not \\
D7 Form & American English, plain Markdown, pseudocode in \texttt{text},
no path, date or machine, sizes in range \\
D8 Epistemic honesty & distinguishes method, prior instance and domain;
marks hypotheses \\
\end{longtable}
}

Its severity scale: \texttt{blocking} for an invented fact about the
prior instance, a contradiction with the principle that would change the
application, a missing mandatory document, or an instance that cannot be
executed; \texttt{high} for a gap that prevents an agent from executing
a step, a signal gate or unjustifiable limit, a term used in a sense
different from the canon's, or an instance element filled with generic
text; \texttt{medium} for a domain imprecision without a
\texttt{[hypothesis]} mark; \texttt{low} for style.

\section{The verdict rule}\label{the-verdict-rule}

The \textbf{verdict} has three values --- \texttt{rejected},
\texttt{accepted with conditions}, \texttt{accepted} --- and the rule
that produces it is mechanical:

\begin{Shaded}
\begin{Highlighting}[]
\NormalTok{if a blocking finding exists or some dimension has a score ≤ 2:}
\NormalTok{    verdict = rejected}
\NormalTok{else if no high finding and every dimension has a score ≥ 4:}
\NormalTok{    verdict = accepted}
\NormalTok{else:}
\NormalTok{    verdict = accepted with conditions}
\end{Highlighting}
\end{Shaded}

\begin{figure}
\centering
\pandocbounded{\includegraphics[keepaspectratio,alt={The mechanical verdict rule of the frozen rubric.}]{diagrams/verdict-rule.pdf}}
\caption{The mechanical verdict rule of the frozen
rubric.}\label{fig-verdict-rule}
\end{figure}

Experiment E3 is the measure of what this rule is worth: without the
rubric and the severity table, two reviewers of different families still
named both injected defects but classed an invented fact about a prior
as \texttt{high} and accepted the product with conditions; with them,
the same finding was \texttt{blocking} and the verdict was
\texttt{rejected} (\texttt{[experiment:\allowbreak{}E3]}). That the
severity table and the mechanical rule are the parts that did it is a
hypothesis the design does not separate \texttt{[hypothesis]}.

\section{The report format}\label{the-report-format}

\begin{enumerate}
\def\labelenumi{\arabic{enumi}.}
\tightlist
\item
  \textbf{Verdict}, with the application of the rule.
\item
  \textbf{Scores by dimension}: table with dimension, score, evidence.
\item
  \textbf{Findings by severity}: each with target artifact (a name from
  the graph, or \texttt{corpus}), severity, text, evidence.
\item
  \textbf{Verified and correct.}
\item
  \textbf{Not verified}, and why.
\item
  \textbf{Out of scope.}
\item
  \textbf{Notes in JSON}, in the note contract (\texttt{origin} =
  \texttt{reviewer}).
\end{enumerate}

From the second round on, the report adds the item-by-item table of the
previous round (resolved, partial, remained, worsened).

\section{Freezing, and who writes the
rubric}\label{freezing-and-who-writes-the-rubric}

The rubric is frozen with
\texttt{validator.\allowbreak{}py -\allowbreak{}-\allowbreak{}freeze-\allowbreak{}rubric FILE},
which writes the sha256 of the rubric as the layout resolves it (a
dedicated \texttt{07-\allowbreak{}rubric.\allowbreak{}md} in a cell; the
``Rubric'' section of the gates file in the corpus). The command
\textbf{refuses to overwrite an existing pin}: replacing a frozen rubric
is a deliberate act via
\texttt{-\allowbreak{}-\allowbreak{}re-\allowbreak{}freeze-\allowbreak{}rubric},
with the reason recorded in the changelog, because re-freezing after a
review would make the verdicts incomparable.

The provenance of the rubric matters. In the reflexive application the
rubric was written by the producing agent in the first round, against
step 8 of the recipe, which gives it to the owner and the reviewer; the
deviation is recorded in the lessons, and the round record names the
author of the rubric (\texttt{rubric\_\allowbreak{}author}). A rubric
written by the producer is recorded as such --- a gap, not a neutral
fact --- and the validator adds a low note when a chained round line has
\texttt{rubric\_\allowbreak{}applied:\allowbreak{} true} and
\texttt{rubric\_\allowbreak{}author:\allowbreak{} producer}.

\chapter{The golden rules and the warning
signs}\label{the-golden-rules-and-the-warning-signs}

\section{The ten golden rules}\label{the-ten-golden-rules}

From \texttt{01-\allowbreak{}principle.\allowbreak{}md}, section 5, each
with its reason in short:

\begin{enumerate}
\def\labelenumi{\arabic{enumi}.}
\tightlist
\item
  \textbf{Measure what the system delivered, not what was asked for.}
  The effective value comes from a named instrument applied to the
  produced artifact, never from the parameter that requested it.
\item
  \textbf{Every type of content has one declared source of truth, and it
  is never the writer.} A number taken from an earlier execution is a
  fact about it, never a limit a new one inherits.
\item
  \textbf{Every false positive is an instrument defect until proven
  otherwise.} The easy path is to loosen the limit; the right path is to
  understand why the instrument read wrongly. A third way out is
  legitimate when the instrument cannot accept exceptions: work around
  it and record the workaround as a note, so that it does not become a
  silent rule.
\item
  \textbf{Every regression that escaped all the gates becomes a gate, or
  a recorded decision not to make one.} Before fixing the defect, create
  the detector. Since v0.6.0 the state ledger records which
  (\texttt{escaped}, with \texttt{gate} or
  \texttt{no\_\allowbreak{}gate}), and the validator refuses an escaped
  regression closed with neither.
\item
  \textbf{The contract refuses ambiguity before the expensive work
  starts.}
\item
  \textbf{Dates, machine names and absolute paths do not enter a durable
  artifact} if anonymity or reproducibility matters. The checking date
  is a command-line parameter, not content.
\item
  \textbf{A human decision is an artifact too.} What belongs to the
  owner is not invented.
\item
  \textbf{Convert judgment into code only when the criterion is
  measurable and the answer is unique.} The rest stays with an agent
  under a gate, or with a human under a contract.
\item
  \textbf{Version is method state; date of making is not.}
\item
  \textbf{What the review points out and the round does not fix is
  written down.}
\end{enumerate}

\section{Signs that an instance is
wrong}\label{signs-that-an-instance-is-wrong}

The principle (section 8) and the recipe (``Signs that the instance is
wrong, and which step catches them'') give the same list from two sides:

{\def\LTcaptype{none} % do not increment counter
\begin{longtable}[]{@{}
  >{\raggedright\arraybackslash}p{(\linewidth - 2\tabcolsep) * \real{0.5000}}
  >{\raggedright\arraybackslash}p{(\linewidth - 2\tabcolsep) * \real{0.5000}}@{}}
\toprule\noalign{}
\begin{minipage}[b]{\linewidth}\raggedright
Sign
\end{minipage} & \begin{minipage}[b]{\linewidth}\raggedright
Step of the recipe that should have caught it
\end{minipage} \\
\midrule\noalign{}
\endhead
\bottomrule\noalign{}
\endlastfoot
a role approves its own artifact without a gate & 5 (exit gate) and 7 \\
the review of round N is not comparable with that of N−1 & 8 (frozen
rubric) \\
a finding was ``resolved'' by loosening a limit & 12 (root-cause fix) \\
two artifacts state different numbers for the same fact & 3 and 9 \\
the owner gets a question the brief already answered, or a role decided
what its card said ``does not decide'' & 1 and 5 \\
the round gets longer and the notes do not fall & 12 (origin, not
symptom) \\
someone redid by hand what the graph should have recompiled & 4 and 10
(hash) \\
\end{longtable}
}

\chapter{Evidence and epistemic
discipline}\label{evidence-and-epistemic-discipline}

The material applies its own doctrine of evidence to itself. Four
mechanisms carry it: the registry of prior instances with its marks, the
labels a claim of efficacy has to carry, the experiments with their
provenance, and the table of what would promote or withdraw a claim.

\section{\texorpdfstring{Prior instances and the
\texttt{[prior:\allowbreak{}<id>]}
mark}{Prior instances and the  mark}}\label{prior-instances-and-the-priorid-mark}

A \textbf{prior instance} is an instance of the method that ran before
the current one and is registered in the \textbf{instance registry}
(\texttt{instances/\allowbreak{}registry.\allowbreak{}json}). A
statement about a prior carries \texttt{[prior:\allowbreak{}<id>]}, and
a number inside the marked unit is checked against that prior's method
document (\texttt{p-\allowbreak{}prior}). A statement that uses one of
the canon's prior-statement trigger phrases (``reference case'', ``prior
instance'', ``of the case'', ``the executed instance'', ``instance 0'')
without the mark is a finding of
\texttt{p-\allowbreak{}prior-\allowbreak{}mark}.

The registry at 0.6.0:

{\def\LTcaptype{none} % do not increment counter
\begin{longtable}[]{@{}
  >{\raggedright\arraybackslash}p{(\linewidth - 8\tabcolsep) * \real{0.2000}}
  >{\raggedright\arraybackslash}p{(\linewidth - 8\tabcolsep) * \real{0.2000}}
  >{\raggedright\arraybackslash}p{(\linewidth - 8\tabcolsep) * \real{0.2000}}
  >{\raggedright\arraybackslash}p{(\linewidth - 8\tabcolsep) * \real{0.2000}}
  >{\raggedright\arraybackslash}p{(\linewidth - 8\tabcolsep) * \real{0.2000}}@{}}
\toprule\noalign{}
\begin{minipage}[b]{\linewidth}\raggedright
Id
\end{minipage} & \begin{minipage}[b]{\linewidth}\raggedright
Status
\end{minipage} & \begin{minipage}[b]{\linewidth}\raggedright
Document
\end{minipage} & \begin{minipage}[b]{\linewidth}\raggedright
Locator
\end{minipage} & \begin{minipage}[b]{\linewidth}\raggedright
Method version
\end{minipage} \\
\midrule\noalign{}
\endhead
\bottomrule\noalign{}
\endlastfoot
\texttt{000-\allowbreak{}av-\allowbreak{}production} & executed & the
owner's method document, pinned by sha256 & owner-held, not
distributable & --- \\
\texttt{001-\allowbreak{}car},
\texttt{002-\allowbreak{}text-\allowbreak{}2027},
\texttt{003-\allowbreak{}hpc} & designed & none & none & --- \\
\texttt{004-\allowbreak{}adoption-\allowbreak{}guide} & executed &
\texttt{metodo-\allowbreak{}de-\allowbreak{}producao.\allowbreak{}md},
pinned by sha256 &
\texttt{instances/\allowbreak{}004-\allowbreak{}adoption-\allowbreak{}guide}
& 0.2.2 \\
\end{longtable}
}

An executed entry carries \texttt{document}, \texttt{sha256},
\texttt{locator} and \texttt{number\_\allowbreak{}format}
(\texttt{point-\allowbreak{}decimal} or
\texttt{comma-\allowbreak{}decimal}), none of them null; a designed
entry keeps the four null, so it can never pass for a pinned prior. The
entry of \texttt{000-\allowbreak{}av-\allowbreak{}production} also
declares a \texttt{heading\_\allowbreak{}style} (the shape of its
section, subsection and step headings), which the synthetic prior of the
sensitivity test follows (chapter 14).

Numbers about a prior do not enter the documents before they have a row
in the canon's ``Recorded facts'' table, which names the number, its
context (with the mark) and the sections of the prior's document where
it lives. The canon also commits, by sha256, to the text of each cited
section of the owner-held document, so that a later rewrite of a cited
section is detectable without publishing the text (chapter 15).

\section{Claims of efficacy and their
labels}\label{claims-of-efficacy-and-their-labels}

A unit that uses one of the canon's efficacy trigger phrases (``pays
off'', ``the method found'', ``the gates caught'', ``the review found'',
``err less'', ``the only way'', and the rest of the closed list) claims
that the method works, and has to carry one of five labels outside code
spans (\texttt{p-\allowbreak{}efficacy-\allowbreak{}mark}):

{\def\LTcaptype{none} % do not increment counter
\begin{longtable}[]{@{}
  >{\raggedright\arraybackslash}p{(\linewidth - 2\tabcolsep) * \real{0.5000}}
  >{\raggedright\arraybackslash}p{(\linewidth - 2\tabcolsep) * \real{0.5000}}@{}}
\toprule\noalign{}
\begin{minipage}[b]{\linewidth}\raggedright
Label
\end{minipage} & \begin{minipage}[b]{\linewidth}\raggedright
Where the evidence is registered
\end{minipage} \\
\midrule\noalign{}
\endhead
\bottomrule\noalign{}
\endlastfoot
\texttt{[prior:\allowbreak{}<id>]} & the method document of that
registered prior \\
\texttt{[experiment:\allowbreak{}<id>]} & the section of that id in
\texttt{60-\allowbreak{}experiments.\allowbreak{}md}, whose status is
\texttt{executed} and whose ``Provenance'' table gives the sha256 of at
least one input and one output \\
\texttt{[against:\allowbreak{}<id>]} & an executed experiment whose
result goes against the claim: counter-evidence is labeled like
evidence \\
\texttt{[lessons:\allowbreak{}<title>]} & the entry of
\texttt{99-\allowbreak{}lessons.\allowbreak{}md} whose heading is that
title \\
\texttt{[hypothesis]} & nowhere yet: an assumption, to be measured \\
\end{longtable}
}

The gate is lexical and says so: a paraphrase of a claim escapes it, and
no list closes that gap; the mitigation is a review with a frozen
rubric. A claim labeled \texttt{[hypothesis]} that states a number or a
universal quantifier gets a low note: the label stays valid, the claim
has the shape of a measurement with none behind it.

\section{The experiments}\label{the-experiments}

\texttt{02-\allowbreak{}reflexive-\allowbreak{}application/\allowbreak{}60-\allowbreak{}experiments.\allowbreak{}md}
registers the measurements about the method itself, each with its
question, object, conditions, executors by family, protocol, predictions
written before measuring, results, what it shows and does not show,
cost, what became a gate or a lesson, and a provenance table. Only an
executed experiment backs a claim.

{\def\LTcaptype{none} % do not increment counter
\begin{longtable}[]{@{}
  >{\raggedright\arraybackslash}p{(\linewidth - 6\tabcolsep) * \real{0.2500}}
  >{\raggedright\arraybackslash}p{(\linewidth - 6\tabcolsep) * \real{0.2500}}
  >{\raggedright\arraybackslash}p{(\linewidth - 6\tabcolsep) * \real{0.2500}}
  >{\raggedright\arraybackslash}p{(\linewidth - 6\tabcolsep) * \real{0.2500}}@{}}
\toprule\noalign{}
\begin{minipage}[b]{\linewidth}\raggedright
Id
\end{minipage} & \begin{minipage}[b]{\linewidth}\raggedright
Status
\end{minipage} & \begin{minipage}[b]{\linewidth}\raggedright
Question
\end{minipage} & \begin{minipage}[b]{\linewidth}\raggedright
Result in one line
\end{minipage} \\
\midrule\noalign{}
\endhead
\bottomrule\noalign{}
\endlastfoot
E1 & executed & does the reviewer with the frozen rubric name two
injected defects, and does a reviewer of another family differ? & both
executors named 2 of 2 and gave \texttt{rejected} on the perturbed copy;
one control run ended for length and is a tool failure, not a verdict \\
E2 & designed & the counterfactual: the same target produced with and
without the method & not run; two productions per target \\
E3 & executed & the same reviewers without the rubric and the reviewer
card & both still named 2 of 2, but classed the invented fact
\texttt{high} and accepted with conditions: the apparatus turned the
finding into a verdict \\
E4 & designed & seeded escape: does the apparatus let fewer of ten
seeded defects escape than a careful baseline, and does another family
find more? & protocol and predictions written; awaits the owner's
thresholds and budget \\
E5 & designed & can a third party set up a cell for a target with domain
measurements and reach a verdict without improvising instrument? &
awaits the owner finding the third party \\
\end{longtable}
}

\section{What would promote or withdraw a
claim}\label{what-would-promote-or-withdraw-a-claim}

Written before the data, as the gates are
(\texttt{01-\allowbreak{}principle.\allowbreak{}md}, section 9), so that
the material leaves the state of a working hypothesis by a measurement
and not by a decision. The thresholds are the owner's to accept or
adjust before E4 runs:

{\def\LTcaptype{none} % do not increment counter
\begin{longtable}[]{@{}
  >{\raggedright\arraybackslash}p{(\linewidth - 4\tabcolsep) * \real{0.3333}}
  >{\raggedright\arraybackslash}p{(\linewidth - 4\tabcolsep) * \real{0.3333}}
  >{\raggedright\arraybackslash}p{(\linewidth - 4\tabcolsep) * \real{0.3333}}@{}}
\toprule\noalign{}
\begin{minipage}[b]{\linewidth}\raggedright
Claim
\end{minipage} & \begin{minipage}[b]{\linewidth}\raggedright
Promoted when
\end{minipage} & \begin{minipage}[b]{\linewidth}\raggedright
Withdrawn when
\end{minipage} \\
\midrule\noalign{}
\endhead
\bottomrule\noalign{}
\endlastfoot
a reviewer of another family adds to one of the producer's family & in
E4 or an equivalent design, the arm of the other family finds at least 1
more defect per 10 seeded, in at least 2 of 3 repetitions & the
difference is nil in 3 repetitions \\
the review apparatus adds to a careful review without it & in E4, the
escape rate of arm A is lower than that of arm B by at least 20 points &
the difference is under 5 points \\
the method is portable outside the documentary domain & one execution
outside it, by someone who is not the owner, with its domain gates read
by the validator and a verdict of the canon & that execution cannot
reach a verdict without improvising instrument \\
\end{longtable}
}

\part{The material and the instrument}

\chapter{The reflexive application}\label{the-reflexive-application}

\section{What it is, and what round counts as
evidence}\label{what-it-is-and-what-round-counts-as-evidence}

\texttt{02-\allowbreak{}reflexive-\allowbreak{}application/\allowbreak{}}
is the method applied to itself: a production cell whose target is the
set of documents that shows how to assemble a production cell. Its brief
(\texttt{00-\allowbreak{}brief.\allowbreak{}md}) declares the target,
the two readers, five non-negotiables with ids
(\texttt{nn-\allowbreak{}method},
\texttt{nn-\allowbreak{}prior-\allowbreak{}facts},
\texttt{nn-\allowbreak{}hypothesis}, \texttt{nn-\allowbreak{}form},
\texttt{nn-\allowbreak{}anonymity}), the owner decisions, a budget (two
returns per document per failed gate, then escalation) and the stopping
criterion: \textbf{exactly one independent review round}.

That first round, and only that round, was a production under the
method. The review gave \texttt{rejected} for the version it read (one
blocking, two high); the blocking and high notes were fixed and checked
by the gates, without a second review, as the brief prescribed; the 22
remaining notes are kept, frozen, in
\texttt{90-\allowbreak{}review/\allowbreak{}notes.\allowbreak{}json}.
Every round after it (v0.2 on) is a \textbf{maintenance round} of the
material: recorded in the round record with
\texttt{kind:\allowbreak{} maintenance}, its scope and the reviewer's
relation to the producer, emitting no verdict of the canon, with its
decisions in \texttt{CHANGELOG.\allowbreak{}md}. The reflexive
application is evidence about the method only for round 1; the
maintenance rounds are evidence about the material.

\section{The folder}\label{the-folder}

{\def\LTcaptype{none} % do not increment counter
\begin{longtable}[]{@{}
  >{\raggedright\arraybackslash}p{(\linewidth - 4\tabcolsep) * \real{0.3333}}
  >{\raggedright\arraybackslash}p{(\linewidth - 4\tabcolsep) * \real{0.3333}}
  >{\raggedright\arraybackslash}p{(\linewidth - 4\tabcolsep) * \real{0.3333}}@{}}
\toprule\noalign{}
\begin{minipage}[b]{\linewidth}\raggedright
File or folder
\end{minipage} & \begin{minipage}[b]{\linewidth}\raggedright
Artifact
\end{minipage} & \begin{minipage}[b]{\linewidth}\raggedright
What it is
\end{minipage} \\
\midrule\noalign{}
\endhead
\bottomrule\noalign{}
\endlastfoot
\texttt{00-\allowbreak{}brief.\allowbreak{}md} & \texttt{brief} &
target, non-negotiables with ids, owner decisions, budget and stop, the
protocol of the maintenance rounds \\
\texttt{01-\allowbreak{}canon.\allowbreak{}md} & \texttt{canon} & method
terms, terms to avoid, the fixed lists the validator reads, form rules,
recorded facts, section commitments, historical vocabulary \\
\texttt{02-\allowbreak{}artifacts.\allowbreak{}md} & \texttt{artifacts}
& the artifact table (the one source of the gate map) and the graph \\
\texttt{03-\allowbreak{}roles/\allowbreak{}} & \texttt{cards} & the four
role cards: writer, verifier, runner, reviewer \\
\texttt{04-\allowbreak{}gates-\allowbreak{}and-\allowbreak{}rubric.\allowbreak{}md}
& \texttt{gates-\allowbreak{}and-\allowbreak{}rubric} & the gate table
with magnitude, limit and justification; the measurement
\texttt{m-\allowbreak{}determinism}; the size ranges; the frozen
``Rubric'' section \\
\texttt{05-\allowbreak{}contracts/\allowbreak{}} & \texttt{contracts} &
six JSON Schemas (appendix B) \\
\texttt{06-\allowbreak{}gate-\allowbreak{}semantics.\allowbreak{}md} &
\texttt{gate-\allowbreak{}semantics} & what each gate reads in detail;
the hash policy; domain gates; the false greens table; the gate
lifecycle; the provenance of the rubric \\
\texttt{10-\allowbreak{}generic-\allowbreak{}recipe.\allowbreak{}md} &
\texttt{recipe} & the fifteen steps, with ``Reads'', ``Creates'',
``Does'' and ``Done when'' each \\
\texttt{20-\allowbreak{}instance-\allowbreak{}car.\allowbreak{}md},
\texttt{21-\allowbreak{}instance-\allowbreak{}text-\allowbreak{}2027.\allowbreak{}md},
\texttt{22-\allowbreak{}instance-\allowbreak{}hpc.\allowbreak{}md} &
\texttt{instance-\allowbreak{}*} & the three designed instances \\
\texttt{60-\allowbreak{}experiments.\allowbreak{}md} &
\texttt{experiments} & E1 to E5 \\
\texttt{80-\allowbreak{}validator/\allowbreak{}} &
\texttt{instrument-\allowbreak{}validator},
\texttt{validator-\allowbreak{}tests},
\texttt{validation-\allowbreak{}report} & the scripts, the test suite,
\texttt{report.\allowbreak{}md} and
\texttt{sensitivity.\allowbreak{}md} \\
\texttt{90-\allowbreak{}review/\allowbreak{}} & \texttt{review},
\texttt{replies}, \texttt{open-\allowbreak{}notes} & the frozen records
of round 1: \texttt{rubric.\allowbreak{}sha256},
\texttt{auditor-\allowbreak{}report.\allowbreak{}md},
\texttt{auditor-\allowbreak{}notes.\allowbreak{}json},
\texttt{notes.\allowbreak{}json}, \texttt{replies.\allowbreak{}md} \\
\texttt{99-\allowbreak{}lessons.\allowbreak{}md} & \texttt{lessons} &
one entry per round; append-only \\
\texttt{state/\allowbreak{}round-\allowbreak{}record.\allowbreak{}jsonl}
& \texttt{round-\allowbreak{}record} & one line per round \\
\texttt{state/\allowbreak{}notes-\allowbreak{}status.\allowbreak{}jsonl}
& \texttt{notes-\allowbreak{}status} & the state of each tracked note
and panel finding \\
\texttt{state/\allowbreak{}owner-\allowbreak{}decisions.\allowbreak{}jsonl}
& \texttt{owner-\allowbreak{}decisions} & the decisions only the owner
can take \\
\texttt{README.\allowbreak{}md} & \texttt{readme} & the index and the
runbook \\
\end{longtable}
}

Four artifacts live at the root and are declared in the same table:
\texttt{version} (\texttt{VERSION}), \texttt{license},
\texttt{changelog}, \texttt{registry}
(\texttt{instances/\allowbreak{}registry.\allowbreak{}json}),
\texttt{readme-\allowbreak{}root} and \texttt{skill} (the three files of
the skill). The gate \texttt{p-\allowbreak{}declared} checks that every
Markdown document of the folder belongs to an artifact or is listed in
the canon as outside the table; this production lists none.

\section{The artifact graph}\label{the-artifact-graph}

Each arrow is ``is read by''. The graph is acyclic and the validator
checks it on every pass (\texttt{p-\allowbreak{}references}).

\begin{figure}
\centering
\pandocbounded{\includegraphics[keepaspectratio,alt={The artifact graph of the reflexive application; an arrow is ``is read by''.}]{diagrams/artifact-graph.pdf}}
\caption{The artifact graph of the reflexive application; an arrow is
``is read by''.}\label{fig-artifact-graph}
\end{figure}

Two things the graph says about the order of the work: the three
instances only start after the recipe passes its gates, because it is
the recipe that fixes what an instance has to fill in; and the review
only starts after the validator passes on every document, so that the
one review round is not spent on defects the code already finds.

\section{The canon's tables the validator
reads}\label{the-canons-tables-the-validator-reads}

The validator reads the tables of
\texttt{01-\allowbreak{}canon.\allowbreak{}md} directly: changing a
table changes what it checks, and the canon's ``Canon schema'' section
is versioned (\texttt{v1}) so that a validator refuses a canon of
another version.

{\def\LTcaptype{none} % do not increment counter
\begin{longtable}[]{@{}
  >{\raggedright\arraybackslash}p{(\linewidth - 2\tabcolsep) * \real{0.5000}}
  >{\raggedright\arraybackslash}p{(\linewidth - 2\tabcolsep) * \real{0.5000}}@{}}
\toprule\noalign{}
\begin{minipage}[b]{\linewidth}\raggedright
Section of the canon
\end{minipage} & \begin{minipage}[b]{\linewidth}\raggedright
Read by
\end{minipage} \\
\midrule\noalign{}
\endhead
\bottomrule\noalign{}
\endlastfoot
Method terms & the glossary; not a gate \\
Terms to avoid & \texttt{p-\allowbreak{}terms} \\
Instantiation table rows, Mandatory sections of an instance &
\texttt{p-\allowbreak{}table}, \texttt{p-\allowbreak{}sections} \\
Role card fields & \texttt{p-\allowbreak{}card} \\
Origin labels of a limit & \texttt{p-\allowbreak{}limit} \\
Method components & \texttt{p-\allowbreak{}components} \\
Prior-statement triggers, Efficacy triggers and their labels &
\texttt{p-\allowbreak{}prior-\allowbreak{}mark},
\texttt{p-\allowbreak{}efficacy-\allowbreak{}mark} \\
Efficacy mark in append-only documents & the heading from which
\texttt{99-\allowbreak{}lessons.\allowbreak{}md} is measured \\
Documents outside the artifact table &
\texttt{p-\allowbreak{}declared} \\
Chained ledgers & \texttt{p-\allowbreak{}contract-\allowbreak{}shape}
(the lines before each chain) \\
Frozen records & \texttt{p-\allowbreak{}frozen-\allowbreak{}rubric} \\
Accepted paths, Month names, Machine roots &
\texttt{p-\allowbreak{}anonymity} \\
Recorded facts, Section commitments & \texttt{p-\allowbreak{}prior} \\
\end{longtable}
}

A cell's JSON canon carries the same lists under the keys the scaffold
writes (\texttt{terms to avoid}, \texttt{instantiation table rows},
\texttt{mandatory sections}, \texttt{role card fields},
\texttt{origin labels}, \texttt{method components},
\texttt{prior-\allowbreak{}statement triggers}, \texttt{recorded facts},
\texttt{efficacy triggers},
\texttt{documents outside the artifact list}), each a list of plain
strings except the recorded facts, which are objects.

\chapter{The validator}\label{the-validator}

\section{Purpose}\label{purpose}

\texttt{validator.\allowbreak{}py} is the \textbf{runner}'s instrument:
it reads every document the artifact list assigns a gate to, measures
each gate on it, and writes a deterministic report with the magnitude,
the limit and the result of each (document, gate) pair, plus the notes
the failures produce. It also runs the global gates (the note ledgers,
the frozen records, the sensitivity report, the declared files, the
domain results), carries an embedded self-test, freezes the rubric, and
measures its own determinism.

\begin{figure}
\centering
\pandocbounded{\includegraphics[keepaspectratio,alt={One run of the validator: inputs, the checker, the report and the exit code.}]{diagrams/validator-run.pdf}}
\caption{One run of the validator: inputs, the checker, the report and
the exit code.}\label{fig-validator-run}
\end{figure}

\section{Inputs and the layout}\label{inputs-and-the-layout}

The validator is written for the corpus's Markdown layout and remapped
to a cell's JSON layout by
\texttt{-\allowbreak{}-\allowbreak{}layout layout.\allowbreak{}json}.
The layout has a closed set of keys; a key the validator does not read
is refused (exit 2) instead of silently falling back to a default it
meant to override.

{\def\LTcaptype{none} % do not increment counter
\begin{longtable}[]{@{}
  >{\raggedright\arraybackslash}p{(\linewidth - 6\tabcolsep) * \real{0.2500}}
  >{\raggedright\arraybackslash}p{(\linewidth - 6\tabcolsep) * \real{0.2500}}
  >{\raggedright\arraybackslash}p{(\linewidth - 6\tabcolsep) * \real{0.2500}}
  >{\raggedright\arraybackslash}p{(\linewidth - 6\tabcolsep) * \real{0.2500}}@{}}
\toprule\noalign{}
\begin{minipage}[b]{\linewidth}\raggedright
Layout key
\end{minipage} & \begin{minipage}[b]{\linewidth}\raggedright
Meaning
\end{minipage} & \begin{minipage}[b]{\linewidth}\raggedright
Default (the corpus)
\end{minipage} & \begin{minipage}[b]{\linewidth}\raggedright
A scaffolded cell
\end{minipage} \\
\midrule\noalign{}
\endhead
\bottomrule\noalign{}
\endlastfoot
\texttt{version} & the layout format & --- & \texttt{1} \\
\texttt{artifacts.\allowbreak{}canon} & the canon &
\texttt{01-\allowbreak{}canon.\allowbreak{}md} &
\texttt{02-\allowbreak{}canon.\allowbreak{}json} \\
\texttt{artifacts.\allowbreak{}artifacts} & the artifact list &
\texttt{02-\allowbreak{}artifacts.\allowbreak{}md} &
\texttt{03-\allowbreak{}artifacts.\allowbreak{}json} \\
\texttt{artifacts.\allowbreak{}gates-\allowbreak{}and-\allowbreak{}rubric}
& the gates file &
\texttt{04-\allowbreak{}gates-\allowbreak{}and-\allowbreak{}rubric.\allowbreak{}md}
& \texttt{06-\allowbreak{}gates.\allowbreak{}md} \\
\texttt{artifacts.\allowbreak{}roles-\allowbreak{}dir} & the role cards
& \texttt{03-\allowbreak{}roles} & \texttt{04-\allowbreak{}roles} \\
\texttt{artifacts.\allowbreak{}rubric} & a dedicated rubric file (else
the ``Rubric'' section of the gates file) & absent &
\texttt{07-\allowbreak{}rubric.\allowbreak{}md} \\
\texttt{artifacts.\allowbreak{}sensitivity} & the sensitivity report &
\texttt{80-\allowbreak{}validator/\allowbreak{}sensitivity.\allowbreak{}md}
&
\texttt{08-\allowbreak{}validator/\allowbreak{}sensitivity.\allowbreak{}md} \\
\texttt{artifacts.\allowbreak{}perturbations} & the perturbation data
file of a cell & absent (the corpus keeps
\texttt{perturbation\_\allowbreak{}set.\allowbreak{}py}) &
\texttt{08-\allowbreak{}validator/\allowbreak{}perturbations.\allowbreak{}json} \\
\texttt{artifacts.\allowbreak{}report} & the validator report &
\texttt{80-\allowbreak{}validator/\allowbreak{}report.\allowbreak{}md} &
\texttt{08-\allowbreak{}validator/\allowbreak{}report.\allowbreak{}md} \\
\texttt{artifacts.\allowbreak{}experiments},
\texttt{gate-\allowbreak{}semantics},
\texttt{notes-\allowbreak{}status}, \texttt{round-\allowbreak{}record},
\texttt{lessons}, \texttt{brief}, \texttt{owner-\allowbreak{}decisions}
& the other documents the global gates read & the corpus's files & as
the scaffold writes, or the defaults \\
\texttt{gates} & the parser of each gate:
\texttt{builtin-\allowbreak{}markdown} or
\texttt{builtin-\allowbreak{}json-\allowbreak{}canon}; every gate of the
map has to use the same one & all \texttt{builtin-\allowbreak{}markdown}
& all \texttt{builtin-\allowbreak{}json-\allowbreak{}canon} \\
\texttt{registry} & the instance registry, relative to the folder &
\texttt{.\allowbreak{}.\allowbreak{}/\allowbreak{}instances/\allowbreak{}registry.\allowbreak{}json}
& \texttt{priors.\allowbreak{}json} \\
\texttt{domain-\allowbreak{}results} & the result files of the
\texttt{d-\allowbreak{}*} gates & \texttt{[]} & \texttt{[]} \\
\texttt{version-\allowbreak{}label} & what \texttt{\{version\}} expands
to in paths and commands & absent & --- \\
\end{longtable}
}

The other inputs are command-line options (appendix C):
\texttt{-\allowbreak{}-\allowbreak{}case} for each prior method document
(repeatable; each is identified by its sha256 against the registry, and
the number check of a \texttt{[prior:\allowbreak{}<id>]} unit reads the
document identified as \texttt{<id>}, in the number format its registry
entry declares); \texttt{-\allowbreak{}-\allowbreak{}notes} for the note
ledger; \texttt{-\allowbreak{}-\allowbreak{}rubric-\allowbreak{}hash}
for the frozen pin; \texttt{-\allowbreak{}-\allowbreak{}forbidden} for
the list of forbidden identifiers;
\texttt{-\allowbreak{}-\allowbreak{}allow-\allowbreak{}pending} for the
gates whose pending rows the caller accepts by name.

\section{Architecture}\label{architecture}

Since v0.6.0 the checker is one class per family of gates over a core
that holds the state of a run and the interface every family uses. The
split moved the methods unchanged, and the report of the corpus was
byte-identical before and after it.

\begin{figure}
\centering
\pandocbounded{\includegraphics[keepaspectratio,alt={The classes of the validator: one family per kind of gate over a core.}]{diagrams/gate-families.pdf}}
\caption{The classes of the validator: one family per kind of gate over
a core.}\label{fig-gate-families}
\end{figure}

A run proceeds as follows (\texttt{CheckerCore.\allowbreak{}run}):

\begin{enumerate}
\def\labelenumi{\arabic{enumi}.}
\tightlist
\item
  The controls are loaded through the layout: the canon, the artifact
  list, the gates file (its merged gate tables, one row per gate, else
  exit 2; its size ranges), the registry. Every gate id of the table or
  of the map has to be a builtin, a global gate or a
  \texttt{d-\allowbreak{}*} gate, else exit 2.
\item
  The gate map is built from the artifact list: for each gate, the
  artifacts that list it; for each artifact, the files its paths resolve
  to. Domain gates are skipped here and read from their result files in
  the global pass.
\item
  The documents that carry \texttt{p-\allowbreak{}prior} are scanned for
  \texttt{[prior:\allowbreak{}<id>]} ids, so that
  \texttt{p-\allowbreak{}prior-\allowbreak{}pin} knows which executed
  priors it has to pin.
\item
  A run that measures no documents writes a blocking note on
  \texttt{artifacts}: a production with no documents is not green, so a
  freshly scaffolded cell fails until it has real artifacts.
\item
  For each document in sorted order, for each (gate, artifact) pair in
  sorted order, the handler of the gate records a measurement and its
  notes. A hit is a note of the gate's severity (high unless the gate
  says otherwise) and fails the row; an advisory is a low note that does
  not. Notes are capped at 50 per gate, with one note saying how many
  were omitted.
\item
  The global gates run once (\texttt{check\_\allowbreak{}global}): the
  brief's languages, the frozen records, the declared files and
  documents, the sensitivity report, the note ledgers, the domain
  results.
\item
  The Limit column of the gate table is compared with the limits the
  validator applied
  (\texttt{check\_\allowbreak{}limit\_\allowbreak{}column}).
\end{enumerate}

\section{The report}\label{the-report}

\texttt{report.\allowbreak{}md} has three sections:
\textbf{Measurements} (one row per document and gate, sorted, with
magnitude, limit and \texttt{pass}, \texttt{fail} or \texttt{pending},
then the lines ``Pending accepted by the caller'', ``Pending not
accepted'' and, when an accepted gate had nothing pending, ``Accepted
but not pending''); \textbf{Instrument measurement} (the
\texttt{m-\allowbreak{}determinism} row, the count of the embedded
self-test, and the sha256 of the test suite next to the validator); and
\textbf{Notes} (the count by severity and a table with target, severity,
text and evidence).

The report is deterministic by construction: no date, relative paths,
rows sorted by document and gate, notes sorted and deduplicated. With
\texttt{-\allowbreak{}-\allowbreak{}output FILE} the report goes to the
file and one summary line goes to the error stream
(\texttt{validator:\allowbreak{} exit N; notes:\allowbreak{} .\allowbreak{}.\allowbreak{}.\allowbreak{}; pending not accepted:\allowbreak{} .\allowbreak{}.\allowbreak{}.\allowbreak{}}),
so that a failing run with \texttt{-\allowbreak{}-\allowbreak{}output}
is never silent.

\texttt{-\allowbreak{}-\allowbreak{}determinism} runs the validator
twice more as child processes, with \texttt{PYTHONHASHSEED} set to 1 and
2 and, on the second run, \texttt{LC\_\allowbreak{}ALL=C}, and compares
the sha256 of the two reports; a difference is a high note. The working
folder and the Python version (3.9 and 3.12) were varied by hand, not by
the measurement, and gave the same report.

\section{Exit codes}\label{exit-codes}

\begin{figure}
\centering
\pandocbounded{\includegraphics[keepaspectratio,alt={How the validator decides its exit code.}]{diagrams/exit-codes.pdf}}
\caption{How the validator decides its exit code.}\label{fig-exit-codes}
\end{figure}

{\def\LTcaptype{none} % do not increment counter
\begin{longtable}[]{@{}
  >{\raggedright\arraybackslash}p{(\linewidth - 2\tabcolsep) * \real{0.5000}}
  >{\raggedright\arraybackslash}p{(\linewidth - 2\tabcolsep) * \real{0.5000}}@{}}
\toprule\noalign{}
\begin{minipage}[b]{\linewidth}\raggedright
Code
\end{minipage} & \begin{minipage}[b]{\linewidth}\raggedright
Meaning
\end{minipage} \\
\midrule\noalign{}
\endhead
\bottomrule\noalign{}
\endlastfoot
0 & every gate was measured, or its pending was accepted by name, and no
note is high or blocking \\
1 & at least one high or blocking note (a
\texttt{-\allowbreak{}-\allowbreak{}determinism} mismatch included) \\
2 & invalid input or tool failure: a layout key the validator does not
read, a file that is not UTF-8, a truncated table row, a gate it does
not implement, a gate with two rows, a gate named by
\texttt{-\allowbreak{}-\allowbreak{}allow-\allowbreak{}pending} outside
its closed list, a domain result that breaks its contract, a Python
older than 3.9, an unexpected exception \\
4 & no such note, but a gate not measured that
\texttt{-\allowbreak{}-\allowbreak{}allow-\allowbreak{}pending} does not
accept \\
\end{longtable}
}

\texttt{-\allowbreak{}-\allowbreak{}allow-\allowbreak{}pending} accepts
only four names: \texttt{p-\allowbreak{}anonymity},
\texttt{p-\allowbreak{}prior} and
\texttt{p-\allowbreak{}prior-\allowbreak{}pin} (the halves that need an
input the material does not distribute: the forbidden list, the
owner-held prior document) and
\texttt{p-\allowbreak{}sensitivity-\allowbreak{}synthetic} (a
sensitivity report a clone regenerated with a synthetic prior). Any
other name is exit 2, \texttt{p-\allowbreak{}sensitivity} included; a
missing sensitivity report is blocking; and
\texttt{p-\allowbreak{}anonymity-\allowbreak{}language} (a language the
brief declares without month names) cannot be accepted either. The
history behind the closed list: the panel that reviewed v0.5.0 deleted
the sensitivity report, accepted the gate by name and got exit 0.

\section{The hash policy}\label{the-hash-policy}

Every sha256 the instrument computes --- the documents passed with
\texttt{-\allowbreak{}-\allowbreak{}case}, the rubric, the gate section,
the canon, the resolved gate map, the validator itself, the inputs of a
domain result and the lines of a chained ledger --- is the sha256 of the
text as the validator reads it: UTF-8, a byte-order mark dropped, every
line break read as LF. A file that is not UTF-8 text (a domain input
only) is hashed as its bytes. The scripts write UTF-8 to the console
whatever its code page.

\section{The self-test and the test
suite}\label{the-self-test-and-the-test-suite}

\texttt{validator.\allowbreak{}py -\allowbreak{}-\allowbreak{}self-\allowbreak{}test}
runs 185 checks (at 0.6.0) embedded in the script, with no inputs; the
runbook and every cell use it without pytest, and the report prints its
count. The pytest suite
(\texttt{80-\allowbreak{}validator/\allowbreak{}test\_\allowbreak{}validator.\allowbreak{}py},
280 tests together with the scaffold's
\texttt{test\_\allowbreak{}scaffold.\allowbreak{}py} at 0.6.0) runs the
self-test and one case per rule: one that passes and one that is
rejected. Both are run with
\texttt{python3 -\allowbreak{}B -\allowbreak{}m pytest -\allowbreak{}p no:\allowbreak{}cacheprovider},
because pytest imports the files before their own code can switch the
bytecode cache off.

\section{The scaffold}\label{the-scaffold}

\texttt{skills/\allowbreak{}cooperatio-\allowbreak{}teamwork/\allowbreak{}scripts/\allowbreak{}scaffold.\allowbreak{}py}
creates an empty production cell for a target, refusing to overwrite
anything: it builds the cell in a staging folder and moves it into
place, so a failure leaves no half-written cell. It copies the
validator, the sensitivity script, the perturbation module and five
contracts from the material, so it has to run from inside the material
(or with the skill folder kept inside it); copied alone, the skill
creates the cell without them, says so, and the cell has no validator
until one is copied in. Chapter 17 lists what it writes.

\chapter{The gate catalogue}\label{the-gate-catalogue}

This chapter is the reference of the gates the validator implements at
version 0.6.0: 17 builtin gates measured per document, 4 global gates
measured once per run, one measurement of the instrument, and the family
of domain gates a cell supplies. For each gate: the artifacts the
reflexive application assigns it to, what it reads, its limit, when it
is pending, and the semantics worth knowing before reading a note. The
authoritative text is the gate table of
\texttt{04-\allowbreak{}gates-\allowbreak{}and-\allowbreak{}rubric.\allowbreak{}md}
and the declared semantics of
\texttt{06-\allowbreak{}gate-\allowbreak{}semantics.\allowbreak{}md}.

\section{Conventions shared by every
gate}\label{conventions-shared-by-every-gate}

\begin{itemize}
\tightlist
\item
  \textbf{Units.} A gate reads a document in units: a paragraph, a list
  item, a table row or a heading. A mark or a label applies to the unit
  it sits in.
\item
  \textbf{Code spans and fences.} Since v0.6.0 one tokenizer reads
  fenced blocks for every gate, in the CommonMark rules (three or more
  backticks or tildes, indented up to three spaces, closed by a line of
  the same character at least as long; an unclosed fence runs to the end
  and is a note). Marks, labels and triggers are read outside code
  spans: a mark between backticks is notation, not a mark.
\item
  \textbf{Placeholders.} An angle-bracket word (\texttt{<fill in>}),
  \texttt{TBD} or \texttt{TODO} outside code spans is an unfilled field
  in exactly four places: a card (\texttt{p-\allowbreak{}card}), a gate
  row (\texttt{p-\allowbreak{}limit}), an instantiation cell
  (\texttt{p-\allowbreak{}table}) and an artifact record
  (\texttt{p-\allowbreak{}references}). A placeholder in prose passes
  every gate and is the review's to find.
\item
  \textbf{Empty lists.} A gate whose object is a list of the canon is
  pending when that list is empty, never a pass of 0 of 0.
\item
  \textbf{Severity.} A hit is a high note unless the gate's semantics
  say blocking (a cycle, an invented prior number, a malformed mark, a
  missing declared file, a broken chain, a changed frozen record, a
  stale domain result) or low (an advisory).
\item
  \textbf{Cap.} 50 notes per gate; the 51st is replaced by a count.
\end{itemize}

\section{Gates measured per document}\label{gates-measured-per-document}

\subsection{\texorpdfstring{\texttt{p-\allowbreak{}terms} --- terms to
avoid}{ --- terms to avoid}}\label{p-terms-terms-to-avoid}

\begin{itemize}
\tightlist
\item
  \textbf{Assigned to:} every prose document of the material, the
  changelog, the replies, the experiments and the skill's files.
\item
  \textbf{Reads:} occurrences of each term of the canon's ``Terms to
  avoid'' table (\texttt{quality gate}, \texttt{checkpoint},
  \texttt{pipeline}, \texttt{feedback}, \texttt{iteration},
  \texttt{single source of truth}, \texttt{bible}, \texttt{ruler},
  \texttt{review board}), as whole words, case-insensitive, with a
  plural suffix, outside code and link targets; in the canon, outside
  the table itself.
\item
  \textbf{Limit:} 0. \texttt{[method]}: a synonym outside the canon is
  the first sign of drift, and replacing it costs one word.
\item
  \textbf{Exempt:} the frozen ledger under
  \texttt{90-\allowbreak{}review/\allowbreak{}}, by design: vocabulary
  rules govern living documents.
\end{itemize}

\subsection{\texorpdfstring{\texttt{p-\allowbreak{}canon-\allowbreak{}link}
--- a link to the
canon}{ --- a link to the canon}}\label{p-canon-link-a-link-to-the-canon}

\begin{itemize}
\tightlist
\item
  \textbf{Assigned to:} every document except the canon and the frozen
  records.
\item
  \textbf{Reads:} links naming
  \texttt{01-\allowbreak{}canon.\allowbreak{}md} (or the cell's canon
  file) that resolve, relative to the document's own folder, to the
  canon of the layout; and links naming it that do not resolve.
\item
  \textbf{Limit:} ≥ 1 resolved, 0 unresolved. \texttt{[method]}: a
  reader who lands on a loose document needs to know where the
  definitions are.
\end{itemize}

\subsection{\texorpdfstring{\texttt{p-\allowbreak{}anonymity} --- paths,
dates,
identifiers}{ --- paths, dates, identifiers}}\label{p-anonymity-paths-dates-identifiers}

\begin{itemize}
\tightlist
\item
  \textbf{Assigned to:} every document, the registry, the three ledgers,
  the principle, the root README and the skill.
\item
  \textbf{Reads, the path-and-date half:} absolute paths (a segment
  under a common root such as \texttt{/\allowbreak{}tmp},
  \texttt{/\allowbreak{}home}, \texttt{/\allowbreak{}opt}; a path of two
  or more segments; file URLs; home folders by user name; UNC paths);
  dates in year-month-day, day/month/year, year/month/day and
  day.month.year formats; dates written with a month's name in English,
  Portuguese, Spanish, French, German or Italian, and in a language the
  canon's ``Month names'' table adds; the expressions ``made on'',
  ``written on'', ``generated on'' followed by a date; e-mail addresses;
  host names under a private suffix (\texttt{.\allowbreak{}local},
  \texttt{.\allowbreak{}lan}, \texttt{.\allowbreak{}internal},
  \texttt{.\allowbreak{}intranet}, \texttt{.\allowbreak{}corp}); a root
  of one segment the canon's ``Machine roots'' table names. A month and
  a year with no day, a version number, a relative path and a public
  domain name are not hits.
\item
  \textbf{Reads, the identifier half:} the entries of the
  \texttt{-\allowbreak{}-\allowbreak{}forbidden} list, by name or as
  \texttt{sha256:\allowbreak{}<hex>} entries compared with the sha256 of
  each lowercased word. Model identifiers belong to that list. The list
  is kept outside the material, since it contains what it forbids; the
  corpus runs with none, so this half is \textbf{pending by design} and
  accepted by name.
\item
  \textbf{Limit:} 0. \texttt{[brief]}: the brief forbids absolute paths,
  authoring dates, machine names and user names; in the method the gate
  is an option of the instance.
\item
  \textbf{Valves:} an interpreter line (\texttt{\#!}) on line 1 is not
  read; a path hit equal to an entry of the canon's ``Accepted paths''
  table for that document is dropped (the entry has to be a path and
  nothing else; the canon cannot accept a path for itself); a final
  period is not part of the path compared.
\item
  \textbf{Languages:} the brief declares the languages of its documents
  on a line \texttt{Languages:\allowbreak{} en, pt}; a declared language
  with no built-in months and no row of the canon's table is pending on
  the brief under
  \texttt{p-\allowbreak{}anonymity-\allowbreak{}language}, which no
  caller accepts.
\end{itemize}

\subsection{\texorpdfstring{\texttt{p-\allowbreak{}form} --- plain
Markdown}{ --- plain Markdown}}\label{p-form-plain-markdown}

\begin{itemize}
\tightlist
\item
  \textbf{Assigned to:} every prose document and the changelog.
\item
  \textbf{Reads:} HTML tags and comments outside code; fenced blocks
  (backtick or tilde, indented up to three spaces) whose info string is
  not exactly \texttt{text}.
\item
  \textbf{Limit:} 0. \texttt{[brief]}: plain Markdown and pseudocode
  only in \texttt{text} blocks are non-negotiables of the brief.
\end{itemize}

\subsection{\texorpdfstring{\texttt{p-\allowbreak{}size} --- words in
range}{ --- words in range}}\label{p-size-words-in-range}

\begin{itemize}
\tightlist
\item
  \textbf{Assigned to:} the documents of the ``Size ranges'' table of
  the gates file.
\item
  \textbf{Reads:} words separated by whitespace; since v0.6.0 the pipes
  and separators of a table are not words.
\item
  \textbf{Limit:} within the range declared, before writing, in the
  ``Size ranges'' table. \texttt{[hypothesis]}: a range is raised before
  writing, never met by cutting text.
\end{itemize}

\subsection{\texorpdfstring{\texttt{p-\allowbreak{}table} --- the
instantiation
table}{ --- the instantiation table}}\label{p-table-the-instantiation-table}

\begin{itemize}
\tightlist
\item
  \textbf{Assigned to:} the three instances.
\item
  \textbf{Reads:} the rows of the canon's ``Instantiation table rows''
  list present in the two-column table that opens the instance; empty
  cells or cells holding a placeholder.
\item
  \textbf{Limit:} 14 of 14, 0 empty. \texttt{[method]}: the table is the
  summary that makes instances comparable.
\end{itemize}

\subsection{\texorpdfstring{\texttt{p-\allowbreak{}sections} --- the
mandatory
sections}{ --- the mandatory sections}}\label{p-sections-the-mandatory-sections}

\begin{itemize}
\tightlist
\item
  \textbf{Assigned to:} the three instances.
\item
  \textbf{Reads:} the headings of the canon's ``Mandatory sections of an
  instance'' list.
\item
  \textbf{Limit:} 14 of 14. \texttt{[method]}: each section corresponds
  to a step of the recipe.
\end{itemize}

\subsection{\texorpdfstring{\texttt{p-\allowbreak{}card} --- the seven
fields}{ --- the seven fields}}\label{p-card-the-seven-fields}

\begin{itemize}
\tightlist
\item
  \textbf{Assigned to:} the role cards and the three instances (which
  carry cards of their own).
\item
  \textbf{Reads:} per card, the fields of the canon's ``Role card
  fields'' list missing (as bold labels), and placeholders left in a
  field.
\item
  \textbf{Limit:} 0 missing, 0 placeholders. \texttt{[method]}: a role
  without a gate and without a note is a step; the seven fields make
  that verifiable.
\end{itemize}

\subsection{\texorpdfstring{\texttt{p-\allowbreak{}limit} --- justified
limits}{ --- justified limits}}\label{p-limit-justified-limits}

\begin{itemize}
\tightlist
\item
  \textbf{Assigned to:} the gates file and the three instances.
\item
  \textbf{Reads:} gate table rows with an empty magnitude, limit or
  justification; a justification without an origin label or with fewer
  than 6 words (the number is read from the gate's own row); a
  placeholder in a field.
\item
  \textbf{Limit:} 0. \texttt{[method]}: 6 words is the minimum to say
  where the number comes from and why not another one. It is a lexical
  proxy: six words of any kind pass, and the review reads whether they
  justify.
\end{itemize}

\subsection{\texorpdfstring{\texttt{p-\allowbreak{}references} --- the
graph
closes}{ --- the graph closes}}\label{p-references-the-graph-closes}

\begin{itemize}
\tightlist
\item
  \textbf{Assigned to:} the artifact table, the cards, the gates file
  and the three instances.
\item
  \textbf{Reads:} artifacts cited in gates, dependencies or notes a card
  emits and not declared; owners that are neither a role nor
  \texttt{owner}; cycles in the graph (blocking); gates the artifact
  table assigns with no row in the gate table, or rows naming other
  artifacts than it assigns; placeholders in an artifact record; paths
  an artifact's \texttt{command} names that do not exist; Limit cells
  whose numbers are not the numbers of the limit the validator applied.
\item
  \textbf{Limit:} 0. \texttt{[method]}: a loose reference is the written
  form of ``the old text nobody knew had changed''; a cycle means a
  wrong graph.
\end{itemize}

\subsection{\texorpdfstring{\texttt{p-\allowbreak{}prior} --- numbers
about a registered
prior}{ --- numbers about a registered prior}}\label{p-prior-numbers-about-a-registered-prior}

\begin{itemize}
\tightlist
\item
  \textbf{Assigned to:} the principle, the canon, the recipe, the three
  instances, the lessons and the experiments.
\item
  \textbf{Reads:} prior labels whose id is not registered, or whose id
  breaks the grammar
  \texttt{[0-\allowbreak{}9a-\allowbreak{}z-\allowbreak{}]+} (a
  malformed mark, blocking); every number of a unit marked
  \texttt{[prior:\allowbreak{}<id>]}, code spans included, which needs a
  row of that prior in the canon's ``Recorded facts'' and is checked by
  presence in the union of the sections the row cites, in the method
  document passed with \texttt{-\allowbreak{}-\allowbreak{}case} whose
  sha256 is the pin of \texttt{<id>}, read in the
  \texttt{number\_\allowbreak{}format} the registry declares; section
  references (``section 2.1'', ``step 16'') checked by the existence of
  that label; on the canon, the section commitments compared with the
  sections of the documents passed.
\item
  \textbf{Limit:} 0. \texttt{[brief]}: a number absent from the prior's
  method document is an invented fact until proven otherwise.
\item
  \textbf{Pending:} without a registry; without the
  \texttt{-\allowbreak{}-\allowbreak{}case} document of a cited prior
  (pending by name: a cited owner-held prior reads ``owner-held
  evidence, not measured here'').
\item
  \textbf{Exceptions:} the full id of a registered instance
  (\texttt{004-\allowbreak{}adoption-\allowbreak{}guide}) is a name, a
  bare \texttt{004} is a number; a digit attached to a letter
  (\texttt{p95}, \texttt{v7}) is a name; numbers written in words are
  not read. The check verifies the presence of the number's tokens in
  the cited sections, not the claim around it; the ``Recorded facts''
  section of the canon is exempt by name.
\end{itemize}

\subsection{\texorpdfstring{\texttt{p-\allowbreak{}prior-\allowbreak{}mark}
--- statements about a prior carry the
mark}{ --- statements about a prior carry the mark}}\label{p-prior-mark-statements-about-a-prior-carry-the-mark}

\begin{itemize}
\tightlist
\item
  \textbf{Assigned to:} the principle, the recipe, the three instances
  and the experiments.
\item
  \textbf{Reads:} paragraphs, list items and table rows that use a
  phrase of the canon's ``Prior-statement triggers'' without a
  \texttt{[prior:\allowbreak{}.\allowbreak{}.\allowbreak{}.\allowbreak{}]}
  mark outside code spans (a mark in a code span is no mark; a mark
  outside the grammar is no mark either).
\item
  \textbf{Limit:} 0. \texttt{[method]}: created after the first review,
  where a qualitative statement about the prior, absent from its
  document, passed every gate because only numeric statements needed a
  mark.
\item
  \textbf{Exempt:} the frozen ``Rubric'' section of the gates file, by
  design.
\end{itemize}

\subsection{\texorpdfstring{\texttt{p-\allowbreak{}prior-\allowbreak{}pin}
--- the pin of each executed
prior}{ --- the pin of each executed prior}}\label{p-prior-pin-the-pin-of-each-executed-prior}

\begin{itemize}
\tightlist
\item
  \textbf{Assigned to:} the registry.
\item
  \textbf{Reads:} for each executed prior cited in a document that
  carries \texttt{p-\allowbreak{}prior}, the registered sha256 compared
  with the documents passed with
  \texttt{-\allowbreak{}-\allowbreak{}case}, prior by prior: a cited
  owner-held prior whose document was not passed is pending by name; a
  cited distributed prior whose pin matches none of the documents passed
  is blocking; a document passed that is the pin of no executed instance
  is blocking (with such a document even the owner-held prior is
  blocking, since the stray may be that prior, edited); a cited executed
  instance without a pin is blocking.
\item
  \textbf{Limit:} equal. \texttt{[method]}: a pin nobody can re-verify
  is not a pin.
\end{itemize}

\subsection{\texorpdfstring{\texttt{p-\allowbreak{}contract-\allowbreak{}shape}
--- the contracts
hold}{ --- the contracts hold}}\label{p-contract-shape-the-contracts-hold}

\begin{itemize}
\tightlist
\item
  \textbf{Assigned to:} the artifact table, the contracts, the registry,
  the note ledger, the state ledger, the round record and the owner
  decisions.
\item
  \textbf{Reads:} schema cells of the artifact table citing a
  \texttt{*.\allowbreak{}schema.\allowbreak{}json} file that is missing
  or does not parse; keywords of each contract outside the subset the
  validator applies, or a \texttt{contract\_\allowbreak{}version} other
  than the one it reads (\texttt{1}); the registry, the note ledger, the
  state ledger, the round record and the owner decisions validated line
  by line against their schemas; registry entries that break the
  contract of their status; the hash chain of the three append-only
  ledgers (chapter 15) and the rules of a chained line of the round
  record (its \texttt{kind}, its
  \texttt{reviewer\_\allowbreak{}relation}, no verdict on a maintenance
  line, a \texttt{rubric\_\allowbreak{}hash} equal to the frozen rubric
  when \texttt{rubric\_\allowbreak{}applied} is true, the form of
  \texttt{family\_\allowbreak{}map\_\allowbreak{}commitment}).
\item
  \textbf{Limit:} 0. \texttt{[method]}: a contract that exists but is
  never applied is not a contract.
\end{itemize}

\subsection{\texorpdfstring{\texttt{p-\allowbreak{}recipe} --- the steps
of the
recipe}{ --- the steps of the recipe}}\label{p-recipe-the-steps-of-the-recipe}

\begin{itemize}
\tightlist
\item
  \textbf{Assigned to:} the recipe.
\item
  \textbf{Reads:} steps missing one of the fields ``Reads'',
  ``Creates'', ``Does'', ``Done when''; the number of steps (the minimum
  is read from the gate's row).
\item
  \textbf{Limit:} 0 missing a field; ≥ 12 steps. \texttt{[method]}: 12
  is the number of steps in the principle's recipe, which this one
  details.
\end{itemize}

\subsection{\texorpdfstring{\texttt{p-\allowbreak{}components} --- the
eleven components
named}{ --- the eleven components named}}\label{p-components-the-eleven-components-named}

\begin{itemize}
\tightlist
\item
  \textbf{Assigned to:} the three instances.
\item
  \textbf{Reads:} the method components of the canon's list named in the
  document (inflected forms accepted).
\item
  \textbf{Limit:} 11 of 11. \texttt{[method]}: an instance without one
  of them is not the method.
\end{itemize}

\subsection{\texorpdfstring{\texttt{p-\allowbreak{}efficacy-\allowbreak{}mark}
--- claims of efficacy carry their
evidence}{ --- claims of efficacy carry their evidence}}\label{p-efficacy-mark-claims-of-efficacy-carry-their-evidence}

\begin{itemize}
\tightlist
\item
  \textbf{Assigned to:} the principle, the brief, the recipe, the three
  instances, the lessons (from the heading the canon declares on), the
  README and the experiments.
\item
  \textbf{Reads:} units that use a phrase of the canon's ``Efficacy
  triggers'' without one of the labels
  \texttt{[prior:\allowbreak{}<id>]},
  \texttt{[experiment:\allowbreak{}<id>]},
  \texttt{[against:\allowbreak{}<id>]},
  \texttt{[lessons:\allowbreak{}<title>]}, \texttt{[hypothesis]} outside
  code spans; \texttt{[experiment:\allowbreak{}<id>]} and
  \texttt{[against:\allowbreak{}<id>]} labels that do not name an
  executed experiment with a ``Provenance'' table in
  \texttt{60-\allowbreak{}experiments.\allowbreak{}md} (a label citing a
  designed or pending experiment is a note; one without provenance reads
  as \texttt{[hypothesis]}, with a note);
  \texttt{[lessons:\allowbreak{}<title>]} labels that name no level-2 or
  level-3 heading of the lessons file; a \texttt{[hypothesis]} claim
  with a number or a universal quantifier (a low note).
\item
  \textbf{Limit:} 0. \texttt{[method]}: a claim that the method itself
  works is a statement of the same kind as a statement about a prior,
  and exempting it was the double standard a critical panel found.
\end{itemize}

\section{Gates measured once per run}\label{gates-measured-once-per-run}

\subsection{\texorpdfstring{\texttt{p-\allowbreak{}notes} --- the note
ledgers}{ --- the note ledgers}}\label{p-notes-the-note-ledgers}

\begin{itemize}
\tightlist
\item
  \textbf{Assigned to:} \texttt{open-\allowbreak{}notes} and
  \texttt{notes-\allowbreak{}status}.
\item
  \textbf{Reads:} notes of the ledger passed with
  \texttt{-\allowbreak{}-\allowbreak{}notes} missing one of the six
  fields, with a severity outside the four, or with a target outside the
  graph; notes without a line in the state ledger, or with two, or with
  a closing status and no version; evidence naming a file of the
  material (\texttt{.\allowbreak{}md}, \texttt{.\allowbreak{}json},
  \texttt{.\allowbreak{}jsonl}, \texttt{.\allowbreak{}py}) that does not
  exist; the rules of the state (chapter 15): an escaped regression
  closed with neither a gate of the gate table nor a reason (high); an
  open high or blocking note that affects a non-negotiable and no
  decided owner decision covers (blocking); an open note older than 3
  rounds and not scheduled (low).
\item
  \textbf{Limit:} 0. \texttt{[hypothesis]}: a note aimed at an artifact
  the graph does not know is undeliverable; the note router of the first
  instance rejected such a target
  \texttt{[prior:\allowbreak{}000-\allowbreak{}av-\allowbreak{}production]}.
\item
  \textbf{Pending:} without \texttt{-\allowbreak{}-\allowbreak{}notes};
  that pending cannot be accepted.
\end{itemize}

\subsection{\texorpdfstring{\texttt{p-\allowbreak{}frozen-\allowbreak{}rubric}
--- the rubric and the frozen
records}{ --- the rubric and the frozen records}}\label{p-frozen-rubric-the-rubric-and-the-frozen-records}

\begin{itemize}
\tightlist
\item
  \textbf{Assigned to:} the gates file and the review.
\item
  \textbf{Reads:} the sha256 recorded at the launch of the review
  (\texttt{-\allowbreak{}-\allowbreak{}rubric-\allowbreak{}hash})
  against the current sha256 of the rubric as the layout resolves it;
  and the sha256 of each record the canon's ``Frozen records'' table
  pins (the reviewer's report and notes of round 1, the derived note
  ledger, the replies, the pin of the rubric, and the part of the
  lessons above the heading the efficacy gate reads from).
\item
  \textbf{Limit:} equal (blocking when a record changed).
  \texttt{[method]}: a rubric changed after the review makes the verdict
  incomparable. The panel that reviewed v0.5.0 rewrote ``Rejected'' as
  ``Accepted'' in the reviewer's report and no gate saw it; since v0.6.0
  the records are pinned.
\end{itemize}

\subsection{\texorpdfstring{\texttt{p-\allowbreak{}sensitivity} --- the
gate on the
gates}{ --- the gate on the gates}}\label{p-sensitivity-the-gate-on-the-gates}

\begin{itemize}
\tightlist
\item
  \textbf{Assigned to:} \texttt{instrument-\allowbreak{}validator}.
\item
  \textbf{Reads:} the sensitivity report named by the layout: a missing
  report (blocking); an incomplete one; a table that does not add up to
  its ``Caught'' line; a provenance that does not match (the running
  validator and the one next to the report, the ``Gates'' section, the
  canon, the resolved gate map, the perturbation set, the
  \texttt{case sha256} of each prior document and its
  \texttt{case attestation}); known defects the expected gate let
  through (missed) or that changed another row (noisy); an assigned gate
  without a perturbation, or fewer than 25 \% of the (document, gate)
  pairs perturbed (pending with the reason, never a pass); a report made
  with a synthetic prior (pending under
  \texttt{p-\allowbreak{}sensitivity-\allowbreak{}synthetic} in a run
  without that prior's document, refused in a run with it).
\item
  \textbf{Limit:} every gate covered, ≥ 25 \% of pairs, 0 missed, 0
  noisy. \texttt{[method]}: a gate that never fails proves nothing, and
  a report this instrument cannot have produced proves nothing either;
  \texttt{[hypothesis]}: the floor of a quarter of the pairs is the
  first this version commits to, above the 18 \% a panel measured, to be
  raised and never lowered. The corpus covers 50 of 182 pairs.
\end{itemize}

\subsection{\texorpdfstring{\texttt{p-\allowbreak{}declared} --- what is
declared exists, and what exists is
declared}{ --- what is declared exists, and what exists is declared}}\label{p-declared-what-is-declared-exists-and-what-exists-is-declared}

\begin{itemize}
\tightlist
\item
  \textbf{Assigned to:} the artifact table.
\item
  \textbf{Reads:} declared paths that match nothing on disk (each
  literal path, and each glob at least once; blocking, naming the path);
  Markdown documents of the folder no artifact declares and the canon
  does not list as outside the table; files that belong to two
  artifacts; documents of the ``Size ranges'' table no artifact
  declares; a declared artifact marked \texttt{proposed} (pending with
  its reason) or declared with no paths (pending).
\item
  \textbf{Limit:} 0. \texttt{[method]}: a declared document that
  silently matches no file is a promise the instrument never kept.
\end{itemize}

\section{The measurement of the
instrument}\label{the-measurement-of-the-instrument}

\subsection{\texorpdfstring{\texttt{m-\allowbreak{}determinism}}{}}\label{m-determinism}

Not a gate on a document: a measurement of the validator itself, which
cannot fail a document. It is the equality of the report's sha256 across
two consecutive runs on the same input, the second with another hash
seed and the C locale; measured only with
\texttt{-\allowbreak{}-\allowbreak{}determinism}, and a difference is a
high note. \texttt{[hypothesis]}: a report that changes by itself is not
evidence; the continuity ledger the validator of the first instance
wrote was deterministic
\texttt{[prior:\allowbreak{}000-\allowbreak{}av-\allowbreak{}production]}.

\section{Domain gates}\label{domain-gates}

\subsection{\texorpdfstring{\texttt{d-\allowbreak{}<name>}}{}}\label{d-name}

A gate of the target's domain, declared as a row of the cell's gate
table (every gate table of the gates file is read, so domain gates may
keep a table of their own) and never run by the validator. The cell's
own code measures it and writes a result file in the contract
\texttt{domain-\allowbreak{}result.\allowbreak{}schema.\allowbreak{}json};
the layout lists the result files under
\texttt{domain-\allowbreak{}results}. The validator reads each result
and:

{\def\LTcaptype{none} % do not increment counter
\begin{longtable}[]{@{}
  >{\raggedright\arraybackslash}p{(\linewidth - 2\tabcolsep) * \real{0.5000}}
  >{\raggedright\arraybackslash}p{(\linewidth - 2\tabcolsep) * \real{0.5000}}@{}}
\toprule\noalign{}
\begin{minipage}[b]{\linewidth}\raggedright
Case
\end{minipage} & \begin{minipage}[b]{\linewidth}\raggedright
What the validator does
\end{minipage} \\
\midrule\noalign{}
\endhead
\bottomrule\noalign{}
\endlastfoot
an input changed or is missing since the result was written (its sha256
recomputed) & blocking: the result is stale \\
the result names no input & blocking \\
the inputs do not include every file of the artifacts the gate's row
names & blocking (a result over the canon alone passed while the
candidate changed) \\
the result's own numbers contradict it, when the limit starts with a
comparison and a number (\texttt{≤ 2 s},
\texttt{at least 1.\allowbreak{}5 x}) or is a bare zero
(\texttt{0 divergent}) & blocking \\
\texttt{result:\allowbreak{} fail} & a note with the severity of the
gate's row (an optional Severity column, else high); a severity below
high counts only for a gate a decided owner decision covers \\
\texttt{result:\allowbreak{} tool-\allowbreak{}failure}, or a result
that breaks the contract & exit 2 \\
a declared \texttt{d-\allowbreak{}*} gate with no result & pending,
which no caller can accept \\
\end{longtable}
}

Declared reach: a result whose inputs were recomputed by hand, the
candidate changed and its new sha256 written in, passes; it is the same
act as running the gate again without running it, which the validator,
reading results and never running the cell's code, cannot tell apart.
The sensitivity of a domain gate is the cell's own (mutants, a reference
with a known defect), and the coverage of
\texttt{p-\allowbreak{}sensitivity} leaves domain gates out. Chapter 21
shows how a cell writes one.

\section{The false greens}\label{the-false-greens}

\texttt{06-\allowbreak{}gate-\allowbreak{}semantics.\allowbreak{}md}
keeps a failure-mode analysis of the validator's gates for the one
failure mode that matters to a verdict: the \textbf{false green}, a
defect the gate lets through with a pass. For each gate it names the
cause, what detects it today, the probe that first found it and the
perturbation that covers it now. The table is reproduced in appendix D;
the summary is that most of the causes found so far came from the
critical panels that reviewed v0.4.2 and v0.5.0 and from the independent
review of v0.6.0, each probe turned into a regression test and, where it
is a defect of a document, a perturbation. The causes are the ones found
so far, not an exhaustive list; a new probe that passes a gate adds a
row and a perturbation.

\chapter{The sensitivity test}\label{the-sensitivity-test-1}

\section{Purpose}\label{purpose-1}

A gate that never fails proves nothing.
\texttt{sensitivity.\allowbreak{}py} is the gate on the gates: it
injects known defects, one at a time, into a copy of the material and
checks that the expected gate rejects each on the expected document, and
nothing else moves. Its report, \texttt{sensitivity.\allowbreak{}md}, is
an artifact the validator reads (\texttt{p-\allowbreak{}sensitivity}),
with a provenance header that ties it to the exact validator, gate
table, canon and gate map it was made against. The validator refuses a
report that the instrument cannot have produced, so the test cannot be
satisfied by an old report or a hand-written one.

\begin{figure}
\centering
\pandocbounded{\includegraphics[keepaspectratio,alt={One run of the sensitivity test, the gate on the gates.}]{diagrams/sensitivity-flow.pdf}}
\caption{One run of the sensitivity test, the gate on the
gates.}\label{fig-sensitivity-flow}
\end{figure}

\section{A perturbation}\label{a-perturbation}

The corpus keeps its set in
\texttt{perturbation\_\allowbreak{}set.\allowbreak{}py} (65
perturbations at 0.6.0), a module the validator reads without importing
the sensitivity script; a cell keeps its own as data,
\texttt{08-\allowbreak{}validator/\allowbreak{}perturbations.\allowbreak{}json},
named by the layout key \texttt{artifacts.\allowbreak{}perturbations}.
The fields:

{\def\LTcaptype{none} % do not increment counter
\begin{longtable}[]{@{}
  >{\raggedright\arraybackslash}p{(\linewidth - 2\tabcolsep) * \real{0.2500}}
  >{\raggedright\arraybackslash}p{(\linewidth - 2\tabcolsep) * \real{0.7500}}@{}}
\toprule\noalign{}
\begin{minipage}[b]{\linewidth}\raggedright
Field
\end{minipage} & \begin{minipage}[b]{\linewidth}\raggedright
Meaning
\end{minipage} \\
\midrule\noalign{}
\endhead
\bottomrule\noalign{}
\endlastfoot
\texttt{name} & unique, without a comma (the report lists the names on
one line) \\
\texttt{document} & the document to perturb, relative to the folder;
\texttt{.\allowbreak{}.\allowbreak{}/\allowbreak{}instances/\allowbreak{}registry.\allowbreak{}json}
is the copied registry \\
\texttt{old}, \texttt{new} & the text replaced once in the document;
with \texttt{delete:\allowbreak{} true} the document is removed
instead \\
\texttt{gate} & the gate that has to react on that document; one the
validator implements, never a domain gate \\
\texttt{expect} & \texttt{fail} (default), or \texttt{pending} for a
defect that empties what the gate reads (one perturbation of the corpus
empties a list of the canon) \\
\texttt{needs\_\allowbreak{}case},
\texttt{needs\_\allowbreak{}rubric\_\allowbreak{}hash},
\texttt{needs\_\allowbreak{}notes},
\texttt{needs\_\allowbreak{}registry} & the validator inputs the
expected gate cannot run without; a row whose input was not passed is
not measured, never missed \\
\texttt{prior} & the registered prior whose method document the gate
needs; without a \texttt{-\allowbreak{}-\allowbreak{}case} document
identified as it, the row is not measured, whatever other prior was
passed \\
\texttt{also} & gates the defect also reaches by design,
comma-separated; their rows are not counted as noise, and the report
lists them (eleven perturbations of the corpus declare one: a pin
decides which document is the prior, a list of the canon feeds a gate on
every document, a placeholder in angle brackets also reads as an HTML
tag, an edited old line of a chained ledger breaks its chain) \\
\end{longtable}
}

Two examples from the corpus's set:

\begin{Shaded}
\begin{Highlighting}[]
\NormalTok{Perturbation("avoid{-}word inserted", "10{-}generic{-}recipe.md",}
\NormalTok{             "\#\# The orchestrator in pseudocode",}
\NormalTok{             "\#\# The orchestrator of the pipeline", "p{-}terms")}
\NormalTok{Perturbation("prior number altered", "10{-}generic{-}recipe.md",}
\NormalTok{             "the validator gave 38", "the validator gave 39", "p{-}prior",}
\NormalTok{             needs\_case=True, prior="000{-}av{-}production")}
\end{Highlighting}
\end{Shaded}

\section{The copy is the material}\label{the-copy-is-the-material}

Since v0.6.0 each perturbation runs on a copy of the whole material the
folder belongs to: the folder and everything its artifact table and
registry reach above it (the principle, the root README, the skill, the
changelog, the registry). Until then the copy was the folder alone, so
no perturbation could reach those documents. The copies go to
\texttt{-\allowbreak{}-\allowbreak{}work SCRATCH\_\allowbreak{}FOLDER},
or to a temporary folder; \texttt{.\allowbreak{}git},
\texttt{\_\allowbreak{}\_\allowbreak{}pycache\_\allowbreak{}\_\allowbreak{}}
and \texttt{.\allowbreak{}pytest\_\allowbreak{}cache} are not copied.

\section{Specificity, coverage and the
floor}\label{specificity-coverage-and-the-floor}

The perturbed report is compared with the baseline row by row. A
perturbation is \textbf{caught} when the expected (document, gate) row
went from pass to fail (or to pending, when \texttt{expect} says so),
\textbf{missed} when it did not, and \textbf{noisy} when another row
changed its result --- not counting the expected row, the row of
\texttt{p-\allowbreak{}sensitivity} itself (which reads the provenance
an edit of an input pins on purpose), the rows of a document that left
the report, and the gates the perturbation declares in \texttt{also}. A
noisy perturbation fails the gate like a missed one.

Coverage is measured twice: by gate (every gate the artifact table
assigns, except \texttt{p-\allowbreak{}sensitivity}, the measurements
and the domain gates, has at least one perturbation) and by pair (the
(document, gate) rows of the baseline report that at least one measured
perturbation targets). The magnitude reads ``k of n gates covered, p of
q pairs (r \%), 0 missed of m, 0 noisy''; at 0.6.0 the corpus reads 20
of 20 gates and 50 of 182 pairs (27 \%). An assigned gate without a
perturbation, or fewer than 25 \% of the pairs, leaves
\texttt{p-\allowbreak{}sensitivity} pending with the reason, never a
pass; the floor is a hypothesis, the first this version commits to,
raised and never lowered. The three generic perturbations of a new cell
cover three gates and three pairs; each gate a cell adds needs one of
its own.

\section{Provenance}\label{provenance}

The header of the report records what it was made against, and the
validator recomputes every line and compares:

\begin{Shaded}
\begin{Highlighting}[]
\NormalTok{validator sha256: \textless{}sha of validator.py: the running one, and the one next to the report\textgreater{}}
\NormalTok{gates sha256: \textless{}sha of the "Gates" section; a cell\textquotesingle{}s whole gates file\textgreater{}}
\NormalTok{canon sha256: \textless{}sha of the canon\textgreater{}}
\NormalTok{map sha256: \textless{}sha of the sorted (gate, document path) pairs of the gate map\textgreater{}}
\NormalTok{perturbations: \textless{}count\textgreater{} (\textless{}names, comma{-}separated\textgreater{})}
\NormalTok{case sha256: \textless{}prior id\textgreater{} \textless{}sha of the document passed\textgreater{}   (one per prior; "none")}
\NormalTok{case attestation: \textless{}prior id\textgreater{} \textless{}sha of the text + the four digests above\textgreater{}}
\end{Highlighting}
\end{Shaded}

Editing the validator, the ``Gates'' section, the canon or the gate map
makes the report stale: regenerate it (and then the validator report)
before the next run, or \texttt{p-\allowbreak{}sensitivity} fails on
provenance. Deleting the report does not help: a missing report is
blocking, and \texttt{p-\allowbreak{}sensitivity} cannot be accepted by
name.

\section{The synthetic prior, for a
clone}\label{the-synthetic-prior-for-a-clone}

Five rows of the corpus's set need the method document of the first
prior, which the owner keeps. A maintainer who edits the instrument
without it regenerates a full report with
\texttt{-\allowbreak{}-\allowbreak{}synthetic-\allowbreak{}prior 000-\allowbreak{}av-\allowbreak{}production}:
the script generates a document from the canon's recorded facts (one
numbered section per label the facts cite, holding the facts' numbers in
the prior's number format), in the \texttt{heading\_\allowbreak{}style}
the registry declares for the real document, passes it as
\texttt{-\allowbreak{}-\allowbreak{}case}, and patches its sha256 and
its section digests into the registry and the canon of every copy, never
of the material.

The validator tells a synthetic prior by its \texttt{case sha256}
differing from the registry's pin, never by a line that declares it (the
first draft's declaring line could be deleted, and the report then read
as the owner's). A report made with a synthetic prior is never a pass:
in a run without that prior's document its row is
\texttt{p-\allowbreak{}sensitivity-\allowbreak{}synthetic}, pending,
accepted only when the caller names that gate; in a run with the
document it is refused. It shows that every gate fails its known defect;
it does not show that the owner's document holds the facts, nor that the
validator reads that document's format beyond what the registry
declares.

\section{The attestation, and what it
leaves}\label{the-attestation-and-what-it-leaves}

The pin of the registry is public, so a \texttt{case sha256} line can be
written by hand with the real pin over a report made with the synthetic
document. Until the closing pass of v0.6.0 that forged line passed every
run. Since then the report records, for each document it was made with,
a \texttt{case attestation}: the sha256 of the document's text followed
by the validator, gates, canon and map digests of the report, which only
whoever holds the text can compute. A run that passes the document
recomputes it and refuses a mismatch; that is the owner's run for
\texttt{000} and every run for \texttt{004}. What it leaves, declared: a
run without the document compares nothing, so \textbf{in a clone the
forged line still passes}, and a frozen record re-pinned in the canon
shows there only in the diff against a signed version.

\chapter{Ledgers, records and
integrity}\label{ledgers-records-and-integrity}

\section{The three ledgers}\label{the-three-ledgers}

Three append-only ledgers under \texttt{state/\allowbreak{}} hold the
state the gates read, one JSON object per line, each with a contract
(appendix B):

{\def\LTcaptype{none} % do not increment counter
\begin{longtable}[]{@{}
  >{\raggedright\arraybackslash}p{(\linewidth - 4\tabcolsep) * \real{0.3333}}
  >{\raggedright\arraybackslash}p{(\linewidth - 4\tabcolsep) * \real{0.3333}}
  >{\raggedright\arraybackslash}p{(\linewidth - 4\tabcolsep) * \real{0.3333}}@{}}
\toprule\noalign{}
\begin{minipage}[b]{\linewidth}\raggedright
Ledger
\end{minipage} & \begin{minipage}[b]{\linewidth}\raggedright
One line per
\end{minipage} & \begin{minipage}[b]{\linewidth}\raggedright
Written by
\end{minipage} \\
\midrule\noalign{}
\endhead
\bottomrule\noalign{}
\endlastfoot
\texttt{round-\allowbreak{}record.\allowbreak{}jsonl} & round: method
version, round label, the executor of each role by family and effort, a
summary, and since v0.6.0 its \texttt{kind}, the
\texttt{reviewer\_\allowbreak{}relation},
\texttt{rubric\_\allowbreak{}applied},
\texttt{rubric\_\allowbreak{}hash},
\texttt{rubric\_\allowbreak{}author}, the
\texttt{family\_\allowbreak{}map\_\allowbreak{}commitment}, and optional
cost fields (tokens per role, calls, finish reason,
\texttt{wall\_\allowbreak{}seconds}, \texttt{notes\_\allowbreak{}count},
\texttt{crew}, \texttt{rounds\_\allowbreak{}to\_\allowbreak{}stop}) &
the verifier, at the close of a round \\
\texttt{notes-\allowbreak{}status.\allowbreak{}jsonl} & state change of
a tracked note: \texttt{ref} (\texttt{notes:\allowbreak{}<n>} for the
n-th note of the frozen ledger,
\texttt{panel:\allowbreak{}<version>:\allowbreak{}<id>} for a panel
finding), \texttt{status} (\texttt{open}, \texttt{fixed},
\texttt{disputed}, \texttt{out of scope}),
\texttt{closed\_\allowbreak{}in}, \texttt{summary}, and the fields of
the state rules below & the verifier \\
\texttt{owner-\allowbreak{}decisions.\allowbreak{}jsonl} & decision only
the owner can take: \texttt{name}, \texttt{status} (\texttt{open},
\texttt{decided}), \texttt{blocks}, \texttt{version}, \texttt{summary},
and once decided the notes it \texttt{covers} & the owner, through the
orchestrator \\
\end{longtable}
}

\section{The hash chain}\label{the-hash-chain}

Each line written from v0.6.0 on carries \texttt{prev}, the sha256 of
every line above it, each followed by LF (for a file with no blank line,
the sha256 of the file before the line was appended).
\texttt{p-\allowbreak{}contract-\allowbreak{}shape} recomputes the
chain: an earlier line edited, removed or reordered breaks the first
chained line after it. The canon's ``Chained ledgers'' table counts the
lines written before each chain (10 in the round record, 48 in the state
ledger, 0 in the owner decisions), which are never edited, and every
line after them has to carry \texttt{prev}, so stripping the chain is a
finding too.

The reach is declared: the chain detects an edit made without
recomputing it --- the careless maintainer and the honest error. It does
not detect an edit that recomputed it, which costs one sha256 per line,
and it does not cover the last line of a ledger, which can be rewritten
with no sign until a line is appended after it. Only an anchor outside
the tree --- a signed tag per version, which is the owner's to make ---
answers that.

\section{The rules of the state of a
note}\label{the-rules-of-the-state-of-a-note}

Since v0.6.0 the last line of each tracked note is read against four
rules (\texttt{p-\allowbreak{}notes}):

\begin{itemize}
\tightlist
\item
  \textbf{Golden rule 4.} A note that is a regression that escaped every
  gate (\texttt{escaped:\allowbreak{} true}) closes with the gate it
  became (\texttt{gate}, a row of the cell's gate table) or with the
  recorded reason not to make one (\texttt{no\_\allowbreak{}gate}, at
  least 8 distinct words, or \texttt{decision} exactly the name of a
  decided owner decision); closed with neither, it is a high note. The
  count of words measures the form of the reason, not its content.
\item
  \textbf{Age.} An open note is counted in rounds from the version it
  was opened in (\texttt{opened\_\allowbreak{}in}, or the version of a
  panel ref, or 0.1.0 for the frozen ledger). Past 3 rounds, and not
  scheduled for a version still ahead (\texttt{due}), it is a low note:
  schedule it, dispute it or close it with the reason. Three rounds is a
  hypothesis: a panel counted 16 of the 22 notes of round 1 open nine
  rounds later.
\item
  \textbf{Non-negotiables.} The brief gives each non-negotiable an id in
  a code span (\texttt{nn-\allowbreak{}<name>}); a note names those it
  affects (\texttt{affects}). An open note of severity high or blocking
  that affects one is blocking, unless a decided entry of the owner
  decisions covers it (\texttt{covers}). The fields a line's writer
  controls cannot undo the protection: the severity is the highest of
  the note ledger's and of every line of the ref; the non-negotiables
  are the union of every line's \texttt{affects}; a decision covers only
  while its last line is \texttt{decided}; a high or blocking note
  opened in 0.6.0 or later has to name what it affects (an empty list
  when nothing).
\item
  \textbf{Evidence.} A note whose evidence names a file of the material
  needs that file to exist; the frozen ledger keeps the names of its
  time and is exempt.
\end{itemize}

\section{Frozen records}\label{frozen-records}

The records of the first review and the history of the lessons are never
edited. The canon's ``Frozen records'' table pins their sha256
(\texttt{90-\allowbreak{}review/\allowbreak{}auditor-\allowbreak{}report.\allowbreak{}md},
\texttt{auditor-\allowbreak{}notes.\allowbreak{}json},
\texttt{notes.\allowbreak{}json}, \texttt{replies.\allowbreak{}md},
\texttt{rubric.\allowbreak{}sha256}, and the part of
\texttt{99-\allowbreak{}lessons.\allowbreak{}md} above the heading ``The
v0.5.0 round: exit 0 is not green''), and
\texttt{p-\allowbreak{}frozen-\allowbreak{}rubric} compares them on
every run. An append-only rule that no gate checked is how the panel
that reviewed v0.5.0 rewrote a verdict unseen.

\section{The family map and its
commitment}\label{the-family-map-and-its-commitment}

The round record names the executors by family letter; the map from
letter to family, and a salt, stay with the owner. Since the rework of
v0.6.0 a chained line of the round record of method version 0.6.0 or
later carries
\texttt{family\_\allowbreak{}map\_\allowbreak{}commitment}, the sha256,
over UTF-8, of

\begin{Shaded}
\begin{Highlighting}[]
\NormalTok{json.dumps(map, sort\_keys=True, separators=(\textquotesingle{},\textquotesingle{}, \textquotesingle{}:\textquotesingle{}), ensure\_ascii=False) + \textquotesingle{}\textbackslash{}n\textquotesingle{} + salt}
\end{Highlighting}
\end{Shaded}

where \texttt{map} is the object from letter to the description of the
family and \texttt{salt} a string of 64 hex characters.
\texttt{p-\allowbreak{}contract-\allowbreak{}shape} checks the form only
(64 hex characters); the owner proves the commitment by revealing the
map and the salt, and nothing in the material lets anyone else check it.
A pseudonym made of the hash of a host name or a user name is not
anonymity (a dictionary reverses it); this is why the material uses
letters and a salted commitment.

\section{What a clone verifies, and what rests on the
owner}\label{what-a-clone-verifies-and-what-rests-on-the-owner}

\begin{figure}
\centering
\pandocbounded{\includegraphics[keepaspectratio,alt={What pins what: the integrity structure of the material at version 0.6.0.}]{diagrams/trust-chain.pdf}}
\caption{What pins what: the integrity structure of the material at
version 0.6.0.}\label{fig-trust-chain}
\end{figure}

{\def\LTcaptype{none} % do not increment counter
\begin{longtable}[]{@{}
  >{\raggedright\arraybackslash}p{(\linewidth - 2\tabcolsep) * \real{0.5000}}
  >{\raggedright\arraybackslash}p{(\linewidth - 2\tabcolsep) * \real{0.5000}}@{}}
\toprule\noalign{}
\begin{minipage}[b]{\linewidth}\raggedright
A clone verifies
\end{minipage} & \begin{minipage}[b]{\linewidth}\raggedright
Only the owner verifies
\end{minipage} \\
\midrule\noalign{}
\endhead
\bottomrule\noalign{}
\endlastfoot
every gate on every document; the pin and the numbers of \texttt{004},
whose document is distributed; the frozen rubric; the provenance of the
sensitivity report, the sha256 of each prior document included; the hash
chain of the three ledgers & the numbers stated about \texttt{000},
whose document is owner-held (its cited sections are committed by
sha256, so it cannot be swapped unseen, but a clone cannot read them);
that the sensitivity report was made with that document, by regenerating
it \\
that a ledger line written since v0.6.0 was not edited without
recomputing the chain & an edit that recomputed the chain; the lines
before the chain; the history of the lessons and the record of the
experiments' raw files \\
the family letter each round line records, and the form of its salted
commitment & the map from a family letter to a family, and its salt \\
\end{longtable}
}

The residual risks the material declares, as they stand at 0.6.0:

\begin{itemize}
\tightlist
\item
  in a clone without the owner's document, a forged \texttt{case sha256}
  line passes;
\item
  the 8 distinct words of a \texttt{no\_\allowbreak{}gate} reason
  measure form, not content;
\item
  the attestation does not prove the perturbations actually ran, only
  that the report was made over those inputs;
\item
  the git history distributes the name of the author and the dates of
  the commits; the gates measure the content, not the history;
\item
  the anchor outside the tree, a signed tag per version, is recommended
  and not yet made (an open owner decision).
\end{itemize}

\chapter{The instances}\label{the-instances}

\section{Instance 000: the audiovisual production
method}\label{instance-000-the-audiovisual-production-method}

The first executed instance, from which the method was generalized: an
audiovisual production method, written in Brazilian Portuguese, held by
the owner and not distributable. The registry pins its method document
by sha256 and declares its number format
(\texttt{comma-\allowbreak{}decimal}) and its heading style (a section
\texttt{\#\# 6.\allowbreak{}}, a subsection
\texttt{\#\#\# 1.\allowbreak{}1}, a step
\texttt{\#\#\# Passo 15.\allowbreak{}}), and the canon commits, section
by section, to the text of the twelve sections the recorded facts cite.
The facts the material states about it are the rows of the canon's
``Recorded facts'' table --- among them 18 roles with a card, 49
artifacts in the graph of its first pilot, 2 maximum rounds per
artifact, the review of v1 with 3 blocking, 4 high, 12 medium and 12 low
findings, 26 low notes open at v5, a drift gate whose wrong instrument
read a p95 of 0.546 s at v1 and 0.10 s once fixed, and an automatic
reviewer of which 10 of 12 notes were false
\texttt{[prior:\allowbreak{}000-\allowbreak{}av-\allowbreak{}production]}.
Several rules of the recipe carry the error of this instance that the
step prevents, as a measured size of the error.

Its historical vocabulary is recorded in the canon so that the frozen
records of round 1 can be read: the validator was its ``continuity
checker''; the verdicts were ``publishable'', ``publishable with minor
corrections'' and ``rejected''; the production cell carried a media
term; the quarantine mark of a bracketed word became
\texttt{[prior:\allowbreak{}<id>]}.

\section{Instance 004: the adoption
guide}\label{instance-004-the-adoption-guide}

The second executed instance: the adoption guide of the method itself,
in Brazilian Portuguese, of 900 to 1,400 words by its brief, produced
under method version 0.2.2 with a writer of one family and a review in a
separate dispatch without the producer's reasoning, which gave
\texttt{accepted} with no dimension below 4 and a guide of 1,054 words
\texttt{[prior:\allowbreak{}004-\allowbreak{}adoption-\allowbreak{}guide]}.
That review was first recorded as coming from another family; it came
from the producer's own model through another route, so it separated the
session and the reasoning, not the family.

Its cell is distributed under
\texttt{instances/\allowbreak{}004-\allowbreak{}adoption-\allowbreak{}guide/\allowbreak{}}
since v0.5.0, as a \textbf{record, not an example to copy}:

\begin{itemize}
\tightlist
\item
  it does not run under the current validator without a migration (the
  migration notes of \texttt{CHANGELOG.\allowbreak{}md});
\item
  its own copy of the validator,
  \texttt{08-\allowbreak{}validator/\allowbreak{}validator.\allowbreak{}py},
  is the instrument of 0.2.2 and since v0.6.0 refuses to run unless
  called with \texttt{-\allowbreak{}-\allowbreak{}historical}, because
  its exit 0 does not mean what the current validator's does (under
  0.2.2 a gate that was not measured left the run with exit 0); that
  flag is the only change made to the cell since it was distributed;
\item
  its \texttt{ERRATA.\allowbreak{}md} lists what the guide says that is
  no longer true when read against the material as it stands (the rubric
  is frozen before the first review, not before the production; the
  registry holds two executed instances; the skill copied alone creates
  a cell without the validator; independence across families is a
  recommendation, not a rule);
\item
  its list of forbidden identifiers stayed with the owner, outside the
  material, since it contains what it forbids.
\end{itemize}

The cell shows the shape of a JSON-layout cell at 0.2.2:
\texttt{00-\allowbreak{}brief.\allowbreak{}md},
\texttt{01-\allowbreak{}sources-\allowbreak{}of-\allowbreak{}truth.\allowbreak{}md},
\texttt{02-\allowbreak{}canon.\allowbreak{}json},
\texttt{03-\allowbreak{}artifacts.\allowbreak{}json},
\texttt{04-\allowbreak{}roles/\allowbreak{}} (four cards),
\texttt{05-\allowbreak{}contracts/\allowbreak{}},
\texttt{06-\allowbreak{}gates.\allowbreak{}md},
\texttt{07-\allowbreak{}rubric.\allowbreak{}md} with its pin,
\texttt{08-\allowbreak{}validator/\allowbreak{}} with its reports,
\texttt{09-\allowbreak{}engine/\allowbreak{}},
\texttt{review/\allowbreak{}} (two wave reports and the round-1 notes),
\texttt{state/\allowbreak{}}, \texttt{lessons.\allowbreak{}md},
\texttt{open-\allowbreak{}notes.\allowbreak{}json},
\texttt{priors.\allowbreak{}json}, \texttt{layout.\allowbreak{}json},
the guide, the method document and the errata.

Whether a new cell under 0.6.0 replaces it as the example is an open
owner decision (appendix F).

\section{The three designed
instances}\label{the-three-designed-instances}

Three instances were chosen because they differ from each other, and
each fills the fourteen rows of the instantiation table and the fourteen
mandatory sections with domain content marked \texttt{[hypothesis]}
where it is an assumption. They are designed sketches, not executions,
and carry the design, not a result.

{\def\LTcaptype{none} % do not increment counter
\begin{longtable}[]{@{}
  >{\raggedright\arraybackslash}p{(\linewidth - 6\tabcolsep) * \real{0.2500}}
  >{\raggedright\arraybackslash}p{(\linewidth - 6\tabcolsep) * \real{0.2500}}
  >{\raggedright\arraybackslash}p{(\linewidth - 6\tabcolsep) * \real{0.2500}}
  >{\raggedright\arraybackslash}p{(\linewidth - 6\tabcolsep) * \real{0.2500}}@{}}
\toprule\noalign{}
\begin{minipage}[b]{\linewidth}\raggedright
Element
\end{minipage} & \begin{minipage}[b]{\linewidth}\raggedright
Car (\texttt{20-\allowbreak{}instance-\allowbreak{}car.\allowbreak{}md})
\end{minipage} & \begin{minipage}[b]{\linewidth}\raggedright
Text about 2027
(\texttt{21-\allowbreak{}instance-\allowbreak{}text-\allowbreak{}2027.\allowbreak{}md})
\end{minipage} & \begin{minipage}[b]{\linewidth}\raggedright
HPC library
(\texttt{22-\allowbreak{}instance-\allowbreak{}hpc.\allowbreak{}md})
\end{minipage} \\
\midrule\noalign{}
\endhead
\bottomrule\noalign{}
\endlastfoot
Target & the design of a battery-electric passenger vehicle released for
tooling, with approved validation prototypes and a homologation dossier
& a publishable nonfiction manuscript of about 90 thousand words in 12
chapters on the world in 2027, saying on every page which statement is
fact, projection or opinion & a C++ library with a C interface that
replaces the reference implementation of a 3D finite-difference
operator, with measured speedup and byte-for-byte identical outputs \\
Owner & the program director & the author & the library's technical
lead \\
Sources of truth & requirement → brief and regulatory matrix; geometry →
released CAD; mass → weighing; cost → supplier quotation; physics →
correlated simulation; standard → the regulation's text & fact → primary
source; projection → declared model; number → claims register; quotation
→ original; figure → data plus script & correctness → oracle (frozen
reference and hashed outputs); performance → measurement on a pinned
node with the configuration obtained recorded; where the time goes →
attribution instrument \\
Main artifacts & requirements, architecture, budgets, subsystem CAD,
virtual vehicle, BOM, simulations, safety analysis, validation plan,
test reports, homologation dossier & thesis, theme map, domain dossier,
claims register, projections, outline, scenarios, chapters, figures,
fact-check, manuscript & specification, interface, oracle, attribution
report, lever decision, experiment plan, implementation, tests, build,
equivalence, measurements, performance report, documentation \\
Main gates & mass against target with reserve, range with the real mass,
charging time, clearances, crash margin, thermal propagation,
traceability, coverage, quoted cost, simulation × test correlation &
every claim with a register identifier, every projection marked, primary
source, source density, readability, size per chapter, overlap, recheck
of volatile claims & byte-for-byte equivalence, independence from order
and thread count, sensitivity of the gate to mutations, configuration
obtained, noise, speedup with confidence interval, size sweep,
prediction written beforehand, sanitizers, third-party reproduction \\
Review rubric & a design review by a panel from outside the program &
blind reading by readers from the target audience & code review and
benchmark reproduction by someone who did not implement \\
What stays with the human & targets, styling, supplier choice, residual
risk, safety sign-off, the decision to launch & thesis, tone, opinion,
the label of projection or opinion, legal risk, release & what to
optimize, the target, acceptable memory, new dependencies, the budget,
the release \\
\end{longtable}
}

What each sketch adds to the method's vocabulary: the car shows the
gates of a physical product and the fidelity ladder of its versions
(paper, virtual, engineering prototype, validation prototype); the text
shows the separation of fact, projection and opinion as a contract of
the claims register and a gate on the chapters; the HPC library shows
the danger moving from the role that writes badly to the instrument that
measures wrongly and passes --- the equivalence gate that never rejects,
the benchmark that records the threads requested and not the threads
obtained, the speedup measured at an extreme point and generalized ---
and makes ``measure what the system gave, not what was asked for'' the
first gate of the loop.

\section{The registry as a contract}\label{the-registry-as-a-contract}

\texttt{instances/\allowbreak{}registry.\allowbreak{}json} is one list;
the status of an entry decides its fields, and
\texttt{p-\allowbreak{}contract-\allowbreak{}shape} enforces that
contract on every run. A cell keeps its own registry,
\texttt{priors.\allowbreak{}json}, in the same contract, and registers
itself in it at the close of its production (step 14 of the recipe),
which is also what makes the token of its own id read as a name rather
than as a number.

\part{Using the method}

\chapter{Quick start}\label{quick-start}

\section{The ten-minute path}\label{the-ten-minute-path}

The root \texttt{README.\allowbreak{}md} gives a four-step path; here it
is with the commands run from the root of the material.

\begin{enumerate}
\def\labelenumi{\arabic{enumi}.}
\item
  Read the first section and section 7 of
  \texttt{01-\allowbreak{}principle.\allowbreak{}md} (the idea and the
  recipe).
\item
  Scaffold a cell:

\begin{Shaded}
\begin{Highlighting}[]
\ExtensionTok{python3}\NormalTok{ skills/cooperatio{-}teamwork/scripts/scaffold.py }\DataTypeTok{\textbackslash{}}
    \AttributeTok{{-}{-}target} \StringTok{"My target"} \AttributeTok{{-}{-}roles}\NormalTok{ writer reviewer }\AttributeTok{{-}{-}out}\NormalTok{ ./teamwork}
\end{Highlighting}
\end{Shaded}
\item
  Fill the brief, give each non-negotiable an id
  (\texttt{nn-\allowbreak{}<name>}), fill the lists of
  \texttt{02-\allowbreak{}canon.\allowbreak{}json} the validator reads
  and one card per role (the skill's
  \texttt{references/\allowbreak{}templates.\allowbreak{}md} has a
  filled example).
\item
  Freeze the rubric, run the sensitivity test, then the validator, with
  the commands of the skill (``The commands of a cell'', reproduced in
  chapter 19); read every pending row before reading the exit code.
\end{enumerate}

\section{What the scaffold writes}\label{what-the-scaffold-writes}

With two roles the scaffold writes 24 files, refusing to overwrite
anything that exists; of the folders it creates,
\texttt{09-\allowbreak{}engine/\allowbreak{}} and
\texttt{review/\allowbreak{}} stay empty:

{\def\LTcaptype{none} % do not increment counter
\begin{longtable}[]{@{}
  >{\raggedright\arraybackslash}p{(\linewidth - 2\tabcolsep) * \real{0.5000}}
  >{\raggedright\arraybackslash}p{(\linewidth - 2\tabcolsep) * \real{0.5000}}@{}}
\toprule\noalign{}
\begin{minipage}[b]{\linewidth}\raggedright
File
\end{minipage} & \begin{minipage}[b]{\linewidth}\raggedright
What it is
\end{minipage} \\
\midrule\noalign{}
\endhead
\bottomrule\noalign{}
\endlastfoot
\texttt{00-\allowbreak{}brief.\allowbreak{}md} & the brief template:
target, audience, non-negotiables, owner decisions, budget and stop \\
\texttt{01-\allowbreak{}sources-\allowbreak{}of-\allowbreak{}truth.\allowbreak{}md}
& the table of sources per content type \\
\texttt{02-\allowbreak{}canon.\allowbreak{}json} & the canon as data:
four sections, the lists the validator reads (empty), the recorded facts
(empty) \\
\texttt{03-\allowbreak{}artifacts.\allowbreak{}json} & one placeholder
artifact record \\
\texttt{04-\allowbreak{}roles/\allowbreak{}<role>.\allowbreak{}md} & one
card per role, seven bold fields with placeholders, a link to the
canon \\
\texttt{05-\allowbreak{}contracts/\allowbreak{}02-\allowbreak{}canon.\allowbreak{}schema.\allowbreak{}json}
& the contract of the JSON canon (written by the scaffold) \\
\texttt{05-\allowbreak{}contracts/\allowbreak{}notes.\allowbreak{}schema.\allowbreak{}json},
\texttt{notes-\allowbreak{}status.\allowbreak{}schema.\allowbreak{}json},
\texttt{round-\allowbreak{}record.\allowbreak{}schema.\allowbreak{}json},
\texttt{domain-\allowbreak{}result.\allowbreak{}schema.\allowbreak{}json}
& copied from the material \\
\texttt{06-\allowbreak{}gates.\allowbreak{}md} & an empty gate table and
an empty ``Size ranges'' table \\
\texttt{07-\allowbreak{}rubric.\allowbreak{}md},
\texttt{07-\allowbreak{}rubric.\allowbreak{}sha256} & the rubric
template and the pin placeholder \texttt{not-\allowbreak{}frozen} \\
\texttt{08-\allowbreak{}validator/\allowbreak{}validator.\allowbreak{}py},
\texttt{sensitivity.\allowbreak{}py},
\texttt{perturbation\_\allowbreak{}set.\allowbreak{}py} & copied from
the material \\
\texttt{08-\allowbreak{}validator/\allowbreak{}perturbations.\allowbreak{}json}
& three generic perturbations (an avoid-word, an absolute path, a
deleted card field) \\
\texttt{layout.\allowbreak{}json} & the layout of the cell (chapter
12) \\
\texttt{lessons.\allowbreak{}md},
\texttt{open-\allowbreak{}notes.\allowbreak{}json},
\texttt{priors.\allowbreak{}json} & empty \\
\texttt{state/\allowbreak{}notes-\allowbreak{}status.\allowbreak{}jsonl},
\texttt{state/\allowbreak{}round-\allowbreak{}record.\allowbreak{}jsonl}
& empty \\
\end{longtable}
}

\section{The first run is red, by
design}\label{the-first-run-is-red-by-design}

Verified on a cell scaffolded with three roles at 0.6.0: the self-test
passes (185 checks), and the first validator run

\begin{Shaded}
\begin{Highlighting}[]
\ExtensionTok{python3}\NormalTok{ 08{-}validator/validator.py }\AttributeTok{{-}{-}folder}\NormalTok{ . }\AttributeTok{{-}{-}layout}\NormalTok{ layout.json }\DataTypeTok{\textbackslash{}}
    \AttributeTok{{-}{-}rubric{-}hash}\NormalTok{ 07{-}rubric.sha256 }\AttributeTok{{-}{-}notes}\NormalTok{ open{-}notes.json }\DataTypeTok{\textbackslash{}}
    \AttributeTok{{-}{-}allow{-}pending}\NormalTok{ p{-}anonymity}
\end{Highlighting}
\end{Shaded}

exits 1 with 3 blocking and 8 high notes: ``no documents were measured;
a production with no documents is not green'' (the placeholder artifact
matches no file), ``declared artifact missing:
\texttt{<relative path>}'', ``no sensitivity report at
\texttt{08-\allowbreak{}validator/\allowbreak{}sensitivity.\allowbreak{}md}'',
and one high note per Markdown document of the cell that ``belongs to no
artifact and is not declared outside the artifact table''. This is the
contract: a newly scaffolded cell fails until it has real artifacts.

\section{The shortest path to green,
verified}\label{the-shortest-path-to-green-verified}

The following edits, made on a cell scaffolded with
\texttt{-\allowbreak{}-\allowbreak{}roles writer reviewer}, took it from
that red run to exit 0 with no notes in one pass (verified on 0.6.0; the
full chapter 18 is the real procedure, this is the smallest cell the
instrument accepts):

\begin{enumerate}
\def\labelenumi{\arabic{enumi}.}
\tightlist
\item
  \textbf{\texttt{00-\allowbreak{}brief.\allowbreak{}md}}: a link to the
  canon
  (\texttt{Canon:\allowbreak{} [02-\allowbreak{}canon.\allowbreak{}json](02-\allowbreak{}canon.\allowbreak{}json)}),
  a \texttt{Languages:\allowbreak{} en} line, two non-negotiables with
  ids (\texttt{nn-\allowbreak{}source}, \texttt{nn-\allowbreak{}size}),
  one decided owner decision, a budget and a stop.
\item
  \textbf{\texttt{02-\allowbreak{}canon.\allowbreak{}json}}:
  \texttt{terms to avoid} (\texttt{pipeline}, \texttt{checkpoint}), the
  seven \texttt{role card fields}, the three \texttt{origin labels}, and
  \texttt{documents outside the artifact list} set to
  \texttt{01-\allowbreak{}sources-\allowbreak{}of-\allowbreak{}truth.\allowbreak{}md},
  \texttt{07-\allowbreak{}rubric.\allowbreak{}md},
  \texttt{lessons.\allowbreak{}md},
  \texttt{08-\allowbreak{}validator/\allowbreak{}*.\allowbreak{}md}.
\item
  \textbf{\texttt{03-\allowbreak{}artifacts.\allowbreak{}json}}: three
  records --- \texttt{brief}
  (\texttt{00-\allowbreak{}brief.\allowbreak{}md}, gates
  \texttt{p-\allowbreak{}terms},
  \texttt{p-\allowbreak{}canon-\allowbreak{}link},
  \texttt{p-\allowbreak{}anonymity}, \texttt{p-\allowbreak{}form}),
  \texttt{cards}
  (\texttt{04-\allowbreak{}roles/\allowbreak{}*.\allowbreak{}md}, those
  four plus \texttt{p-\allowbreak{}card}) and \texttt{gates}
  (\texttt{06-\allowbreak{}gates.\allowbreak{}md}, gates
  \texttt{p-\allowbreak{}limit}, \texttt{p-\allowbreak{}terms},
  \texttt{p-\allowbreak{}canon-\allowbreak{}link},
  \texttt{p-\allowbreak{}form}).
\item
  \textbf{\texttt{04-\allowbreak{}roles/\allowbreak{}writer.\allowbreak{}md},
  \texttt{reviewer.\allowbreak{}md}}: the seven fields filled with one
  line each, the executor named by family and effort in the degree
  field.
\item
  \textbf{\texttt{06-\allowbreak{}gates.\allowbreak{}md}}: a link to the
  canon and one row per gate assigned above, each naming the same
  artifacts, with a magnitude, a limit and a justification of at least
  six words starting with an origin label; the ``Size ranges'' table
  left empty.
\item
  \textbf{\texttt{08-\allowbreak{}validator/\allowbreak{}perturbations.\allowbreak{}json}}:
  the three generic perturbations kept, plus one each for
  \texttt{p-\allowbreak{}form} (an HTML tag in the brief),
  \texttt{p-\allowbreak{}canon-\allowbreak{}link} (a broken link in a
  card) and \texttt{p-\allowbreak{}limit} (a justification deleted in
  the gates file), so that every assigned gate has one.
\end{enumerate}

Then the four commands of the cell (chapter 19): the self-test (185
checks); the freeze, which printed the pin and replaced
\texttt{not-\allowbreak{}frozen}; the sensitivity test, whose report
read ``Caught 6 of 6. Noisy: 0 of 6. Pairs covered: 6 of 22.''; and the
validator, which wrote a report of 23 rows with
\texttt{p-\allowbreak{}sensitivity} at ``6 of 6 gates covered, 6 of 22
pairs (27 \%), 0 missed of 6, 0 noisy'', three
\texttt{p-\allowbreak{}anonymity} rows pending and accepted (the
identifier half, with no forbidden list),
\texttt{m-\allowbreak{}determinism} passed, and exit 0. Had the three
generic perturbations been left alone,
\texttt{p-\allowbreak{}sensitivity} would have been pending for the
three uncovered gates and the run would have exited 4.

\chapter{Setting up a cell}\label{setting-up-a-cell}

This chapter condenses the fifteen steps of
\texttt{02-\allowbreak{}reflexive-\allowbreak{}application/\allowbreak{}10-\allowbreak{}generic-\allowbreak{}recipe.\allowbreak{}md}
and the procedure of the skill into one how-to. Each step says what it
creates and when it is done; the recipe adds, for each, who performs it,
what it reads and the trap the step exists to avoid.

\begin{figure}
\centering
\pandocbounded{\includegraphics[keepaspectratio,alt={The fifteen steps of the generic recipe and the file each one creates.}]{diagrams/recipe-steps.pdf}}
\caption{The fifteen steps of the generic recipe and the file each one
creates.}\label{fig-recipe-steps}
\end{figure}

\section{Step 0 --- the folder}\label{step-0-the-folder}

Create the tree (the scaffold does), decide where everything that stays
out of version control goes and record it in the brief, and decide how
machine-dependent paths enter: always as a command-line variable, never
written into an artifact. Done when the tree exists, the ledgers exist
and are empty, and \texttt{priors.\allowbreak{}json} parses. Trap:
starting to write the brief in a conversation. What is not in a file has
no hash, and what has no hash makes nobody stale.

\section{Step 1 --- the brief}\label{step-1-the-brief}

\texttt{00-\allowbreak{}brief.\allowbreak{}md}, with fixed sections:
target; for whom; non-negotiables, each with an id in a code span
(\texttt{nn-\allowbreak{}<name>}), written as sentences a gate could
measure (``mass ≤ X kg'', ``no fact without a source'', not ``safe'',
``good'', ``fast''); owner decisions as a table with a status (and,
since v0.6.0, as lines of
\texttt{state/\allowbreak{}owner-\allowbreak{}decisions.\allowbreak{}jsonl});
available sources of truth; budget (maximum rounds per artifact, review
rounds); a provisional stopping criterion; what is not an objective;
assumptions labeled \texttt{[hypothesis]}. The brief also declares the
languages of the documents on a line \texttt{Languages:\allowbreak{} en}
(two-letter codes), which the anonymity gate reads. Done when every
section is filled or says ``owner decision pending'', and the owner
recorded the approval.

\section{Step 2 --- the sources of
truth}\label{step-2-the-sources-of-truth}

\texttt{01-\allowbreak{}sources-\allowbreak{}of-\allowbreak{}truth.\allowbreak{}md}:
a table with content type, the source that errs least, who consults it,
and what is forbidden as a source. The writer is never a source of fact;
a generative tool is never a source of exact text; a number never comes
from the paragraph that cites it. Done when every content type of the
brief has exactly one source.

\section{Step 3 --- the canon}\label{step-3-the-canon}

\texttt{02-\allowbreak{}canon.\allowbreak{}json}, with its schema in
\texttt{05-\allowbreak{}contracts/\allowbreak{}}. Sections with owner,
version and state; canonical numbers (spelling by medium, unit,
rounding, nature, source); accepted variants; names and terms; rules.
Fill the lists the validator reads (a list left empty leaves its gate
pending): \texttt{terms to avoid}, \texttt{role card fields},
\texttt{origin labels} (with their brackets), and the others the
scaffold writes empty. List under
\texttt{documents outside the artifact list} the Markdown documents of
the cell that belong to no artifact
(\texttt{01-\allowbreak{}sources-\allowbreak{}of-\allowbreak{}truth.\allowbreak{}md},
\texttt{07-\allowbreak{}rubric.\allowbreak{}md},
\texttt{lessons.\allowbreak{}md},
\texttt{08-\allowbreak{}validator/\allowbreak{}*.\allowbreak{}md}), or
\texttt{p-\allowbreak{}declared} reports them. Every number the cell
states about a registered prior gets a row of \texttt{recorded facts}:

\begin{Shaded}
\begin{Highlighting}[]
\NormalTok{\{ "number": "18", "context": "roles with a card [prior:\textless{}id\textgreater{}]", "section": "1.1" \}}
\end{Highlighting}
\end{Shaded}

Done when every section has an owner and a version, every number of the
brief is in the canon, and the schema validates.

\section{Step 4 --- the artifacts and the
graph}\label{step-4-the-artifacts-and-the-graph}

\texttt{03-\allowbreak{}artifacts.\allowbreak{}json}: one record per
artifact (chapter 6). Start from the target and walk backward; give each
artifact a single owner; turn every owner decision into an artifact with
resource \texttt{human}; mark as \texttt{proposed} an artifact that has
no file yet; compute the topological order. The \texttt{gates} field of
each record is the one source of the gate map. Done when the graph is
acyclic, every artifact has one owner, at least one gate and (except the
target) at least one consumer. Trap: a finding that lands on the nearest
artifact because the right target is not in the graph.

\section{Step 5 --- the role cards}\label{step-5-the-role-cards}

\texttt{04-\allowbreak{}roles/\allowbreak{}<role>.\allowbreak{}md}, one
per artifact owner, with the seven fields as bold labels
(\texttt{**Decides:\allowbreak{}**}), plus ``Questions it always asks''
and the base prompt the orchestrator hands to the agent. In ``Does not
decide'', write the neighboring decisions the role would be tempted to
make and whose they are. In the degree field, name the executor: a
person, code, or a kind of language model by family and effort; a role
that audits another gets an executor different from the one that
produced the artifact. Done when every card has the seven fields, no
decision appears in ``Decides'' of two cards, and no auditing role
shares an executor with the role it audits.

\section{Step 6 --- the contracts}\label{step-6-the-contracts}

\texttt{05-\allowbreak{}contracts/\allowbreak{}<artifact>.\allowbreak{}schema.\allowbreak{}json}
for each artifact that passes between roles, rejecting unknown fields
and the domain's most likely ambiguity, with one valid and one invalid
example; the invalid one has to fail. Done when validation runs in
seconds, before any expensive work. Trap: the ambiguous position key.

\section{Step 7 --- the gates}\label{step-7-the-gates}

\texttt{06-\allowbreak{}gates.\allowbreak{}md}: one row per gate of the
artifact list, naming the same artifacts, with magnitude, limit and a
justification that starts with an origin label and has at least six
words; a ``Size ranges'' table with the word range of each document,
declared before writing. Give every signal gate a magnitude limit on
both sides; separate ``rejects'' from ``tool failure''; never let a
limit inherit a prior's number. A gate of the domain is a row
\texttt{d-\allowbreak{}<name>} (chapter 21). The ids have to be ones the
validator implements, or \texttt{d-\allowbreak{}*}: any other id exits
2. Done when every limit has a justification with an origin, and every
code gate rejected its known perturbation.

\section{Step 8 --- the rubric, frozen}\label{step-8-the-rubric-frozen}

\texttt{07-\allowbreak{}rubric.\allowbreak{}md}: six to ten dimensions
with anchors for 1, 3 and 5; the severity scale; the mechanical verdict
rule; the report format. Then:

\begin{Shaded}
\begin{Highlighting}[]
\ExtensionTok{python3}\NormalTok{ 08{-}validator/validator.py }\AttributeTok{{-}{-}folder}\NormalTok{ . }\AttributeTok{{-}{-}layout}\NormalTok{ layout.json }\DataTypeTok{\textbackslash{}}
    \AttributeTok{{-}{-}freeze{-}rubric}\NormalTok{ 07{-}rubric.sha256}
\end{Highlighting}
\end{Shaded}

The scaffold's pin holds the placeholder
\texttt{not-\allowbreak{}frozen}, which the first freeze replaces; a
real pin is refused, and a deliberate replacement goes through
\texttt{-\allowbreak{}-\allowbreak{}re-\allowbreak{}freeze-\allowbreak{}rubric}.
Done when the hash is recorded before the first review and the round
record will name who wrote the rubric. Written by the owner with whoever
will review; never by whoever will produce, alone.

\section{Step 9 --- the validator}\label{step-9-the-validator}

The scaffold copied \texttt{validator.\allowbreak{}py},
\texttt{sensitivity.\allowbreak{}py} and
\texttt{perturbation\_\allowbreak{}set.\allowbreak{}py} into
\texttt{08-\allowbreak{}validator/\allowbreak{}}. Adapt the three
generic perturbations of \texttt{perturbations.\allowbreak{}json} to the
cell's text and add one per gate the cell assigns. Run the sensitivity
test, then the validator (chapter 19). Done when two runs give an
identical report, every gate has its perturbation caught, and the
validator refuses a canon of another schema version.

\section{Step 10 --- the assembly
engine}\label{step-10-the-assembly-engine}

\texttt{09-\allowbreak{}engine/\allowbreak{}}: the graph core (state,
hash per artifact, plan, execute), the note router, the assembly of the
target and the tests. The content hash of an artifact is the sha256 over
the canonical form of the file (UTF-8, the whitespace at the end of each
line normalized, a single final newline). The plan marks each artifact
\texttt{current}, \texttt{stale}, \texttt{missing}, \texttt{returned} or
\texttt{escalated}; only what is stale runs, in topological order; the
router refuses an unknown target, deduplicates, and only high and
blocking return the artifact; the assembly validates every contract
before the expensive work and names the output with the version label.
Done when assembling twice with the same inputs gives the same product.
Trap: the namesake file from an old version, which every reader silently
picks.

\section{Step 11 --- the quick first version and its
review}\label{step-11-the-quick-first-version-and-its-review}

Produce every artifact, even bad, instead of polishing the first ones;
assemble; hand the reviewer only the files and the rubric; record in the
dispatch the output limit and the reason the generation ended; import
the review's notes into the ledger. Done when every artifact exists, the
report follows the rubric's format and the notes entered the ledger
without refusal. Trap: polishing v1.

\section{Step 12 --- the round loop}\label{step-12-the-round-loop}

Chapter 20.

\section{Step 13 --- the stop}\label{step-13-the-stop}

The owner declares the criterion before the last round, in terms of
rounds or verdict, never calendar; the review of the last version
becomes open notes in
\texttt{open-\allowbreak{}notes.\allowbreak{}json}, in the contract; a
violated non-negotiable is not an open note but an undelivered target.

\section{Step 14 --- the case method document and the
registration}\label{step-14-the-case-method-document-and-the-registration}

\texttt{method-\allowbreak{}TARGET.\allowbreak{}md}: the ideas with the
defect each prevents, the graph, the step by step with the real
commands, the pseudocode per role with its degree, the owner decisions,
the known false positives with a status, \texttt{[unverified]} on
whatever was not checked. Register the instance in
\texttt{priors.\allowbreak{}json} with id, status, the document's path
relative to its own repository root, its sha256, its locator and its
number format. Trap: a document that repeats the READMEs.

\section{Before handing over}\label{before-handing-over}

Run through
\texttt{skills/\allowbreak{}cooperatio-\allowbreak{}teamwork/\allowbreak{}references/\allowbreak{}review-\allowbreak{}checklist.\allowbreak{}md}.
At minimum: every document points to the canon; every gate has
magnitude, limit and a justified origin label; the graph is acyclic;
every owner is a role or the owner; every card has the seven fields; the
rubric hash is recorded; every non-negotiable of the domain is a
\texttt{d-\allowbreak{}*} gate whose result the cell's code writes; the
validator exits 0 and every pending it reports is one of the four it
accepts by name, with the reason; and the files do not claim the method
has been validated on this target --- the first review round does that.

\chapter{Running the instrument}\label{running-the-instrument}

\section{The commands of a cell}\label{the-commands-of-a-cell}

From inside the cell, after the files are filled (the skill, ``The
commands of a cell''):

\begin{Shaded}
\begin{Highlighting}[]
\ExtensionTok{python3}\NormalTok{ 08{-}validator/validator.py }\AttributeTok{{-}{-}folder}\NormalTok{ . }\AttributeTok{{-}{-}layout}\NormalTok{ layout.json }\AttributeTok{{-}{-}self{-}test}
\ExtensionTok{python3}\NormalTok{ 08{-}validator/validator.py }\AttributeTok{{-}{-}folder}\NormalTok{ . }\AttributeTok{{-}{-}layout}\NormalTok{ layout.json }\DataTypeTok{\textbackslash{}}
    \AttributeTok{{-}{-}freeze{-}rubric}\NormalTok{ 07{-}rubric.sha256}
\ExtensionTok{python3}\NormalTok{ 08{-}validator/sensitivity.py }\AttributeTok{{-}{-}folder}\NormalTok{ . }\AttributeTok{{-}{-}layout}\NormalTok{ layout.json }\DataTypeTok{\textbackslash{}}
    \AttributeTok{{-}{-}rubric{-}hash}\NormalTok{ 07{-}rubric.sha256 }\AttributeTok{{-}{-}notes}\NormalTok{ open{-}notes.json }\AttributeTok{{-}{-}work}\NormalTok{ SCRATCH\_FOLDER}
\ExtensionTok{python3}\NormalTok{ 08{-}validator/validator.py }\AttributeTok{{-}{-}folder}\NormalTok{ . }\AttributeTok{{-}{-}layout}\NormalTok{ layout.json }\DataTypeTok{\textbackslash{}}
    \AttributeTok{{-}{-}rubric{-}hash}\NormalTok{ 07{-}rubric.sha256 }\AttributeTok{{-}{-}notes}\NormalTok{ open{-}notes.json }\DataTypeTok{\textbackslash{}}
    \AttributeTok{{-}{-}allow{-}pending}\NormalTok{ p{-}anonymity }\AttributeTok{{-}{-}determinism} \AttributeTok{{-}{-}output}\NormalTok{ 08{-}validator/report.md}
\ExtensionTok{python3} \AttributeTok{{-}B} \AttributeTok{{-}m}\NormalTok{ pytest }\AttributeTok{{-}p}\NormalTok{ no:cacheprovider tests/}
\end{Highlighting}
\end{Shaded}

The last line only if the cell keeps a test suite. Add
\texttt{-\allowbreak{}-\allowbreak{}case} for each registered prior the
cell cites, and \texttt{-\allowbreak{}-\allowbreak{}forbidden} with the
cell's list of forbidden identifiers kept outside the cell (then drop
\texttt{-\allowbreak{}-\allowbreak{}allow-\allowbreak{}pending}). Run
the sensitivity test again after every change to the validator, the
gates file, the canon, the artifact list's gate map or the perturbation
set, and before the validator: the validator reads the report and fails
\texttt{p-\allowbreak{}sensitivity} on provenance otherwise.

\section{The commands of the corpus}\label{the-commands-of-the-corpus}

From
\texttt{02-\allowbreak{}reflexive-\allowbreak{}application/\allowbreak{}}
(its README, ``How to run the validator''):

\begin{Shaded}
\begin{Highlighting}[]
\ExtensionTok{python3}\NormalTok{ 80{-}validator/sensitivity.py }\AttributeTok{{-}{-}folder}\NormalTok{ . }\DataTypeTok{\textbackslash{}}
    \AttributeTok{{-}{-}case}\NormalTok{ PRIOR\_000\_DOCUMENT\_PATH }\DataTypeTok{\textbackslash{}}
    \AttributeTok{{-}{-}case}\NormalTok{ ../instances/004{-}adoption{-}guide/metodo{-}de{-}producao.md }\DataTypeTok{\textbackslash{}}
    \AttributeTok{{-}{-}rubric{-}hash}\NormalTok{ 90{-}review/rubric.sha256 }\AttributeTok{{-}{-}notes}\NormalTok{ 90{-}review/notes.json }\DataTypeTok{\textbackslash{}}
    \AttributeTok{{-}{-}work}\NormalTok{ SCRATCH\_FOLDER }\AttributeTok{{-}{-}output}\NormalTok{ 80{-}validator/sensitivity.md}
\ExtensionTok{python3}\NormalTok{ 80{-}validator/validator.py }\AttributeTok{{-}{-}folder}\NormalTok{ . }\DataTypeTok{\textbackslash{}}
    \AttributeTok{{-}{-}case}\NormalTok{ PRIOR\_000\_DOCUMENT\_PATH }\DataTypeTok{\textbackslash{}}
    \AttributeTok{{-}{-}case}\NormalTok{ ../instances/004{-}adoption{-}guide/metodo{-}de{-}producao.md }\DataTypeTok{\textbackslash{}}
    \AttributeTok{{-}{-}rubric{-}hash}\NormalTok{ 90{-}review/rubric.sha256 }\AttributeTok{{-}{-}notes}\NormalTok{ 90{-}review/notes.json }\DataTypeTok{\textbackslash{}}
    \AttributeTok{{-}{-}allow{-}pending}\NormalTok{ p{-}anonymity }\AttributeTok{{-}{-}determinism} \AttributeTok{{-}{-}output}\NormalTok{ 80{-}validator/report.md}
\ExtensionTok{python3}\NormalTok{ 80{-}validator/validator.py }\AttributeTok{{-}{-}self{-}test}
\ExtensionTok{python3} \AttributeTok{{-}B} \AttributeTok{{-}m}\NormalTok{ pytest }\AttributeTok{{-}p}\NormalTok{ no:cacheprovider 80{-}validator/test\_validator.py }\DataTypeTok{\textbackslash{}}
\NormalTok{    ../skills/cooperatio{-}teamwork/scripts/test\_scaffold.py}
\end{Highlighting}
\end{Shaded}

The published report is made with exactly that command; at 0.6.0 it
exits 0 with no notes, 27 rows of \texttt{p-\allowbreak{}anonymity}
pending and accepted (the identifier half, with no forbidden list),
\texttt{m-\allowbreak{}determinism} passed, the self-test at 185 checks;
the sensitivity report reads 65 of 65 caught, 0 noisy, 50 of 182 pairs;
the test suites pass 280 tests.

\section{The clone run}\label{the-clone-run}

The method document of instance 000 stays with the owner, so a clone of
the material does not have it. The clone passes the document of 004 and
accepts by name the gates that need the owner's:

\begin{Shaded}
\begin{Highlighting}[]
\ExtensionTok{python3}\NormalTok{ 80{-}validator/validator.py }\AttributeTok{{-}{-}folder}\NormalTok{ . }\DataTypeTok{\textbackslash{}}
    \AttributeTok{{-}{-}case}\NormalTok{ ../instances/004{-}adoption{-}guide/metodo{-}de{-}producao.md }\DataTypeTok{\textbackslash{}}
    \AttributeTok{{-}{-}rubric{-}hash}\NormalTok{ 90{-}review/rubric.sha256 }\AttributeTok{{-}{-}notes}\NormalTok{ 90{-}review/notes.json }\DataTypeTok{\textbackslash{}}
    \AttributeTok{{-}{-}allow{-}pending}\NormalTok{ p{-}anonymity,p{-}prior,p{-}prior{-}pin }\AttributeTok{{-}{-}determinism}
\end{Highlighting}
\end{Shaded}

Expected: exit 0, no notes, the report on the standard output, the three
gates listed as pending accepted by the caller, the pin and the numbers
of 004 measured (its edited document, or an invented number about it, is
a blocking note, exit 1), and the rows of 000 saying ``owner-held
evidence, not measured here''. Without any
\texttt{-\allowbreak{}-\allowbreak{}case} the same run still exits 0,
but measures neither pin.

A clone that edits the instrument cannot regenerate the sensitivity
report with the owner's document; it regenerates one with a synthetic
prior:

\begin{Shaded}
\begin{Highlighting}[]
\ExtensionTok{python3}\NormalTok{ 80{-}validator/sensitivity.py }\AttributeTok{{-}{-}folder}\NormalTok{ . }\DataTypeTok{\textbackslash{}}
    \AttributeTok{{-}{-}case}\NormalTok{ ../instances/004{-}adoption{-}guide/metodo{-}de{-}producao.md }\DataTypeTok{\textbackslash{}}
    \AttributeTok{{-}{-}synthetic{-}prior}\NormalTok{ 000{-}av{-}production }\DataTypeTok{\textbackslash{}}
    \AttributeTok{{-}{-}rubric{-}hash}\NormalTok{ 90{-}review/rubric.sha256 }\AttributeTok{{-}{-}notes}\NormalTok{ 90{-}review/notes.json }\DataTypeTok{\textbackslash{}}
    \AttributeTok{{-}{-}work}\NormalTok{ SCRATCH\_FOLDER }\AttributeTok{{-}{-}output}\NormalTok{ 80{-}validator/sensitivity.md}
\end{Highlighting}
\end{Shaded}

and then adds \texttt{p-\allowbreak{}sensitivity-\allowbreak{}synthetic}
to \texttt{-\allowbreak{}-\allowbreak{}allow-\allowbreak{}pending} of
the clone run. The report stays pending until the owner regenerates it
with the document (chapter 14).

\section{Reading a report}\label{reading-a-report}

Read the pending rows before the exit code. A row is \texttt{pass},
\texttt{fail} or \texttt{pending}; the lines after the table say which
pending gates the caller accepted, which it did not (the run exits 4 for
those unless a note exits 1), and which accepted gates had nothing
pending (``Accepted but not pending'', so the caller sees that the
acceptance covered nothing). The notes table gives, per note, the target
artifact, the severity, the text (prefixed by the gate id) and the
evidence (the document). With
\texttt{-\allowbreak{}-\allowbreak{}output}, one summary line goes to
the error stream.

\section{Troubleshooting}\label{troubleshooting}

{\def\LTcaptype{none} % do not increment counter
\begin{longtable}[]{@{}
  >{\raggedright\arraybackslash}p{(\linewidth - 4\tabcolsep) * \real{0.3333}}
  >{\raggedright\arraybackslash}p{(\linewidth - 4\tabcolsep) * \real{0.3333}}
  >{\raggedright\arraybackslash}p{(\linewidth - 4\tabcolsep) * \real{0.3333}}@{}}
\toprule\noalign{}
\begin{minipage}[b]{\linewidth}\raggedright
Message or symptom
\end{minipage} & \begin{minipage}[b]{\linewidth}\raggedright
Cause
\end{minipage} & \begin{minipage}[b]{\linewidth}\raggedright
What to do
\end{minipage} \\
\midrule\noalign{}
\endhead
\bottomrule\noalign{}
\endlastfoot
\texttt{no documents were measured; a production with no documents is not green}
(blocking, on \texttt{artifacts}) & the artifact list has no record
whose paths match a file, or no gates & fill
\texttt{03-\allowbreak{}artifacts.\allowbreak{}json} with real records
and gates \\
\texttt{p-\allowbreak{}declared:\allowbreak{} declared artifact missing:\allowbreak{} <path>}
(blocking) & a literal path of an artifact matches no file, or a glob
matches nothing & create the file or fix the path \\
\texttt{p-\allowbreak{}declared:\allowbreak{} document X belongs to no artifact and is not declared outside the artifact table}
(high) & a Markdown document of the folder is neither an artifact nor
listed in the canon & add it to an artifact's paths, or list it under
\texttt{documents outside the artifact list} with the reason \\
\texttt{p-\allowbreak{}sensitivity:\allowbreak{} no sensitivity report at .\allowbreak{}.\allowbreak{}.\allowbreak{}}
(blocking) & \texttt{sensitivity.\allowbreak{}py} has not run, or the
layout key \texttt{artifacts.\allowbreak{}sensitivity} points elsewhere
& run the sensitivity test \\
\texttt{p-\allowbreak{}sensitivity} fails on provenance (validator,
gates, canon or map sha) & the instrument or its inputs changed after
the report & regenerate the report, then the validator report \\
\texttt{p-\allowbreak{}sensitivity} pending: an assigned gate without a
perturbation, or fewer than 25 \% of pairs & the perturbation set does
not cover the cell's gates & add one perturbation per gate the artifact
list assigns \\
exit 4 with
\texttt{Pending not accepted:\allowbreak{} .\allowbreak{}.\allowbreak{}.\allowbreak{}}
& a gate was not measured and the caller did not accept it & supply the
input (\texttt{-\allowbreak{}-\allowbreak{}notes},
\texttt{-\allowbreak{}-\allowbreak{}rubric-\allowbreak{}hash},
\texttt{-\allowbreak{}-\allowbreak{}case}, the canon's list, the domain
result), or accept it by name when it is one of the four \\
\texttt{-\allowbreak{}-\allowbreak{}allow-\allowbreak{}pending accepts only p-\allowbreak{}anonymity, p-\allowbreak{}prior, p-\allowbreak{}prior-\allowbreak{}pin, p-\allowbreak{}sensitivity-\allowbreak{}synthetic .\allowbreak{}.\allowbreak{}.\allowbreak{}}
(exit 2) & the caller named a gate that has to be measured & measure
it \\
\texttt{-\allowbreak{}-\allowbreak{}allow-\allowbreak{}pending names gates this validator does not implement}
(exit 2) & a typo in the gate name & fix the name \\
\texttt{the rubric is already frozen at .\allowbreak{}.\allowbreak{}.\allowbreak{}}
(exit 2) &
\texttt{-\allowbreak{}-\allowbreak{}freeze-\allowbreak{}rubric} on an
existing pin & leave the pin, or
\texttt{-\allowbreak{}-\allowbreak{}re-\allowbreak{}freeze-\allowbreak{}rubric}
deliberately with the reason in the changelog \\
\texttt{the gate table has more than one row for .\allowbreak{}.\allowbreak{}.\allowbreak{}}
(exit 2) & two rows for one gate & one row per gate \\
a gate id the validator does not implement (exit 2) & a placeholder row,
a typo, or a cell-specific gate not named \texttt{d-\allowbreak{}*} &
use a builtin id, or a \texttt{d-\allowbreak{}*} row with a result
file \\
\texttt{layout:\allowbreak{} .\allowbreak{}.\allowbreak{}.\allowbreak{}}
or a layout key refused (exit 2) & a key the validator does not read (a
typo such as \texttt{sensitivty}) & use the keys of chapter 12 \\
\texttt{more than one parser} (exit 2) & a cell from before 0.4.0 whose
layout lacks \texttt{p-\allowbreak{}efficacy-\allowbreak{}mark} in its
gate map, so the merged map names two parsers & add the gate to the
layout's map \\
\texttt{p-\allowbreak{}prior}: a number of a marked unit unverified
(blocking) & the number has no row in the recorded facts, or is absent
from the sections the row cites, or the prior's document was not passed
& add the row and check the number; pass
\texttt{-\allowbreak{}-\allowbreak{}case} \\
\texttt{p-\allowbreak{}prior}: malformed mark (blocking) & a
\texttt{[prior:\allowbreak{}.\allowbreak{}.\allowbreak{}.\allowbreak{}]}
whose id is not lowercase letters, digits and hyphens & fix the id \\
\texttt{p-\allowbreak{}prior-\allowbreak{}pin} blocking on the registry
& a document passed is the pin of no executed instance, or a distributed
prior's pin matches no document & pass the right document; check the
registry's sha256 \\
\texttt{p-\allowbreak{}contract-\allowbreak{}shape}: a broken chain & an
earlier chained line was edited, removed or reordered & restore the
line; a state change is a new line \\
\texttt{p-\allowbreak{}notes}: an escaped regression closed with neither
a gate nor a reason (high) & the ledger line has
\texttt{escaped:\allowbreak{} true} and no \texttt{gate} of the gate
table, no \texttt{no\_\allowbreak{}gate} of 8 distinct words and no
decided \texttt{decision} & name the gate, or write the reason \\
\texttt{p-\allowbreak{}anonymity-\allowbreak{}language} pending on the
brief & the brief declares a language with no built-in months and no row
in the canon's ``Month names'' & add the row with the twelve names \\
\texttt{Python 3.\allowbreak{}9 or later is required} (exit 2) & an
older interpreter & use 3.9 or later \\
\end{longtable}
}

\chapter{Running a round}\label{running-a-round}

\begin{figure}
\centering
\pandocbounded{\includegraphics[keepaspectratio,alt={A round of a cell, from the first version to the stop.}]{diagrams/round-loop.pdf}}
\caption{A round of a cell, from the first version to the
stop.}\label{fig-round-loop}
\end{figure}

\section{Before the review}\label{before-the-review}

The review only starts after the validator passes on every document and
the sensitivity report is current: reviewing a text that the validator
fails would spend a review round on defects the code already finds.

\section{The dispatch}\label{the-dispatch}

The reviewer receives the files and the rubric, never the reasoning of
whoever produced them. The dispatch records the effective output limit
of the executor and the reason its generation ended. A report that ended
for length or lacks a section of the format is a tool failure and
produces no verdict; the review runs again. The reviewer's relation to
the producer is recorded in the round record
(\texttt{reviewer\_\allowbreak{}relation}), checked on what each alias
or route maps to, not on its name: a second session of the producing
model is \texttt{same family, separate session} or
\texttt{same model, separate route}, never ``independent''.

\section{Importing and classifying}\label{importing-and-classifying}

Import the notes into the ledger (the router refuses an unknown target
and deduplicates); only high and blocking return the artifact. Classify
each finding:

\begin{itemize}
\tightlist
\item
  \textbf{product defect} --- a note to the artifact where it
  originated;
\item
  \textbf{regression that escaped every gate} --- first create the gate
  (which has to reject the current version) with its perturbation, then
  fix; the state ledger records the note as \texttt{escaped}, and it
  closes with the \texttt{gate} it became or a
  \texttt{no\_\allowbreak{}gate} reason;
\item
  \textbf{false positive} --- fix the instrument's rule at the root,
  with a test that reproduces the case, and redo the previous round to
  prove that no real finding disappeared; never loosen the limit.
\end{itemize}

Each owner answers note by note: done (where), partial (what is
missing), contested (reason) or out of scope (for whom). A validator
note the writer disputes goes to the verifier, who decides whether it is
a false positive.

\section{Closing the round}\label{closing-the-round}

The engine recompiles what is stale; the gates run; the validator runs;
the reviewer reviews with the same rubric and makes the item-by-item
table of the previous round. The round's entry goes to
\texttt{lessons.\allowbreak{}md} (what it taught and what became code,
contract or gate), and a line goes to
\texttt{state/\allowbreak{}round-\allowbreak{}record.\allowbreak{}jsonl}:

\begin{Shaded}
\begin{Highlighting}[]
\NormalTok{\{"method\_version": "0.6.0", "round": "v2", "kind": "production",}
\NormalTok{ "models": [\{"role": "writer", "model": "family A, high effort"\},}
\NormalTok{            \{"role": "reviewer", "model": "family B, high effort"\}],}
\NormalTok{ "reviewer\_relation": "another family",}
\NormalTok{ "rubric\_applied": true, "rubric\_hash": "\textless{}the frozen sha256\textgreater{}", "rubric\_author": "owner",}
\NormalTok{ "family\_map\_commitment": "\textless{}64 hex characters\textgreater{}",}
\NormalTok{ "summary": "...", "prev": "\textless{}sha256 of every line above\textgreater{}"\}}
\end{Highlighting}
\end{Shaded}

A production line may carry a \texttt{verdict}; a maintenance line may
not. The optional cost fields (per role: \texttt{tokens} with
\texttt{input}, \texttt{output}, \texttt{reasoning}, \texttt{cached};
\texttt{calls}; \texttt{finish}) and the round's
\texttt{wall\_\allowbreak{}seconds}, \texttt{notes\_\allowbreak{}count},
\texttt{crew} and \texttt{rounds\_\allowbreak{}to\_\allowbreak{}stop}
are what the principle asks a round to measure, so that costs become
comparable.

\section{The stop}\label{the-stop}

Done when the criterion is written in terms of rounds or verdict, the
open notes validate against the contract with their state in the state
ledger, no non-negotiable is among the open notes (an open high or
blocking note that affects one is blocking at the stop unless a decided
owner decision covers it), the lessons have the round's entry, and the
instance is registered and pinned in \texttt{priors.\allowbreak{}json}.

\chapter{Domain gates in a cell}\label{domain-gates-in-a-cell}

The distributed validator measures documents. A target whose
non-negotiables are domain measurements --- identity of results, a
speedup, a mass --- needs the cell's own code to write them as results
the validator reads, or its exit 0 says nothing about them.

\section{Declaring the gate}\label{declaring-the-gate}

A row \texttt{d-\allowbreak{}<name>} in the cell's gate table, with its
magnitude, limit and justification like any other; an optional Severity
column sets the severity of a \texttt{fail} (high by default; a severity
below high counts only for a gate a decided owner decision covers, so a
failed identity of results cannot become a low note by the cell's own
column). The artifact list assigns the gate to the artifacts it
measures, and the result's inputs have to include every file of those
artifacts.

\section{Writing the result}\label{writing-the-result}

The cell's code measures the gate and writes a result file in the
contract
\texttt{05-\allowbreak{}contracts/\allowbreak{}domain-\allowbreak{}result.\allowbreak{}schema.\allowbreak{}json}:

\begin{Shaded}
\begin{Highlighting}[]
\NormalTok{\{}
\NormalTok{  "contract\_version": "1",}
\NormalTok{  "gate": "d{-}identity",}
\NormalTok{  "magnitude": "0 divergent of 38 cases",}
\NormalTok{  "limit": "0 divergent",}
\NormalTok{  "result": "pass",}
\NormalTok{  "inputs": [ \{ "path": "v2/candidate.bin", "sha256": "\textless{}sha256 of that file\textgreater{}" \},}
\NormalTok{              \{ "path": "oracle/reference.bin", "sha256": "\textless{}sha256 of that file\textgreater{}" \} ]}
\NormalTok{\}}
\end{Highlighting}
\end{Shaded}

\texttt{result} is \texttt{pass}, \texttt{fail} or
\texttt{tool-\allowbreak{}failure}. The sha256 of an input is that of
the file's text read as UTF-8 with LF line breaks (for a binary file, of
its bytes). The layout lists the result files under
\texttt{domain-\allowbreak{}results}.

\section{What the validator does with
it}\label{what-the-validator-does-with-it}

It recomputes the sha256 of every input; an input changed or missing
makes the result stale (blocking); a result that names no input, or
whose inputs do not include every file of the gate's artifacts, is
blocking; when the limit starts with a comparison and a number or is a
bare zero, the first number of the magnitude is compared and a
contradiction is blocking; \texttt{fail} is a note of the row's
severity; \texttt{tool-\allowbreak{}failure} or a result outside the
contract is exit 2; a declared \texttt{d-\allowbreak{}*} gate with no
result is pending, which no caller can accept. The validator never runs
the cell's code: a result whose inputs were recomputed by hand after the
candidate changed passes, and the material says so.

\section{The cell's own sensitivity}\label{the-cells-own-sensitivity}

The sensitivity test perturbs documents and leaves domain gates out of
its coverage; the sensitivity of a domain gate is the cell's own:
mutants of the implementation, a reference with a known defect, a seeded
regression that has to leave the run red. Experiment E5 is designed
around exactly this (a candidate that breaks identity has to leave the
validator at exit 1), and a failing test of the cell is a rejection of
its gate, with a note to the implementation where the defect was born.

\chapter{Maintaining the material}\label{maintaining-the-material}

For a maintainer of the material itself, or of a cell that carries a
copy of the instrument.

\section{The order of a change}\label{the-order-of-a-change}

Editing the validator, the ``Gates'' section of the gates file, the
canon or the gate map of the artifact table changes a sha the
sensitivity report records. The order that keeps the provenance
coherent, as the changelog's ``Round infra'' entries record it: the
scripts first; then the canon and the gate section the provenance pins;
then the sensitivity report (regenerated with the owner's document, or
with the synthetic prior in a clone); then the other documents and the
validator report; then a round-record line, chained, with the executors
by family and effort.

\section{Adding a gate or a
perturbation}\label{adding-a-gate-or-a-perturbation}

A new gate needs, in the same change: its implementation and its handler
in the family it belongs to; a row in the gate table; its assignment in
the artifact list; its perturbation in
\texttt{perturbation\_\allowbreak{}set.\allowbreak{}py} (or the cell's
\texttt{perturbations.\allowbreak{}json}); a test in the suite; its
declared semantics in
\texttt{06-\allowbreak{}gate-\allowbreak{}semantics.\allowbreak{}md}
where a cell of the table is not enough; and a regenerated sensitivity
report. The lesson of v0.6.0 is the rule: when a correction moves a
definition, the place it moves to gets the pin, the completeness check
and the perturbation in the same round, or the correction is declared
incomplete.

\section{Retiring a gate}\label{retiring-a-gate}

The row, the assignment and the perturbation leave together, and the
changelog records ``retired gate: \texttt{<id>}'', the reason and the
rounds it stayed silent (chapter 7).

\section{Re-freezing the rubric}\label{re-freezing-the-rubric}

Only when another review round opens: a new rubric is written by the
reviewing side, frozen through
\texttt{-\allowbreak{}-\allowbreak{}re-\allowbreak{}freeze-\allowbreak{}rubric}
with the reason in the changelog, and the round record names its author.

\section{Ledgers and frozen records}\label{ledgers-and-frozen-records}

A change of state is a new chained line, never an edit; the last line of
a ledger is the one the chain does not yet cover, so append before
relying on it. The records of
\texttt{90-\allowbreak{}review/\allowbreak{}} and the history of the
lessons above the measured heading are never edited; the append-only
rule of the lessons is checked by hand at each round (the old file has
to be a prefix of the new one), and the sweep for forbidden identifiers
is run by hand each round, recorded outside the material.

\section{Migration notes}\label{migration-notes}

Each entry of \texttt{CHANGELOG.\allowbreak{}md} carries a migration
note for a cell that moves to that version. The ones a cell from an
older validator meets, in order:

{\def\LTcaptype{none} % do not increment counter
\begin{longtable}[]{@{}
  >{\raggedright\arraybackslash}p{(\linewidth - 2\tabcolsep) * \real{0.5000}}
  >{\raggedright\arraybackslash}p{(\linewidth - 2\tabcolsep) * \real{0.5000}}@{}}
\toprule\noalign{}
\begin{minipage}[b]{\linewidth}\raggedright
To version
\end{minipage} & \begin{minipage}[b]{\linewidth}\raggedright
A cell has to
\end{minipage} \\
\midrule\noalign{}
\endhead
\bottomrule\noalign{}
\endlastfoot
0.3.0 & re-run \texttt{sensitivity.\allowbreak{}py}, since reports
without the provenance header fail \texttt{p-\allowbreak{}sensitivity}
by design; name its rubric in the layout
(\texttt{artifacts.\allowbreak{}rubric}) \\
0.4.0 & add \texttt{p-\allowbreak{}efficacy-\allowbreak{}mark} to the
gate map of its layout (else exit 2, ``more than one parser''); give its
executed priors a \texttt{number\_\allowbreak{}format}; give every
number it states about a prior a row of its JSON canon's recorded facts;
register its own entry in its priors registry; make the card links
resolve; regenerate the sensitivity report \\
0.4.1 & cite an instance inside a marked unit by its full id; give every
executed entry a \texttt{locator} and a
\texttt{number\_\allowbreak{}format}, and keep the four fields null on a
designed one; cite only executed experiments; name executors by family
and effort, model identifiers going to the forbidden list \\
0.4.2 & replace a copy of the validator from before 0.4.2 (it wrote a
bytecode cache); run pytest with
\texttt{-\allowbreak{}B -\allowbreak{}p no:\allowbreak{}cacheprovider} \\
0.5.0 & give its artifact records a \texttt{gates} list and keep the
gate table in step; fill the canon's lists or accept their pending by
name; remove placeholders; write cards in bold fields; add
\texttt{contract\_\allowbreak{}version} to its contracts; give each note
a line in the state ledger; copy \texttt{sensitivity.\allowbreak{}py},
write \texttt{perturbations.\allowbreak{}json}, name it in the layout
and regenerate the report \\
0.6.0 & list its Markdown documents outside the artifact list in the
canon; give every literal path a file; write its
\texttt{d-\allowbreak{}*} gates as result files under
\texttt{domain-\allowbreak{}results}; use
\texttt{[lessons:\allowbreak{}<title>]} and a provenance table for each
cited experiment; freeze its rubric and pass
\texttt{-\allowbreak{}-\allowbreak{}notes} (their pending can no longer
be accepted); regenerate the report (map, specificity, pairs); copy
\texttt{perturbation\_\allowbreak{}set.\allowbreak{}py}; give its
non-negotiables ids and keep the owner's decisions in
\texttt{state/\allowbreak{}owner-\allowbreak{}decisions.\allowbreak{}jsonl};
schedule, dispute or close the open notes older than three rounds;
declare \texttt{"version":\allowbreak{} "v1"} in its JSON canon \\
\end{longtable}
}

\section{The tests}\label{the-tests}

\begin{Shaded}
\begin{Highlighting}[]
\ExtensionTok{python3}\NormalTok{ 02{-}reflexive{-}application/80{-}validator/validator.py }\AttributeTok{{-}{-}self{-}test}
\ExtensionTok{python3} \AttributeTok{{-}B} \AttributeTok{{-}m}\NormalTok{ pytest }\AttributeTok{{-}p}\NormalTok{ no:cacheprovider }\DataTypeTok{\textbackslash{}}
\NormalTok{    02{-}reflexive{-}application/80{-}validator/test\_validator.py }\DataTypeTok{\textbackslash{}}
\NormalTok{    skills/cooperatio{-}teamwork/scripts/test\_scaffold.py}
\end{Highlighting}
\end{Shaded}

The suite runs the self-test, one passing and one failing case per rule,
the regression tests of the probes the panels found (a repointed file, a
mark outside its grammar, a number between backticks, a deleted report
accepted by name, a domain regression the validator never read), and the
scaffold's integration test, which asserts that a fresh cell fails high.

\section{The freeze policy, restated}\label{the-freeze-policy-restated}

After 0.6.0 the instrument changes only for a failure that changes a
verdict, or for a new execution. A maintainer who finds a false green
adds a row to the table of false greens, a regression test and, where it
is a defect of a document, a perturbation --- and declares the change in
the changelog with its migration note.

\part{Appendices}

\chapter*{Appendix A --- Glossary}\label{appendix-a-glossary}
\addcontentsline{toc}{chapter}{Appendix A --- Glossary}

Condensed from the ``Method terms'' table of
\texttt{02-\allowbreak{}reflexive-\allowbreak{}application/\allowbreak{}01-\allowbreak{}canon.\allowbreak{}md},
which is the authoritative definition of each term.

{\def\LTcaptype{none} % do not increment counter
\begin{longtable}[]{@{}
  >{\raggedright\arraybackslash}p{(\linewidth - 2\tabcolsep) * \real{0.5000}}
  >{\raggedright\arraybackslash}p{(\linewidth - 2\tabcolsep) * \real{0.5000}}@{}}
\toprule\noalign{}
\begin{minipage}[b]{\linewidth}\raggedright
Term
\end{minipage} & \begin{minipage}[b]{\linewidth}\raggedright
Definition
\end{minipage} \\
\midrule\noalign{}
\endhead
\bottomrule\noalign{}
\endlastfoot
accepted variant & another spelling of a number or term that the canon
admits in a declared place and for a declared reason \\
artifact & an intermediate or final output, typed, with a name, a single
owner, paths, a schema, dependencies, gates and content identified by
hash; an owner decision is also an artifact \\
artifact graph & the set of artifacts with the dependencies between
them; it has to be acyclic \\
assembly engine & the code that composes the artifacts into the target,
rejects ambiguity, rebuilds by hash, routes notes and names the outputs
by version label \\
brief & the owner's document: what the target is, for whom, what is
non-negotiable, what is an owner decision, the round budget \\
canon & the single source of numbers, names, terms to prefer and to
avoid, accepted variants and rules, with an owner and a version per
section \\
case method document & the document that records, at the end, the ideas,
graph, step by step, pseudocode per role with its degree, the owner's
pending items and the known false positives of an instance \\
code-conversion degree & how much of a role's work the code does: (a)
converted, (b) partial under gate, (c) contract and gate only \\
contract & the schema an artifact has to validate before it is
accepted \\
corpus & the facts, materials, standards and references available before
the production; read, never edited \\
dispatch & one call of an executor with its inputs, recorded with what
the endpoint returned (tokens, the reason the generation ended) \\
domain gate & a gate of the target's domain, named
\texttt{d-\allowbreak{}<name>}, measured by the cell's own code and read
by the validator from a result file \\
escalated & the state of an artifact that exceeded its maximum rounds:
nothing that depends on it runs until the owner decides \\
executor & whoever or whatever exercises a role: a person, code, or a
language model of a declared kind \\
experiment & a measurement about the method itself, registered with its
question, protocol and predictions written before measuring;
\texttt{executed}, \texttt{pending} or \texttt{designed} \\
false negative & a real defect that no instrument reported \\
false positive & a note the instrument emits when there is no defect; an
instrument defect until proven otherwise \\
family & the provider and line of models an executor belongs to, named
by a stable letter \\
first version & the v1: the whole target, end to end, made quickly to be
reviewed, not to be good \\
gate & a measurable function over an artifact that returns pass or fail
together with the measured magnitude; with a limit and a written
justification \\
gate sensitivity test & injecting a known defect and checking that the
gate rejects it, before trusting the gate \\
hand-off sheet & the artifact one role writes for another role to
consume \\
independent review & a reading of the product by someone who did not
produce it and has no access to the producer's reasoning, with the fixed
rubric \\
instance registry & the record of the known instances, each with an id,
a status, the path of its method document and the sha256 that pins it \\
instantiation table & the table of fourteen fixed rows that summarizes
an instance \\
instrument & anything that measures: gate, validator, reviewer,
simulation, benchmark \\
limit & the value that separates pass from fail, with the justification
and the origin label next to it \\
magnitude & the measured value the gate returns, never just a sign \\
maintenance round & a round that corrects the material of the method
itself after its stopping criterion, outside the reflexive production,
with no verdict of the canon \\
maximum rounds & how many returns an artifact accepts before it is
escalated to the owner \\
note & the only change request between roles: target artifact, severity,
text, origin, evidence, round \\
open note & a note recorded and not resolved when the production
stops \\
orchestrator & whoever executes the recipe: passes paths to the roles,
runs the engine and imports notes; a step, not a role \\
owner & the person who commissions the target, approves the brief, makes
the decisions no role can make and declares the stop \\
position key & the value that fixes where a part of the target lands
during assembly; an ambiguous one is a contract error \\
prior instance & an instance of the method that ran before the current
one and is registered in the registry \\
production cell & the organization assembled for one target: roles,
artifacts, gates, canon, validator, review and assembly engine \\
regression & a defect that comes back or appears in a new version; if it
escaped every gate, it gains a gate before it is fixed \\
retired gate & a gate removed deliberately, with its perturbation,
because its object is gone or its rule moved into a contract \\
role & a narrow responsibility exercised by an agent, a person or code,
which communicates only through artifacts and notes \\
role card & the sheet of a role with seven fixed fields \\
root-cause fix & handling a false positive by fixing the instrument's
rule, with a test that reproduces the case, and redoing the previous
round \\
round & one pass of the loop: produce, pass the gates, assemble, review,
notes \\
rubric & the fixed grid of dimensions, anchors and verdict rule that
every independent review uses in every round \\
rubric anchor & the written description of what a score of 1, 3 or 5
means in one rubric dimension \\
runner & the role that runs the validator and routes its notes; it
decides nothing beyond the rules the validator executes \\
severity & \texttt{low}, \texttt{medium}, \texttt{high},
\texttt{blocking}; only the last two send the artifact back \\
signal gate & a gate that only tests existence or ``greater than zero'';
it has to gain a magnitude limit \\
source of truth & for each content type, the only source accepted as
origin; the writer is never a source of fact \\
stale & the state of an artifact whose input changed since it passed its
gates \\
stopping criterion & the condition, declared by the owner before the
last round, under which the review becomes a list of open notes \\
target & the final product that the production cell hands over \\
teamwork & the files that set up a production cell for a target; the
name of the skill that creates them \\
unit & what a gate reads as one statement: a paragraph, a list item, a
table row or a heading \\
validator & a checking instrument, preferably in code, that reads all
artifacts against the canon and emits notes on its own \\
verdict & \texttt{rejected}, \texttt{accepted with conditions},
\texttt{accepted} \\
wave & a batch of dispatches run in parallel, whose outputs are read
together before the next step \\
\end{longtable}
}

\chapter*{Appendix B --- The contracts}\label{appendix-b-the-contracts}
\addcontentsline{toc}{chapter}{Appendix B --- The contracts}

Six JSON Schemas under
\texttt{02-\allowbreak{}reflexive-\allowbreak{}application/\allowbreak{}05-\allowbreak{}contracts/\allowbreak{}},
each declaring \texttt{contract\_\allowbreak{}version} \texttt{1} and
rejecting unknown fields. The validator applies a closed subset of JSON
Schema keywords and reports a contract that uses others.

{\def\LTcaptype{none} % do not increment counter
\begin{longtable}[]{@{}
  >{\raggedright\arraybackslash}p{(\linewidth - 6\tabcolsep) * \real{0.2500}}
  >{\raggedright\arraybackslash}p{(\linewidth - 6\tabcolsep) * \real{0.2500}}
  >{\raggedright\arraybackslash}p{(\linewidth - 6\tabcolsep) * \real{0.2500}}
  >{\raggedright\arraybackslash}p{(\linewidth - 6\tabcolsep) * \real{0.2500}}@{}}
\toprule\noalign{}
\begin{minipage}[b]{\linewidth}\raggedright
Contract
\end{minipage} & \begin{minipage}[b]{\linewidth}\raggedright
File
\end{minipage} & \begin{minipage}[b]{\linewidth}\raggedright
Required fields
\end{minipage} & \begin{minipage}[b]{\linewidth}\raggedright
Optional fields
\end{minipage} \\
\midrule\noalign{}
\endhead
\bottomrule\noalign{}
\endlastfoot
Note & \texttt{notes.\allowbreak{}schema.\allowbreak{}json} &
\texttt{target\_\allowbreak{}artifact}, \texttt{severity} (the four
values), \texttt{text}, \texttt{origin}, \texttt{evidence},
\texttt{round} (string, or the legacy integer) & \texttt{status}
(\texttt{open}, \texttt{fixed}, \texttt{disputed},
\texttt{out of scope}), \texttt{closed\_\allowbreak{}in} \\
State of a note &
\texttt{notes-\allowbreak{}status.\allowbreak{}schema.\allowbreak{}json}
& \texttt{ref} (\texttt{notes:\allowbreak{}<n>} or
\texttt{panel:\allowbreak{}<version>:\allowbreak{}<id>}),
\texttt{status}, \texttt{summary} & \texttt{closed\_\allowbreak{}in},
\texttt{prev}, \texttt{escaped}, \texttt{gate},
\texttt{no\_\allowbreak{}gate}, \texttt{decision},
\texttt{opened\_\allowbreak{}in}, \texttt{due}, \texttt{severity},
\texttt{affects} \\
Instance registry &
\texttt{registry.\allowbreak{}schema.\allowbreak{}json} &
\texttt{instances}: a list of entries with \texttt{id}, \texttt{status}
(\texttt{executed}, \texttt{designed}) & per entry: \texttt{document},
\texttt{sha256}, \texttt{locator}, \texttt{number\_\allowbreak{}format}
(\texttt{point-\allowbreak{}decimal},
\texttt{comma-\allowbreak{}decimal}),
\texttt{heading\_\allowbreak{}style} (\texttt{section},
\texttt{subsection}, \texttt{step}),
\texttt{method\_\allowbreak{}version}; an executed entry carries the
first four non-null, a designed one keeps them null \\
Round record &
\texttt{round-\allowbreak{}record.\allowbreak{}schema.\allowbreak{}json}
& \texttt{method\_\allowbreak{}version}, \texttt{round}, \texttt{models}
(a list of \texttt{role}, \texttt{model}, with optional \texttt{tokens}
of \texttt{input}, \texttt{output}, \texttt{reasoning}, \texttt{cached};
\texttt{calls}; \texttt{finish}), \texttt{summary} & \texttt{verdict},
\texttt{rubric\_\allowbreak{}hash},
\texttt{rubric\_\allowbreak{}author},
\texttt{family\_\allowbreak{}map\_\allowbreak{}commitment},
\texttt{prev}, \texttt{kind} (\texttt{production},
\texttt{maintenance}), \texttt{reviewer\_\allowbreak{}relation}
(\texttt{another family}, \texttt{same family, separate session},
\texttt{same model, separate route}, \texttt{none}),
\texttt{rubric\_\allowbreak{}applied},
\texttt{rounds\_\allowbreak{}to\_\allowbreak{}stop},
\texttt{wall\_\allowbreak{}seconds}, \texttt{notes\_\allowbreak{}count},
\texttt{crew} \\
Domain result &
\texttt{domain-\allowbreak{}result.\allowbreak{}schema.\allowbreak{}json}
& \texttt{contract\_\allowbreak{}version}, \texttt{gate},
\texttt{magnitude}, \texttt{limit}, \texttt{result} (\texttt{pass},
\texttt{fail}, \texttt{tool-\allowbreak{}failure}), \texttt{inputs} (a
list of \texttt{path}, \texttt{sha256}) & --- \\
Owner decisions &
\texttt{owner-\allowbreak{}decisions.\allowbreak{}schema.\allowbreak{}json}
& \texttt{name}, \texttt{status} (\texttt{open}, \texttt{decided}),
\texttt{blocks}, \texttt{version}, \texttt{summary} & \texttt{covers},
\texttt{prev} \\
\end{longtable}
}

The scaffold writes a seventh,
\texttt{02-\allowbreak{}canon.\allowbreak{}schema.\allowbreak{}json},
the contract of a cell's JSON canon, and copies the first five except
the registry's (a cell's \texttt{priors.\allowbreak{}json} follows the
same contract, enforced by the validator).

\chapter*{Appendix C --- Command-line
reference}\label{appendix-c-command-line-reference}
\addcontentsline{toc}{chapter}{Appendix C --- Command-line reference}

\section{\texorpdfstring{\texttt{validator.\allowbreak{}py}}{}}\label{validator.py}

{\def\LTcaptype{none} % do not increment counter
\begin{longtable}[]{@{}
  >{\raggedright\arraybackslash}p{(\linewidth - 2\tabcolsep) * \real{0.5000}}
  >{\raggedright\arraybackslash}p{(\linewidth - 2\tabcolsep) * \real{0.5000}}@{}}
\toprule\noalign{}
\begin{minipage}[b]{\linewidth}\raggedright
Option
\end{minipage} & \begin{minipage}[b]{\linewidth}\raggedright
Meaning
\end{minipage} \\
\midrule\noalign{}
\endhead
\bottomrule\noalign{}
\endlastfoot
\texttt{-\allowbreak{}-\allowbreak{}folder FOLDER} & the production
folder (default \texttt{.\allowbreak{}}) \\
\texttt{-\allowbreak{}-\allowbreak{}layout LAYOUT} & a
\texttt{layout.\allowbreak{}json} that remaps artifacts to paths and
gates to parsers; default is the Markdown layout of the corpus \\
\texttt{-\allowbreak{}-\allowbreak{}registry REGISTRY} & the instance
registry; default
\texttt{.\allowbreak{}.\allowbreak{}/\allowbreak{}instances/\allowbreak{}registry.\allowbreak{}json}
relative to the folder, when present; without a registry
\texttt{p-\allowbreak{}prior} is pending \\
\texttt{-\allowbreak{}-\allowbreak{}case CASE} & a prior method
document, repeatable; each is identified by its sha256 against the
registry \\
\texttt{-\allowbreak{}-\allowbreak{}forbidden FORBIDDEN} & the list of
forbidden identifiers, one per line (names or
\texttt{sha256:\allowbreak{}<hex>}); without it that half of
\texttt{p-\allowbreak{}anonymity} is pending \\
\texttt{-\allowbreak{}-\allowbreak{}rubric-\allowbreak{}hash FILE} & the
file with the frozen sha256 of the rubric \\
\texttt{-\allowbreak{}-\allowbreak{}freeze-\allowbreak{}rubric FILE} &
write the current sha256 of the rubric to the file and exit; refuses an
existing pin (the scaffold's \texttt{not-\allowbreak{}frozen}
placeholder excepted) \\
\texttt{-\allowbreak{}-\allowbreak{}re-\allowbreak{}freeze-\allowbreak{}rubric FILE}
& replace an existing frozen rubric deliberately; mutually exclusive
with the previous \\
\texttt{-\allowbreak{}-\allowbreak{}notes NOTES} & the note ledger to
validate against the contract, with its state ledger \\
\texttt{-\allowbreak{}-\allowbreak{}output FILE} & write the report to
the file; one summary line goes to the error stream \\
\texttt{-\allowbreak{}-\allowbreak{}determinism} & run twice more with
different hash seeds (and the C locale) and compare the reports \\
\texttt{-\allowbreak{}-\allowbreak{}allow-\allowbreak{}pending LIST} &
the gates whose pending rows the caller accepts, comma-separated, from
the closed list \texttt{p-\allowbreak{}anonymity},
\texttt{p-\allowbreak{}prior},
\texttt{p-\allowbreak{}prior-\allowbreak{}pin},
\texttt{p-\allowbreak{}sensitivity-\allowbreak{}synthetic} \\
\texttt{-\allowbreak{}-\allowbreak{}self-\allowbreak{}test} & run the
embedded checks and exit \\
\end{longtable}
}

Exit codes: 0, 1, 2, 4 (chapter 12).

\section{\texorpdfstring{\texttt{sensitivity.\allowbreak{}py}}{}}\label{sensitivity.py}

{\def\LTcaptype{none} % do not increment counter
\begin{longtable}[]{@{}
  >{\raggedright\arraybackslash}p{(\linewidth - 2\tabcolsep) * \real{0.5000}}
  >{\raggedright\arraybackslash}p{(\linewidth - 2\tabcolsep) * \real{0.5000}}@{}}
\toprule\noalign{}
\begin{minipage}[b]{\linewidth}\raggedright
Option
\end{minipage} & \begin{minipage}[b]{\linewidth}\raggedright
Meaning
\end{minipage} \\
\midrule\noalign{}
\endhead
\bottomrule\noalign{}
\endlastfoot
\texttt{-\allowbreak{}-\allowbreak{}folder FOLDER} & the production
folder to copy and perturb \\
\texttt{-\allowbreak{}-\allowbreak{}case CASE} & as in the validator,
repeatable; without any, the rows that need a prior are not measured \\
\texttt{-\allowbreak{}-\allowbreak{}forbidden FORBIDDEN} & the list of
forbidden identifiers \\
\texttt{-\allowbreak{}-\allowbreak{}registry REGISTRY} & the instance
registry \\
\texttt{-\allowbreak{}-\allowbreak{}rubric-\allowbreak{}hash FILE} &
forwarded to the validator; without it the
\texttt{p-\allowbreak{}frozen-\allowbreak{}rubric} row is not
measured \\
\texttt{-\allowbreak{}-\allowbreak{}notes NOTES} & forwarded; without it
the \texttt{p-\allowbreak{}notes} row is not measured \\
\texttt{-\allowbreak{}-\allowbreak{}work WORK} & the scratch folder for
the copies (default: a temporary folder) \\
\texttt{-\allowbreak{}-\allowbreak{}output FILE} & where the report is
written (default: the layout key
\texttt{artifacts.\allowbreak{}sensitivity}) \\
\texttt{-\allowbreak{}-\allowbreak{}synthetic-\allowbreak{}prior ID} &
an executed prior whose method document you do not have: a synthetic
document is generated from the canon's recorded facts and pinned in the
copies only \\
\texttt{-\allowbreak{}-\allowbreak{}layout LAYOUT} & the cell's
\texttt{layout.\allowbreak{}json}, which also names the cell's
perturbation file \\
\end{longtable}
}

An empty perturbation set exits 4; a gate the validator does not
implement in the set is a tool failure (exit 2).

\section{\texorpdfstring{\texttt{scaffold.\allowbreak{}py}}{}}\label{scaffold.py}

{\def\LTcaptype{none} % do not increment counter
\begin{longtable}[]{@{}
  >{\raggedright\arraybackslash}p{(\linewidth - 2\tabcolsep) * \real{0.5000}}
  >{\raggedright\arraybackslash}p{(\linewidth - 2\tabcolsep) * \real{0.5000}}@{}}
\toprule\noalign{}
\begin{minipage}[b]{\linewidth}\raggedright
Option
\end{minipage} & \begin{minipage}[b]{\linewidth}\raggedright
Meaning
\end{minipage} \\
\midrule\noalign{}
\endhead
\bottomrule\noalign{}
\endlastfoot
\texttt{-\allowbreak{}-\allowbreak{}target TARGET} & the name of the
target (required) \\
\texttt{-\allowbreak{}-\allowbreak{}roles ROLE [ROLE .\allowbreak{}.\allowbreak{}.\allowbreak{}]}
& role ids, lowercase letters, digits and hyphens (required) \\
\texttt{-\allowbreak{}-\allowbreak{}out OUT} & the destination folder
(required); the scaffold refuses to overwrite \\
\end{longtable}
}

Run from inside the material, so that it finds the validator, the
sensitivity script, the perturbation module and the contracts to copy.

\section{The historical validator of instance
004}\label{the-historical-validator-of-instance-004}

\texttt{instances/\allowbreak{}004-\allowbreak{}adoption-\allowbreak{}guide/\allowbreak{}08-\allowbreak{}validator/\allowbreak{}validator.\allowbreak{}py}
is the instrument of 0.2.2 and refuses to run without
\texttt{-\allowbreak{}-\allowbreak{}historical}, since its exit 0 does
not mean what the current validator's does.

\chapter*{Appendix D --- The false greens of the
gates}\label{appendix-d-the-false-greens-of-the-gates}
\addcontentsline{toc}{chapter}{Appendix D --- The false greens of the
gates}

From the table of
\texttt{06-\allowbreak{}gate-\allowbreak{}semantics.\allowbreak{}md},
``The false greens of the gates'': for each gate, the cause of a false
green found so far, what detects it today, and the perturbation that
covers it. The causes are the ones found, not an exhaustive list.

{\def\LTcaptype{none} % do not increment counter
\begin{longtable}[]{@{}
  >{\raggedright\arraybackslash}p{(\linewidth - 6\tabcolsep) * \real{0.2500}}
  >{\raggedright\arraybackslash}p{(\linewidth - 6\tabcolsep) * \real{0.2500}}
  >{\raggedright\arraybackslash}p{(\linewidth - 6\tabcolsep) * \real{0.2500}}
  >{\raggedright\arraybackslash}p{(\linewidth - 6\tabcolsep) * \real{0.2500}}@{}}
\toprule\noalign{}
\begin{minipage}[b]{\linewidth}\raggedright
Gate
\end{minipage} & \begin{minipage}[b]{\linewidth}\raggedright
Cause of a false green
\end{minipage} & \begin{minipage}[b]{\linewidth}\raggedright
Current detection
\end{minipage} & \begin{minipage}[b]{\linewidth}\raggedright
Perturbation that covers it
\end{minipage} \\
\midrule\noalign{}
\endhead
\bottomrule\noalign{}
\endlastfoot
\texttt{p-\allowbreak{}terms} & a synonym written in a way the lexical
match does not see & the canon's list and whole-word matching; the
review reads the prose & avoid-word in the changelog \\
\texttt{p-\allowbreak{}canon-\allowbreak{}link} & a link that names the
canon from the wrong folder & the target resolved relative to the
document's folder & link to the canon does not resolve \\
\texttt{p-\allowbreak{}anonymity} & a path, a date or an identifier in a
form the patterns do not know, or hidden by the valve & the patterns,
the months of six languages and the canon's own, the forbidden list by
name or sha256; an accepted path dropped only by equality & date written
with the month's name \\
\texttt{p-\allowbreak{}form} & an HTML tag or comment hidden in a block
the reader skips & one reader of fences, comments read & HTML comment \\
\texttt{p-\allowbreak{}size} & a document pushed past its range, or
words counted wrong & the declared ranges; table pipes are no words &
canon pushed above its size range \\
\texttt{p-\allowbreak{}table} & an instantiation table missing a row &
the canon's fixed list of rows & instantiation table row deleted \\
\texttt{p-\allowbreak{}sections} & an instance missing a mandatory
section & the canon's fixed list of sections & mandatory section heading
removed \\
\texttt{p-\allowbreak{}card} & a card with a field empty, deleted or a
placeholder & the seven fields, the placeholder detector & placeholder
in a card field \\
\texttt{p-\allowbreak{}limit} & a justification of six words of any
kind, or a limit raised in the table & the origin label and the word
count; the Limit column compared & limit raised in the gate table \\
\texttt{p-\allowbreak{}references} & a reference to a role, an artifact
or a gate that does not exist & the artifact table, the graph without
cycles, the commands resolved & limit cell diverges from the applied
limit \\
\texttt{p-\allowbreak{}prior} & an invented number about a prior: in a
code span, under a malformed mark, in a section the reader does not find
& each number checked in the sections its recorded fact cites; section
commitments; the synthetic prior in the registry's heading style & prior
number in a code span \\
\texttt{p-\allowbreak{}prior-\allowbreak{}mark} & a statement about a
prior with no mark, or the mark in a code span & the canon's triggers,
read outside code spans & prior mark inside a code span \\
\texttt{p-\allowbreak{}prior-\allowbreak{}pin} & a prior document
swapped or edited, or a pin no run compares & the registry's sha256
compared with each document passed, prior by prior & prior pin altered
in the registry \\
\texttt{p-\allowbreak{}contract-\allowbreak{}shape} & a ledger line
rewritten, a contract of another version, a registry entry that breaks
its status & the schemas line by line, the hash chain (an edit that
recomputes it is not detected) & old line of the round record edited \\
\texttt{p-\allowbreak{}recipe} & a step of the recipe without one of its
fields & the canon's step fields & step field deleted \\
\texttt{p-\allowbreak{}components} & a component the principle names and
an instance omits & the canon's list & method component omitted \\
\texttt{p-\allowbreak{}efficacy-\allowbreak{}mark} & a claim that the
method works without a label, or with a label citing nothing executed &
the canon's triggers (lexical), the provenance table of an experiment, a
lessons title that resolves & claim backed by a designed experiment \\
\texttt{p-\allowbreak{}notes} & a note without a state, a regression
closed with a gate that does not exist, a protection undone by a later
line & the state ledger, golden rule 4 against the gate table, the
highest severity and the union of \texttt{affects} & note without a
status line \\
\texttt{p-\allowbreak{}frozen-\allowbreak{}rubric} & a rubric or a
frozen record rewritten after the fact & the sha256 of the rubric and of
each frozen record; the attestation in the owner's run & frozen review
report rewritten \\
\texttt{p-\allowbreak{}sensitivity} & a report the instrument did not
produce, a stale one, or one made without the owner's document read as
made with it & the provenance header, the case lines against the pins,
the attestation, a synthetic prior never a pass & none: the baseline
reads the report; its defects are regression tests \\
\texttt{p-\allowbreak{}declared} & a declared path that matches nothing,
a repointed file, a document no artifact declares & each literal path
and each glob, the resolved map pinned in the sensitivity's provenance &
file of an artifact repointed \\
\end{longtable}
}

\chapter*{Appendix E --- File index of the
material}\label{appendix-e-file-index-of-the-material}
\addcontentsline{toc}{chapter}{Appendix E --- File index of the
material}

The 79 files the material tracks at 0.6.0, grouped.

{\def\LTcaptype{none} % do not increment counter
\begin{longtable}[]{@{}
  >{\raggedright\arraybackslash}p{(\linewidth - 2\tabcolsep) * \real{0.5000}}
  >{\raggedright\arraybackslash}p{(\linewidth - 2\tabcolsep) * \real{0.5000}}@{}}
\toprule\noalign{}
\begin{minipage}[b]{\linewidth}\raggedright
Path
\end{minipage} & \begin{minipage}[b]{\linewidth}\raggedright
Contents
\end{minipage} \\
\midrule\noalign{}
\endhead
\bottomrule\noalign{}
\endlastfoot
\texttt{README.\allowbreak{}md}, \texttt{VERSION},
\texttt{CHANGELOG.\allowbreak{}md}, \texttt{LICENSE},
\texttt{.\allowbreak{}gitattributes}, \texttt{.\allowbreak{}gitignore} &
the entry page, the version, the history, the license, the LF rule, the
ignore list \\
\texttt{01-\allowbreak{}principle.\allowbreak{}md} & the method \\
\texttt{02-\allowbreak{}reflexive-\allowbreak{}application/\allowbreak{}README.\allowbreak{}md}
& the index and the runbook of the reflexive application \\
\texttt{02-\allowbreak{}reflexive-\allowbreak{}application/\allowbreak{}00-\allowbreak{}brief.\allowbreak{}md}
\ldots{}
\texttt{06-\allowbreak{}gate-\allowbreak{}semantics.\allowbreak{}md} &
brief, canon, artifacts, four role cards, gates and rubric, gate
semantics \\
\texttt{02-\allowbreak{}reflexive-\allowbreak{}application/\allowbreak{}05-\allowbreak{}contracts/\allowbreak{}*.\allowbreak{}schema.\allowbreak{}json}
& the six contracts \\
\texttt{02-\allowbreak{}reflexive-\allowbreak{}application/\allowbreak{}10-\allowbreak{}generic-\allowbreak{}recipe.\allowbreak{}md}
& the fifteen steps \\
\texttt{02-\allowbreak{}reflexive-\allowbreak{}application/\allowbreak{}20-\allowbreak{}instance-\allowbreak{}car.\allowbreak{}md},
\texttt{21-\allowbreak{}instance-\allowbreak{}text-\allowbreak{}2027.\allowbreak{}md},
\texttt{22-\allowbreak{}instance-\allowbreak{}hpc.\allowbreak{}md} & the
three designed instances \\
\texttt{02-\allowbreak{}reflexive-\allowbreak{}application/\allowbreak{}60-\allowbreak{}experiments.\allowbreak{}md}
& E1 to E5 \\
\texttt{02-\allowbreak{}reflexive-\allowbreak{}application/\allowbreak{}80-\allowbreak{}validator/\allowbreak{}validator.\allowbreak{}py},
\texttt{sensitivity.\allowbreak{}py},
\texttt{perturbation\_\allowbreak{}set.\allowbreak{}py},
\texttt{test\_\allowbreak{}validator.\allowbreak{}py} & the instrument
and its suite \\
\texttt{02-\allowbreak{}reflexive-\allowbreak{}application/\allowbreak{}80-\allowbreak{}validator/\allowbreak{}report.\allowbreak{}md},
\texttt{sensitivity.\allowbreak{}md} & the published reports \\
\texttt{02-\allowbreak{}reflexive-\allowbreak{}application/\allowbreak{}90-\allowbreak{}review/\allowbreak{}rubric.\allowbreak{}sha256},
\texttt{auditor-\allowbreak{}report.\allowbreak{}md},
\texttt{auditor-\allowbreak{}notes.\allowbreak{}json},
\texttt{notes.\allowbreak{}json}, \texttt{replies.\allowbreak{}md} & the
frozen records of round 1 \\
\texttt{02-\allowbreak{}reflexive-\allowbreak{}application/\allowbreak{}99-\allowbreak{}lessons.\allowbreak{}md}
& the lessons, one entry per round \\
\texttt{02-\allowbreak{}reflexive-\allowbreak{}application/\allowbreak{}state/\allowbreak{}round-\allowbreak{}record.\allowbreak{}jsonl},
\texttt{notes-\allowbreak{}status.\allowbreak{}jsonl},
\texttt{owner-\allowbreak{}decisions.\allowbreak{}jsonl} & the three
chained ledgers \\
\texttt{skills/\allowbreak{}cooperatio-\allowbreak{}teamwork/\allowbreak{}SKILL.\allowbreak{}md},
\texttt{references/\allowbreak{}templates.\allowbreak{}md},
\texttt{references/\allowbreak{}review-\allowbreak{}checklist.\allowbreak{}md},
\texttt{scripts/\allowbreak{}scaffold.\allowbreak{}py},
\texttt{scripts/\allowbreak{}test\_\allowbreak{}scaffold.\allowbreak{}py}
& the skill \\
\texttt{instances/\allowbreak{}registry.\allowbreak{}json} & the
instance registry \\
\texttt{instances/\allowbreak{}004-\allowbreak{}adoption-\allowbreak{}guide/\allowbreak{}}
& the cell of the second executed instance: brief, sources of truth,
JSON canon, artifacts, four cards, two contracts, gates, rubric and pin,
the validator of 0.2.2 with its rules and reports, the engine's README,
the two wave reports and the round-1 notes, the state ledgers, lessons,
open notes, priors, layout, the guide
(\texttt{guia-\allowbreak{}de-\allowbreak{}adocao.\allowbreak{}md}), the
method document
(\texttt{metodo-\allowbreak{}de-\allowbreak{}producao.\allowbreak{}md})
and \texttt{ERRATA.\allowbreak{}md} \\
\texttt{docs/\allowbreak{}manual/\allowbreak{}} & this manual:
\texttt{src/\allowbreak{}}, \texttt{diagrams/\allowbreak{}},
\texttt{filters/\allowbreak{}}, \texttt{Makefile}, the built
\texttt{manual.\allowbreak{}md}, \texttt{manual.\allowbreak{}pdf} and
\texttt{build/\allowbreak{}manual.\allowbreak{}tex} \\
\end{longtable}
}

\chapter*{Appendix F --- Open owner
decisions}\label{appendix-f-open-owner-decisions}
\addcontentsline{toc}{chapter}{Appendix F --- Open owner decisions}

The lines of
\texttt{02-\allowbreak{}reflexive-\allowbreak{}application/\allowbreak{}state/\allowbreak{}owner-\allowbreak{}decisions.\allowbreak{}jsonl}
at 0.6.0, all \texttt{open}; each names what it blocks. Only the owner
can close one, with a new chained line.

{\def\LTcaptype{none} % do not increment counter
\begin{longtable}[]{@{}
  >{\raggedright\arraybackslash}p{(\linewidth - 2\tabcolsep) * \real{0.5000}}
  >{\raggedright\arraybackslash}p{(\linewidth - 2\tabcolsep) * \real{0.5000}}@{}}
\toprule\noalign{}
\begin{minipage}[b]{\linewidth}\raggedright
Decision
\end{minipage} & \begin{minipage}[b]{\linewidth}\raggedright
What it blocks
\end{minipage} \\
\midrule\noalign{}
\endhead
\bottomrule\noalign{}
\endlastfoot
sign the tag of the version (create a signing key and sign
\texttt{v0.\allowbreak{}6.\allowbreak{}0}) & the anchor of integrity
outside the tree \\
publish an extract of the first prior, or not & a clone's check of the
numbers about \texttt{000-\allowbreak{}av-\allowbreak{}production} \\
accept or adjust the thresholds to promote or withdraw a claim &
experiment E4 \\
approve the budget and run E4 & the measure of value of the material \\
find the third party for E5 & the measure of portability \\
promote independence across families to a rule & the rule of
independence \\
open the next production of the material as a production, with its own
brief, a rubric of the reviewing side and a stop & a new verdict of the
canon on the material \\
replace the cell of 004 as the example (a new cell under 0.6.0 would be
the third execution) & an executable example of a cell \\
\end{longtable}
}

\chapter*{Appendix G --- Building this
manual}\label{appendix-g-building-this-manual}
\addcontentsline{toc}{chapter}{Appendix G --- Building this manual}

The manual lives in \texttt{docs/\allowbreak{}manual/\allowbreak{}}:

{\def\LTcaptype{none} % do not increment counter
\begin{longtable}[]{@{}
  >{\raggedright\arraybackslash}p{(\linewidth - 2\tabcolsep) * \real{0.5000}}
  >{\raggedright\arraybackslash}p{(\linewidth - 2\tabcolsep) * \real{0.5000}}@{}}
\toprule\noalign{}
\begin{minipage}[b]{\linewidth}\raggedright
Path
\end{minipage} & \begin{minipage}[b]{\linewidth}\raggedright
What it is
\end{minipage} \\
\midrule\noalign{}
\endhead
\bottomrule\noalign{}
\endlastfoot
\texttt{src/\allowbreak{}*.\allowbreak{}md} & the source, one file per
part, concatenated in order \\
\texttt{diagrams/\allowbreak{}*.\allowbreak{}dot} & the Graphviz twin of
each Mermaid diagram; rendered to \texttt{.\allowbreak{}svg} and
\texttt{.\allowbreak{}pdf} \\
\texttt{filters/\allowbreak{}mermaid.\allowbreak{}lua} & the pandoc
filter that swaps a Mermaid block for its twin's PDF in the PDF build \\
\texttt{filters/\allowbreak{}header.\allowbreak{}tex} & the LaTeX
preamble additions (fonts, shaded code blocks, table spacing) \\
\texttt{Makefile} & the build \\
\texttt{manual.\allowbreak{}md} & the concatenated Markdown, for viewers
that render Mermaid \\
\texttt{build/\allowbreak{}manual.\allowbreak{}tex},
\texttt{manual.\allowbreak{}pdf} & the LaTeX and the PDF editions \\
\end{longtable}
}

Each Mermaid block opens with two Mermaid comments,
\texttt{\%\% name:\allowbreak{}} (the twin under
\texttt{diagrams/\allowbreak{}}) and
\texttt{\%\% caption:\allowbreak{}}, which the filter reads; a Markdown
viewer ignores them. The build needs pandoc 3 or later, Graphviz and
LuaLaTeX with the TeX Gyre fonts (TeX Live 2025 was used); no Mermaid
renderer is needed:

\begin{Shaded}
\begin{Highlighting}[]
\BuiltInTok{cd}\NormalTok{ docs/manual}
\FunctionTok{make}            \CommentTok{\# manual.md, the diagrams, build/manual.tex, manual.pdf}
\FunctionTok{make}\NormalTok{ check      }\CommentTok{\# margin overflow check of the PDF, where the local tool exists}
\end{Highlighting}
\end{Shaded}

The manual follows the form rules of the material where it can: no
absolute path, machine name, user name, model identifier or authoring
date; executors by family letter; claims of efficacy with their label.

\end{document}
